% ECONOMIC GEOLOGY: Economic Minerals — Geology Upper Sixth — Cours signé OnBuch+
% Official MINESEC syllabus. Compiler avec : tectonic GEO05-economic-geology-economic-minerals.tex
\newcommand{\DOCMATIERE}{Geology}
\newcommand{\DOCNIVEAU}{Upper Sixth}
\newcommand{\DOCMODULE}{Upper Sixth — Geology}
\newcommand{\DOCLECON}{5}
\newcommand{\DOCTITRE}{ECONOMIC GEOLOGY: Economic Minerals}
% OnBuch+ — préambule « cours signés » ENRICHI (compilé avec tectonic / XeLaTeX)
% Système visuel « Pop » d'OnBuch+ : fond crème (jamais blanc), bordures encre
% épaisses, ombres « dures » décalées, titres Archivo Black en capitales.
% Version MATHÉMATIQUES : graphes (pgfplots/tikz), boîtes pédagogiques
% (définition, propriété, méthode, attention, le savais-tu, exercice résolu,
% exercice, TP, bilan), page de garde et signature OnBuch+ ancrée en bas.
%
% Macros attendues dans main.tex AVANT \input{preamble} :
%   \DOCMATIERE  \DOCNIVEAU  \DOCMODULE  \DOCLECON (numéro)  \DOCTITRE
\documentclass[11pt]{article}

\usepackage{fontspec}
\usepackage[english]{babel}
\usepackage[a4paper,margin=2cm,top=3.4cm,bottom=2.8cm,headheight=1.4cm,footskip=1.3cm]{geometry}
\usepackage{amsmath,amssymb,amsfonts}
\usepackage{mathtools} % \xmapsto, \coloneqq... souvent utilisés par les modèles
\usepackage{cancel} % simplifications barrées (bilans de forces)
% \abs{z} (module, valeur absolue) et \norm{u} (norme), utilisés par les modèles
\providecommand{\abs}{}\let\abs\relax\DeclarePairedDelimiter{\abs}{\lvert}{\rvert}
\providecommand{\norm}{}\let\norm\relax\DeclarePairedDelimiter{\norm}{\lVert}{\rVert}
\usepackage[table]{xcolor} % [table] : fournit aussi \rowcolors (alternance de lignes)
\usepackage{pagecolor}
\usepackage{tikz}
\usetikzlibrary{arrows.meta,positioning,calc,shapes.geometric,shapes.misc,shapes.arrows,decorations.pathmorphing,decorations.pathreplacing,decorations.markings,patterns,shadows,mindmap,trees,fit,backgrounds,matrix,intersections,angles,quotes}
\usepackage{pgfplots}
\usepackage[version=4]{mhchem} % équations nucléaires \ce{^{235}_{92}U}
\usepackage[european,siunitx]{circuitikz} % schémas électriques
\pgfplotsset{compat=1.18}
\usepgfplotslibrary{fillbetween,groupplots} % aires entre courbes (calcul intégral)
\usepackage{enumitem}
\usepackage{fancyhdr}
% [most] charge toutes les bibliothèques tcolorbox (listings, minted, raster,
% documentation...) et gonfle inutilement la mémoire TeX sur un document long
% (~300 boîtes) : on ne charge que ce qui est réellement utilisé.
\usepackage[skins,breakable]{tcolorbox}
\usepackage{graphicx}
\usepackage{colortbl}
\usepackage{booktabs}
\usepackage{tabularx}
\usepackage{diagbox} % case d'en-tête diagonale des tableaux à double entrée
% Colonnes extensibles centrée / gauche / droite (C, L, R), souvent utilisées par les modèles
\newcolumntype{C}{>{\centering\arraybackslash}X}
\newcolumntype{L}{>{\raggedright\arraybackslash}X}
\newcolumntype{R}{>{\raggedleft\arraybackslash}X}
\usepackage{array}
\usepackage{multicol}
\usepackage{siunitx}
% parse-numbers=false : accepte aussi \SI{2,5 \times 10^{-3}}{...}
\sisetup{locale=UK,output-decimal-marker={.},group-minimum-digits=4,parse-numbers=false}
\DeclareSIUnit\ohm{\ensuremath{\symup{\Omega}}} % U+2126 absent de la police
\DeclareSIUnit\division{div} % réglages d'oscilloscope (ms/div, V/div)
\DeclareSIUnit\tr{tr} % tours (vitesse de rotation, ex. \SI{1500}{\tr\per\minute})
\usepackage{qrcode}
% \degree : commande fréquemment utilisée par les modèles pour un angle ou une
% température en dehors de \SI{}{} (non fournie par les paquets chargés).
\providecommand{\degree}{^{\circ}}
% \qed (fin de démonstration), utilisé par les modèles sans amsthm
\providecommand{\qed}{\hfill$\square$}
% \barre : double barre des valeurs interdites dans un tableau de variations (tkz-tab)
\providecommand{\barre}{\ensuremath{\|}}
% environnement proof (amsthm non chargé) utilisé par les modèles
\ifdefined\proof\else\newenvironment{proof}[1][Proof]{\par\noindent\textit{#1.}\ }{\qed\par}\fi
\usepackage{lastpage}
\usepackage{hyperref}
\hypersetup{hidelinks}

% ============================================================
% Palette « Pop » OnBuch+ — mêmes hex que la classe P (plus_ui.dart)
% ============================================================
\definecolor{poporange}{HTML}{F59321}
\definecolor{popdark}{HTML}{E07A0C}
\definecolor{popcream}{HTML}{FFF6E8}
\definecolor{popink}{HTML}{1C1714}
\definecolor{poppurple}{HTML}{7A5AE0}
\definecolor{popgreen}{HTML}{2E9E6A}
\definecolor{poppink}{HTML}{E0457B}
\definecolor{popblue}{HTML}{2F7DE1}
\definecolor{popgold}{HTML}{C98A12}
\definecolor{popmuted}{HTML}{8A7F76}
\definecolor{popline}{HTML}{EADFD2}
\definecolor{popgray}{HTML}{8A7F76}
\colorlet{popgrey}{popgray}
% Alias tolérants (le modèle nomme parfois les couleurs différemment)
\colorlet{popwhite}{white}
\colorlet{popblack}{popink}
\colorlet{popred}{poppink}
\colorlet{popyellow}{popgold}
% Teintes claires pour fonds de tableaux / zones
\colorlet{popblueL}{popblue!10!white}
\colorlet{poporangeL}{poporange!14!white}
\colorlet{popgreenL}{popgreen!12!white}
\colorlet{poppurpleL}{poppurple!12!white}
\colorlet{poppinkL}{poppink!12!white}
\colorlet{popgoldL}{popgold!14!white}

\pagecolor{popcream}

% ============================================================
% Typographie
% ============================================================
\defaultfontfeatures{Ligatures=TeX}
\setmainfont{PlusJakartaSans-Medium.ttf}[Path=../../../fonts/,BoldFont=PlusJakartaSans-Bold.ttf]
% Maths en sans-serif (Fira Math) avec chiffres et lettres droites de Plus
% Jakarta, pour que les formules se fondent dans le texte.
\usepackage[math-style=ISO,bold-style=ISO]{unicode-math}
\setmathfont{FiraMath-Regular.otf}[Path=../../../fonts/]
\setmathfont{PlusJakartaSans-Medium.ttf}[Path=../../../fonts/,range={up/num,up/{Latin,latin},\mathup}]
% (icomma retiré : décimales avec point en anglais)
% Caractères mathématiques Unicode absents des polices de texte (Plus Jakarta,
% Archivo Black) : rendus par leur équivalent mathématique, en texte comme en
% formule — sinon ils s'affichent en carré vide (ex. « Arithmétique dans ℤ »).
\usepackage{newunicodechar}
\newunicodechar{ℝ}{\ensuremath{\mathbb{R}}}
\newunicodechar{ℤ}{\ensuremath{\mathbb{Z}}}
\newunicodechar{ℂ}{\ensuremath{\mathbb{C}}}
\newunicodechar{ℕ}{\ensuremath{\mathbb{N}}}
\newunicodechar{ℚ}{\ensuremath{\mathbb{Q}}}
\newunicodechar{𝔻}{\ensuremath{\mathbb{D}}}
\newunicodechar{α}{\ensuremath{\alpha}}
\newunicodechar{β}{\ensuremath{\beta}}
\newunicodechar{γ}{\ensuremath{\gamma}}
\newunicodechar{δ}{\ensuremath{\delta}}
\newunicodechar{ε}{\ensuremath{\varepsilon}}
\newunicodechar{θ}{\ensuremath{\theta}}
\newunicodechar{λ}{\ensuremath{\lambda}}
\newunicodechar{μ}{\ensuremath{\mu}}
\newunicodechar{σ}{\ensuremath{\sigma}}
\newunicodechar{τ}{\ensuremath{\tau}}
\newunicodechar{φ}{\ensuremath{\varphi}}
\newunicodechar{ψ}{\ensuremath{\psi}}
\newunicodechar{ω}{\ensuremath{\omega}}
\newunicodechar{Σ}{\ensuremath{\Sigma}}
% majuscules produites par \MakeUppercase dans les titres (α-aminés -> Α)
\newunicodechar{Α}{\ensuremath{\alpha}}
\newunicodechar{Β}{\ensuremath{\beta}}
\newunicodechar{Γ}{\ensuremath{\Gamma}}
\newunicodechar{Δ}{\ensuremath{\Delta}}
\newunicodechar{Ε}{\ensuremath{\varepsilon}}
\newunicodechar{Θ}{\ensuremath{\Theta}}
\newunicodechar{Λ}{\ensuremath{\Lambda}}
\newunicodechar{Μ}{\ensuremath{\mu}}
\newunicodechar{Τ}{\ensuremath{\tau}}
\newunicodechar{Φ}{\ensuremath{\Phi}}
\newunicodechar{Ψ}{\ensuremath{\Psi}}
\newunicodechar{Ω}{\ensuremath{\Omega}}
\newunicodechar{↦}{\ensuremath{\mapsto}}
\newunicodechar{∈}{\ensuremath{\in}}
\newunicodechar{∉}{\ensuremath{\notin}}
\newunicodechar{∧}{\ensuremath{\wedge}}
\newunicodechar{⇔}{\ensuremath{\Leftrightarrow}}
\newunicodechar{⟺}{\ensuremath{\Longleftrightarrow}}
\newunicodechar{⇒}{\ensuremath{\Rightarrow}}
\newunicodechar{⟹}{\ensuremath{\Longrightarrow}}
\newunicodechar{∘}{\ensuremath{\circ}}
\newunicodechar{∪}{\ensuremath{\cup}}
\newunicodechar{∩}{\ensuremath{\cap}}
\newunicodechar{≡}{\ensuremath{\equiv}}
\newunicodechar{⊂}{\ensuremath{\subset}}
\newunicodechar{⊥}{\ensuremath{\perp}}
\newunicodechar{∃}{\ensuremath{\exists}}
\newunicodechar{∀}{\ensuremath{\forall}}
\newunicodechar{∛}{\ensuremath{\sqrt[3]{\,}}}
\newunicodechar{∜}{\ensuremath{\sqrt[4]{\,}}}
\newunicodechar{⃗}{\ensuremath{{}^{\to}}}
\newunicodechar{ⁿ}{\ifmmode^{n}\else\textsuperscript{\ensuremath{n}}\fi}
\newunicodechar{ˣ}{\ifmmode^{x}\else\textsuperscript{\ensuremath{x}}\fi}
\newunicodechar{ᵇ}{\ifmmode^{b}\else\textsuperscript{\ensuremath{b}}\fi}
\newunicodechar{ʳ}{\ifmmode^{r}\else\textsuperscript{\ensuremath{r}}\fi}
\newunicodechar{ᵘ}{\ifmmode^{u}\else\textsuperscript{\ensuremath{u}}\fi}
\newunicodechar{ʸ}{\ifmmode^{y}\else\textsuperscript{\ensuremath{y}}\fi}
\newunicodechar{ᵃ}{\ifmmode^{a}\else\textsuperscript{\ensuremath{a}}\fi}
\newunicodechar{ᵉ}{\ifmmode^{e}\else\textsuperscript{\ensuremath{e}}\fi}
\newunicodechar{ᵏ}{\ifmmode^{k}\else\textsuperscript{\ensuremath{k}}\fi}
\newunicodechar{ᵖ}{\ifmmode^{p}\else\textsuperscript{\ensuremath{p}}\fi}
\newunicodechar{ᴺ}{\ifmmode^{N}\else\textsuperscript{\ensuremath{N}}\fi}
\newunicodechar{ᶜ}{\ifmmode^{c}\else\textsuperscript{\ensuremath{c}}\fi}
\newunicodechar{ʰ}{\ifmmode^{h}\else\textsuperscript{\ensuremath{h}}\fi}
\newunicodechar{ᵠ}{\ifmmode^{q}\else\textsuperscript{\ensuremath{q}}\fi}
\newunicodechar{ₙ}{\ifmmode_{n}\else\textsubscript{\ensuremath{n}}\fi}
\newunicodechar{ₐ}{\ifmmode_{a}\else\textsubscript{\ensuremath{a}}\fi}
\newunicodechar{ᵢ}{\ifmmode_{i}\else\textsubscript{\ensuremath{i}}\fi}
\newunicodechar{ᵣ}{\ifmmode_{r}\else\textsubscript{\ensuremath{r}}\fi}
\newunicodechar{ₖ}{\ifmmode_{k}\else\textsubscript{\ensuremath{k}}\fi}
\newunicodechar{ᵦ}{\ifmmode_{\beta}\else\textsubscript{\ensuremath{\beta}}\fi}
\newunicodechar{ₓ}{\ifmmode_{x}\else\textsubscript{\ensuremath{x}}\fi}
\newfontfamily\popdisplay{ArchivoBlack-Regular.ttf}[Path=../../../fonts/]
\newfontfamily\poplogo{SpaceGrotesk-Bold.ttf}[Path=../../../fonts/]
\newfontfamily\popsemi{PlusJakartaSans-SemiBold.ttf}[Path=../../../fonts/]
\newfontfamily\popxbold{PlusJakartaSans-ExtraBold.ttf}[Path=../../../fonts/]
\newfontfamily\popxboldls{PlusJakartaSans-ExtraBold.ttf}[Path=../../../fonts/,LetterSpace=3.0]

\setlist[enumerate]{leftmargin=1.7em,itemsep=5pt,topsep=4pt}
\setlist[itemize]{leftmargin=1.7em,itemsep=4pt,topsep=4pt,label={\color{poporange}\textbullet}}
\setlength{\parindent}{0pt}
\setlength{\parskip}{5pt}
\renewcommand{\arraystretch}{1.25}

% pgfplots : style Pop par défaut
\pgfplotsset{
  every axis/.append style={axis line style={popink,line width=1pt},tick style={popink},
    grid style={popline,line width=0.5pt},label style={font=\small},tick label style={font=\footnotesize}},
  popcurve/.style={poporange,line width=1.8pt,no markers,smooth},
}
% Même style disponible dans un \draw TikZ simple (hors pgfplots)
\tikzset{popcurve/.style={poporange,line width=1.8pt}}
\tikzset{popgrid/.style={popline,very thin}}
\colorlet{popcurve}{poporange} % le nom du style est parfois utilisé comme couleur

% ============================================================
% Pill / squiggle
% ============================================================
\newcommand{\poppill}[3]{%
  \tikz[baseline=-0.55ex]\node[draw=popink,line width=1.2pt,rounded corners=2.6pt,%
    fill=#2,inner xsep=6.5pt,inner ysep=3pt]{%
    \popxboldls\fontsize{7}{7}\selectfont\color{#3}\MakeUppercase{#1}};%
}
\newcommand{\popsquiggle}[1]{%
  \tikz[baseline=-0.4ex]{
    \draw[#1,line width=2.6pt,line cap=round]
      plot [smooth] coordinates {(0,0.09) (0.34,0.01) (0.68,0.21) (1.02,0.01) (1.36,0.10)};
  }%
}
% Mot-clé surligné (marqueur orange sous le texte)
\newcommand{\cle}[1]{{\popxbold\color{popdark}#1}}

% ============================================================
% En-tête / pied
% ============================================================
\pagestyle{fancy}
\fancyhf{}
\renewcommand{\headrulewidth}{0pt}
\renewcommand{\footrulewidth}{0pt}
\lhead{\raisebox{-0.15ex}{\poplogo\fontsize{13}{13}\selectfont\color{popink}OnBuch\color{poporange}+}}
\rhead{\poppill{\DOCMATIERE\ \textbullet\ \DOCNIVEAU\ \textbullet\ Chapter \DOCLECON}{white}{popink}}
\lfoot{\popxbold\fontsize{7.6}{7.6}\selectfont\color{popmuted}\MakeUppercase{Signed course OnBuch+}\ \textbullet\ {\popsemi\DOCTITRE}}
\rfoot{\tikz[baseline=-0.6ex]\node[rounded corners=8.5pt,fill=popink,minimum height=17pt,minimum width=17pt,inner xsep=5pt,inner sep=0]{\popxbold\fontsize{8}{8}\selectfont\color{popcream}\thepage\,/\,\pageref*{LastPage}};}

% ============================================================
% Titres
% ============================================================
\newcommand{\chaptitre}[2]{%
  \begin{tcolorbox}[enhanced,colback=poporange,colframe=popink,boxrule=2.6pt,arc=11pt,
      drop shadow=popink,left=15pt,right=15pt,top=12pt,bottom=11pt,before skip=3pt,after skip=18pt]
    \poppill{#2}{popink}{popcream}\par\vspace{9pt}
    {\raggedright\hyphenpenalty=10000\exhyphenpenalty=10000\popdisplay\fontsize{22}{24}\selectfont\color{popcream}\MakeUppercase{#1}\par}
  \end{tcolorbox}
}
% Section numérotée : pastille orange avec le numéro + titre Archivo
\newcounter{popsec}
\newcommand{\coursec}[1]{%
  \stepcounter{popsec}\setcounter{popsub}{0}%
  \par\vspace{16pt}\noindent
  \begin{minipage}{\linewidth}
  \tikz[baseline=-0.7ex]\node[draw=popink,line width=1.6pt,fill=poporange,rounded corners=4pt,minimum size=22pt,inner sep=0]{\popdisplay\fontsize{12}{12}\selectfont\color{popcream}\thepopsec};%
  \hspace{8pt}{\popdisplay\fontsize{15}{17}\selectfont\color{popink}\MakeUppercase{#1}}\par
  \vspace{4pt}\noindent{\color{poporange}\rule{36pt}{3.6pt}}
  \end{minipage}\par\nopagebreak\vspace{8pt}
}
\newcounter{popsub}
\newcommand{\courssub}[1]{%
  \stepcounter{popsub}%
  \par\vspace{9pt}\noindent{\popxbold\fontsize{11.5}{13}\selectfont\color{popink}{\color{poporange}\thepopsec.\thepopsub}\hspace{6pt}#1}\par\nopagebreak\vspace{4pt}
}

% ============================================================
% Boîtes pédagogiques — même recette : fond blanc, bordure épaisse colorée,
% ombre dure encre, badge plein. Titre optionnel [#1].
% ============================================================
\tcbset{popbase/.style={breakable,enhanced,colback=white,boxrule=2.3pt,arc=9pt,
  shadow={3pt}{-3pt}{0pt}{popink},left=14pt,right=14pt,top=11pt,bottom=11pt,before skip=11pt,after skip=15pt,
  fontupper=\popsemi\fontsize{10.6}{14.5}\selectfont\color{popink},
  fontlower=\popsemi\fontsize{10.6}{14.5}\selectfont\color{popink}}}
% Coche dessinée (le glyphe ✓ n'existe pas dans les polices du document)
\newcommand{\popcheck}[1]{\tikz[baseline=-0.5ex]\draw[#1,line width=1.6pt,line cap=round,line join=round](-0.13,0.0)--(-0.04,-0.09)--(0.13,0.1);}
\providecommand{\coche}{\popcheck{popgreen}}
\providecommand{\resultat}[1]{\par\smallskip\noindent\fcolorbox{popgreen}{popgreenL}{\parbox{\dimexpr\linewidth-2\fboxsep-2\fboxrule\relax}{\textbf{Result.} #1}}\par\smallskip}
\newcommand{\popboxhead}[3]{\poppill{#1}{#2}{white}\ifx\relax#3\relax\else\hspace{6pt}{\popxbold\fontsize{10.6}{12}\selectfont\color{popink}#3}\fi\par\vspace{7pt}}

\newtcolorbox{aretenir}[1][]{popbase,colframe=popgreen,before upper={\popboxhead{Key points}{popgreen}{#1}}}
\newtcolorbox{exemplebox}[1][]{popbase,colframe=poporange,before upper={\popboxhead{Exemple}{poporange}{#1}}}
\newtcolorbox{definition}[1][]{popbase,colframe=popblue,before upper={\popboxhead{Definition}{popblue}{#1}}}
\newtcolorbox{propriete}[1][]{popbase,colframe=poppurple,before upper={\popboxhead{Property}{poppurple}{#1}}}
\newtcolorbox{methode}[1][]{popbase,colframe=popgold,before upper={\popboxhead{Method}{popgold}{#1}}}
\newtcolorbox{attention}[1][]{popbase,colframe=poppink,before upper={\popboxhead{Warning}{poppink}{#1}}}
\newtcolorbox{experience}[1][]{popbase,colframe=popblue,colback=popblueL,before upper={\popboxhead{Experiment}{popblue}{#1}}}
% « Le savais-tu ? » : carte violette pleine (contexte, culture, Cameroun)
\newtcolorbox{savaistu}[1][]{popbase,colback=poppurple,colframe=popink,
  fontupper=\popsemi\fontsize{10.4}{14.2}\selectfont\color{white},
  before upper={\poppill{Did you know?}{white}{poppurple}\ifx\relax#1\relax\else\hspace{6pt}{\popxbold\color{white}#1}\fi\par\vspace{7pt}}}
% Exercice résolu : énoncé en haut, \tcblower puis la solution sur fond crème
\newcounter{popexo}\newcounter{popexr}
\newtcolorbox{exoresolu}[1][]{popbase,colframe=poporange,colbacklower=popcream,
  segmentation style={popink,line width=1.2pt,dashed},
  before upper={\stepcounter{popexr}\popboxhead{Worked example \thepopexr}{poporange}{#1}},
  before lower={\poppill{Solution}{popink}{popcream}\par\vspace{6pt}}}
% Exercice (non résolu en place) — la correction est dans la partie Corrigés
\newtcolorbox{exercice}[1][]{popbase,colframe=popink,
  before upper={\stepcounter{popexo}\popboxhead{Exercise \thepopexo}{popink}{#1}}}
% Corrigé
\newtcolorbox{corrige}[1][]{popbase,colframe=popgreen,colback=popgreenL,
  before upper={\popboxhead{Solution}{popgreen}{#1}}}
% Objectifs / prérequis / bilan
\newtcolorbox{objectifs}{popbase,colframe=popink,colback=poporangeL,
  before upper={\popboxhead{Chapter objectives}{popink}{}}}
\newtcolorbox{prerequis}{popbase,colframe=popink,colback=popblueL,
  before upper={\popboxhead{Prerequisites}{popblue}{}}}
\newtcolorbox{bilan}{popbase,colframe=popink,colback=white,boxrule=2.8pt,
  before upper={\popboxhead{Fiche bilan}{poporange}{L'essentiel à connaître pour l'examen}}}
% Figure encadrée avec légende
\newtcolorbox{popfigure}[1][]{enhanced,breakable=false,colback=white,colframe=popline,boxrule=1.4pt,arc=8pt,
  left=10pt,right=10pt,top=10pt,bottom=8pt,before skip=10pt,after skip=12pt,halign=center,
  lower separated=false,halign lower=center,
  fontlower=\popsemi\fontsize{9}{11.5}\selectfont\color{popmuted},
  #1}
\newcommand{\legende}[1]{\par\vspace{6pt}{\popsemi\fontsize{9}{11.5}\selectfont\color{popmuted}#1}}

% ============================================================
% Signature OnBuch+
% ============================================================
\newcommand{\poplogoinline}[1]{{\poplogo\fontsize{#1}{#1}\selectfont\color{popink}OnBuch\color{poporange}+}}
\newcommand{\signatureonbuch}{%
  \begin{tcolorbox}[enhanced,colback=popink,colframe=popink,boxrule=2.2pt,arc=10pt,
      drop shadow=poporange,left=14pt,right=14pt,top=10pt,bottom=10pt,before skip=0pt,after skip=0pt]
    \tikz[baseline=-0.7ex]\node[circle,fill=popgreen,draw=popcream,line width=1.3pt,minimum size=22pt,inner sep=0]{\popcheck{white}};%
    \hspace{9pt}\begin{minipage}[c]{0.62\linewidth}
      {\popxbold\fontsize{12}{13}\selectfont\color{popcream}Signed\ \textbullet\ {\poplogo OnBuch\color{poporange}+}}\par\vspace{2pt}
      {\raggedright\popsemi\fontsize{8}{10.5}\selectfont\color{popline}\MakeUppercase{OnBuch+ teaching team}\\\MakeUppercase{Follows the official MINESEC syllabus}\par}
    \end{minipage}\hfill
    {\poplogo\fontsize{22}{22}\selectfont\color{popcream}OnBuch\color{poporange}+}
  \end{tcolorbox}%
}

% ============================================================
% Page de garde
% #1 titre · #2 matière · #3 niveau · #4 module · #5 n° leçon · #6 durée ·
% #7 description / objectifs (texte ou itemize)
% ============================================================
\newcommand{\pagegarde}[7]{%
  \thispagestyle{empty}
  \newgeometry{a4paper,margin=1.7cm,top=1.6cm,bottom=1.4cm}%
  \begin{tikzpicture}[remember picture,overlay]
    \fill[poporange!18] ([xshift=-3.2cm,yshift=-3.6cm]current page.north east) circle (4.6cm);
    \fill[poppurple!14] ([xshift=2.4cm,yshift=7.6cm]current page.south west) circle (3.8cm);
  \end{tikzpicture}%
  {\poplogo\fontsize{30}{30}\selectfont\color{popink}OnBuch\color{poporange}+}\hfill
  \raisebox{0.6ex}{\poppill{Signed course \textbullet\ Premium}{popink}{popcream}}\par
  \vspace{1pt}\popsquiggle{poppurple}\par
  \vspace{8pt}
  \begin{tcolorbox}[enhanced,colback=poporange,colframe=popink,boxrule=2.8pt,arc=13pt,
      drop shadow=popink,left=18pt,right=18pt,top=12pt,bottom=13pt,after skip=10pt]
    \poppill{#2 \textbullet\ #3}{popink}{popcream}\hspace{5pt}\poppill{#4}{white}{popink}\par\vspace{8pt}
    {\popdisplay\fontsize{14}{14}\selectfont\color{popink}CHAPTER #5}\par\vspace{4pt}
    {\raggedright\hyphenpenalty=10000\exhyphenpenalty=10000\popdisplay\fontsize{25}{27}\selectfont\color{popcream}\MakeUppercase{#1}\par}
    \vspace{8pt}
    {\popxbold\fontsize{9.5}{9.5}\selectfont\color{popink}\MakeUppercase{Suggested time: #6}}
  \end{tcolorbox}
  \begin{tcolorbox}[enhanced,colback=white,colframe=popink,boxrule=2.2pt,arc=10pt,
      drop shadow=popink,left=16pt,right=16pt,top=10pt,bottom=10pt,after skip=8pt]
    \poppill{In this chapter}{poppurple}{white}\par\vspace{5pt}
    \popsemi\fontsize{9.3}{12.6}\selectfont\color{popink}#7
  \end{tcolorbox}
  \noindent\begin{minipage}[t]{0.487\linewidth}
    \begin{tcolorbox}[enhanced,colback=popgreen,colframe=popink,boxrule=2.2pt,arc=10pt,
        drop shadow=popink,left=11pt,right=11pt,top=10pt,bottom=10pt,height=3.1cm,valign=center]
      \tcbox[colback=white,colframe=popink,boxrule=1.3pt,arc=5pt,boxsep=0pt,nobeforeafter,
        left=3pt,right=3pt,top=3pt,bottom=3pt,valign=center]{\qrcode[height=1.9cm]{https://play.google.com/store/apps/details?id=com.luvvix.onbuch&hl=en}}%
      \hspace{8pt}\begin{minipage}[c]{0.5\linewidth}
        {\popdisplay\fontsize{11}{12.5}\selectfont\color{white}\MakeUppercase{Download OnBuch}}\par\vspace{4pt}
        {\raggedright\popsemi\fontsize{8.4}{11}\selectfont\color{white}Free \textbullet\ Google Play\\AI tutor, past papers, results\par}
      \end{minipage}
    \end{tcolorbox}
  \end{minipage}\hfill
  \begin{minipage}[t]{0.487\linewidth}
    \begin{tcolorbox}[enhanced,colback=poppurple,colframe=popink,boxrule=2.2pt,arc=10pt,
        drop shadow=popink,left=11pt,right=11pt,top=10pt,bottom=10pt,height=3.1cm,valign=center]
      \tikz[baseline=-0.7ex]\node[circle,fill=white,draw=popink,line width=1.3pt,minimum size=1.6cm,inner sep=0]{\popdisplay\fontsize{24}{24}\selectfont\color{poppurple}+};%
      \hspace{8pt}\begin{minipage}[c]{0.6\linewidth}
        {\popdisplay\fontsize{11}{12.5}\selectfont\color{white}\MakeUppercase{Go OnBuch+}}\par\vspace{4pt}
        {\raggedright\popsemi\fontsize{8.4}{11}\selectfont\color{white}All signed courses, unlimited AI tutor, PDF sheets — OnBuch+ tab in the app\par}
      \end{minipage}
    \end{tcolorbox}
  \end{minipage}
  \vspace{8pt}
  \popcontenu
  \vfill
  \signatureonbuch
  \restoregeometry
  \newpage
}

% Rangée « Ce cours contient » (page de garde)
\newcommand{\poptile}[3]{% #1 couleur #2 gros texte #3 légende
  \begin{tcolorbox}[enhanced,colback=#1,colframe=popink,boxrule=2pt,arc=9pt,shadow={2.5pt}{-2.5pt}{0pt}{popink},
      left=6pt,right=6pt,top=9pt,bottom=9pt,halign=center,valign=center,height=1.75cm,width=\linewidth,nobeforeafter]
    {\popdisplay\fontsize{12.5}{13}\selectfont\color{white}\MakeUppercase{#2}}\par\vspace{3pt}
    {\centering\popsemi\fontsize{7.8}{9.5}\selectfont\color{white}#3\par}
  \end{tcolorbox}}
\newcommand{\popcontenu}{%
  \par\noindent{\popxboldls\fontsize{7.5}{7.5}\selectfont\color{popmuted}THIS SIGNED COURSE CONTAINS}\par\vspace{7pt}
  \noindent\begin{minipage}[t]{0.235\linewidth}\poptile{popblue}{Course}{complete, illustrated, step by step}\end{minipage}\hfill
  \begin{minipage}[t]{0.235\linewidth}\poptile{poporange}{Methods}{and worked examples}\end{minipage}\hfill
  \begin{minipage}[t]{0.235\linewidth}\poptile{poppink}{Exercises}{exam style + solutions}\end{minipage}\hfill
  \begin{minipage}[t]{0.235\linewidth}\poptile{popgreen}{Summary sheet}{the essentials to remember}\end{minipage}\par}

% Sommaire
\newtcolorbox{sommaire}{enhanced,colback=white,colframe=popink,boxrule=2.2pt,arc=10pt,
  drop shadow=popink,left=18pt,right=18pt,top=13pt,bottom=13pt,before skip=6pt,after skip=20pt,
  before upper={\poppill{Contents}{poppurple}{white}\par\vspace{9pt}\popsemi\fontsize{10.6}{14.5}\selectfont\color{popink}}}

% Signature de fin de cours — poussée tout en bas de la dernière page
\newcommand{\findecours}{%
  \par\vfill\nopagebreak
  \begin{center}\popsquiggle{poporange}\end{center}\vspace{4pt}
  \signatureonbuch
}
\AtBeginDocument{\def\llbracket{\mathopen{[\mkern-3.5mu[}}\def\rrbracket{\mathclose{]\mkern-3.5mu]}}} % intervalles d'entiers (glyphe ⟦ absent de la police)
% Glyphes absents de Fira Math : redéfinis à partir de glyphes disponibles
\AtBeginDocument{%
  \def\setminus{\mathbin{\backslash}}\let\smallsetminus\setminus
  \def\vdots{\mathord{\vbox{\baselineskip3.2pt\lineskiplimit0pt\kern4pt\hbox{.}\hbox{.}\hbox{.}}}}%
  \def\ddots{\mathinner{\mkern1mu\raise6.5pt\hbox{.}\mkern2mu\raise3.6pt\hbox{.}\mkern2mu\raise0.7pt\hbox{.}\mkern1mu}}%
  \def\triangle{\mathord{\scalebox{0.9}{$\symup{\Delta}$}}}%
  \def\nabla{\mathord{\raisebox{\depth}{\scalebox{1}[-1]{$\symup{\Delta}$}}}}%
  \def\checkmark{\popcheck{popdark}}%
  \def\star{\mathbin{\ast}}\def\bigstar{\mathord{\ast}}%
  \def\diamond{\mathbin{\scalebox{0.6}{$\Diamond$}}}%
  \def\bigcup{\mathop{\mathchoice{\raisebox{-0.3em}{\scalebox{1.7}{$\cup$}}}{\scalebox{1.25}{$\cup$}}{\cup}{\cup}}\displaylimits}%
  \def\bigcap{\mathop{\mathchoice{\raisebox{-0.3em}{\scalebox{1.7}{$\cap$}}}{\scalebox{1.25}{$\cap$}}{\cap}{\cap}}\displaylimits}%
  \def\lesseqqgtr{\mathrel{\lessgtr}}\def\bigoplus{\mathop{\oplus}\displaylimits}%
}
\providecommand{\ssi}{if and only if} % abréviation fréquente
\AtBeginDocument{\providecommand{\wideparen}{\overparen}} % arc de cercle

% Fonctions usuelles que les modèles utilisent sans que amsmath les fournisse
\providecommand{\arsinh}{\operatorname{arsinh}}\providecommand{\arcosh}{\operatorname{arcosh}}
\providecommand{\artanh}{\operatorname{artanh}}\providecommand{\arsech}{\operatorname{arsech}}
\providecommand{\arcsch}{\operatorname{arcsch}}\providecommand{\arcoth}{\operatorname{arcoth}}
\providecommand{\arcsinh}{\operatorname{arsinh}}\providecommand{\arccosh}{\operatorname{arcosh}}
\providecommand{\arctanh}{\operatorname{artanh}}\providecommand{\sech}{\operatorname{sech}}
\providecommand{\csch}{\operatorname{csch}}\providecommand{\arcsec}{\operatorname{arcsec}}
\providecommand{\arccsc}{\operatorname{arccsc}}\providecommand{\arccot}{\operatorname{arccot}}
\providecommand{\sgn}{\operatorname{sgn}}\providecommand{\lcm}{\operatorname{lcm}}
\providecommand{\Var}{\operatorname{Var}}\providecommand{\Cov}{\operatorname{Cov}}

%% FIGFIX-BEGIN v5 — correctifs globaux des figures (docs/onbuch-generation/tools/figfix)
\usepackage{adjustbox}\usepackage{etoolbox}\tcbuselibrary{raster}
% toute figure TikZ/pgfplots plus large que la ligne est réduite à la largeur disponible
\BeforeBeginEnvironment{tikzpicture}{\begin{adjustbox}{max width=\linewidth}}
\AfterEndEnvironment{tikzpicture}{\end{adjustbox}}
% légendes pgfplots lisibles par-dessus les courbes
\pgfplotsset{every axis/.append style={legend style={fill=white,fill opacity=0.92,text opacity=1,draw=black!15,inner sep=3pt}}}
% étiquettes d'arêtes (syntaxe edge["..."]) avec fond blanc
\tikzset{every edge quotes/.append style={fill=white,inner sep=1.2pt}}
%% FIGFIX-END
\begin{document}

\pagegarde{ECONOMIC GEOLOGY: Economic Minerals}{Geology}{Upper Sixth}{Upper Sixth — Geology}{5}{12 periods}{This chapter introduces the concept of economic minerals, the distinction between resources and reserves, and the classification, formation and occurrence of mineral deposits. Using Cameroonian examples such as the bauxite of Minim-Martap, the iron of Mbalam and the gold of the East Region, students learn how geologists evaluate and locate the mineral wealth of a country.
\vspace{4pt}
\begin{multicols}{2}\small
\begin{itemize}
\item By the end of this chapter I can define an economic mineral and distinguish it from a gangue mineral and an ordinary rock.
\item By the end of this chapter I can explain the terms ore, grade, cut-off grade and tonnage with simple calculations.
\item By the end of this chapter I can distinguish mineral resources from mineral reserves and explain the factors that convert one into the other.
\item By the end of this chapter I can classify mineral deposits according to their morphology, their host rock and their mode of origin.
\item By the end of this chapter I can describe the magmatic, hydrothermal, sedimentary and weathering processes that form ore deposits.
\item By the end of this chapter I can relate the occurrence of major deposits to the geological context of Cameroon and the world.
\end{itemize}\end{multicols}}

\begin{sommaire}
\begin{multicols}{2}
\begin{enumerate}
\item The Notion of Economic Mineral
\item Resources and Reserves
\item Classification of Mineral Deposits
\item Formation of Ore Deposits I: Magmatic and Hydrothermal Processes
\item Formation of Ore Deposits II: Sedimentary, Weathering and Metamorphic Processes
\item Occurrence of Mineral Deposits: From the World to Cameroon
\item Integration activity
\item Methods and worked examples
\item Exercises
\item Solutions to the exercises
\item Summary sheet
\end{enumerate}
\end{multicols}
\end{sommaire}

\begin{objectifs}
\begin{itemize}
\item By the end of this chapter I can define an economic mineral and distinguish it from a gangue mineral and an ordinary rock.
\item By the end of this chapter I can explain the terms ore, grade, cut-off grade and tonnage with simple calculations.
\item By the end of this chapter I can distinguish mineral resources from mineral reserves and explain the factors that convert one into the other.
\item By the end of this chapter I can classify mineral deposits according to their morphology, their host rock and their mode of origin.
\item By the end of this chapter I can describe the magmatic, hydrothermal, sedimentary and weathering processes that form ore deposits.
\item By the end of this chapter I can relate the occurrence of major deposits to the geological context of Cameroon and the world.
\item By the end of this chapter I can interpret simple maps, sections and grade-tonnage data to assess the economic potential of a deposit.
\item By the end of this chapter I can discuss the economic and environmental stakes of mineral exploitation in Cameroon.
\end{itemize}
\end{objectifs}

\begin{prerequis}
\begin{itemize}
\item Minerals and rocks: definition of a mineral, common rock-forming minerals (Lower Sixth).
\item The rock cycle and the three great families of rocks (igneous, sedimentary, metamorphic).
\item Basic notions of plate tectonics and magmatic processes.
\item Weathering and sedimentation processes (Lower Sixth, external geodynamics).
\item Simple percentage and unit-conversion calculations from Mathematics.
\end{itemize}
\end{prerequis}

\newpage

\coursec{The Notion of Economic Mineral}

In 2023, a young entrepreneur from Batouri in the East Region of Cameroon sold a few kilograms of gold-bearing gravel for less than the price of a bag of cement, while a foreign company exporting the same gold earned millions of CFA francs. Why does the same mineral have such different values depending on who holds it and how much of it exists underground? How do geologists decide whether a hill of rock is worthless or a fortune? And why does Cameroon, rich in bauxite at Minim-Martap and iron at Mbalam, still import most of its metals? This chapter gives you the geological keys to answer these questions.

\courssub{From Mineral to Ore: the Basic Vocabulary}

In 2023, a young entrepreneur from Batouri in the East Region of Cameroon sold a few kilograms of gold-bearing gravel for less than the price of a bag of cement, while a foreign company exporting the same gold earned millions of CFA francs. Why does the same mineral have such different values depending on who holds it and how much of it exists underground? How do geologists decide whether a hill of rock is worthless or a fortune? And why does Cameroon, rich in bauxite at Minim-Martap and iron at Mbalam, still import most of its metals? This chapter gives you the geological keys to answer these questions.

\courssub{From Mineral to Ore: the Basic Vocabulary}

Imagine you find a rock. Is it a mineral? Is it a metal? Is it an ore? Is it a metal? Is it an ore? Is it an ore? Is it an ore? Is it an ore? Is it an ore? Is it an ore? Is it an ore? Is it an ore? Is it an ore? Is it an ore? Is it an ore? Is it an ore? Is it an ore? Is it an ore? Is it an ore? Is it an ore? Is it an ore? Is it an ore? Is it an ore? Is it an ore? Is it an ore? Is it an ore? Is it an ore? Is it an ore? Is it an ore? Is it an ore? Is it an ore? Is it an ore? Is it an ore? Is it an ore? Is it an ore? Is it an ore

\coursec{Resources and Reserves}

In the previous section we established that an \cle{economic mineral} is a mineral substance whose concentration in the Earth's crust is high enough, and whose extraction is profitable enough, to be exploited. But between ``there is iron in this hill'' and ``this iron mine is profitable'', there is a long road of measurement, calculation and decision. Mining companies, banks and governments do not speak vaguely of ``minerals in the ground'': they use two precise, internationally standardised words --- \cle{resource} and \cle{reserve}. Mastering the difference between them is one of the most frequently tested points of this chapter at the GCE Advanced Level.

\courssub{The Two Fundamental Notions}

\textbf{Intuition first.}\par
Imagine a farmer who knows that mangoes grow somewhere on his large farm. He has \emph{seen} mangoes from a distance in some areas (low confidence), counted the trees in other areas (medium confidence), and weighed and priced the fruit on the trees nearest his house (high confidence). Only the fruit he can actually pick, transport and sell at a profit this season constitutes his real ``harvestable stock''. The rest exists, certainly, but it is not yet money in the pocket. In the same way, geologists distinguish what is \emph{probably there} (a resource) from what is \emph{demonstrably mineable at a profit} (a reserve).

\begin{definition}[Mineral resource]
A \cle{mineral resource} is a concentration or occurrence of solid mineral material in or on the Earth's crust, in such form, grade (quality) and quantity that there are \textbf{reasonable prospects of eventual economic extraction}. Its location, tonnage and characteristics are known, estimated or interpreted from geological evidence and sampling.
\end{definition}

\begin{definition}[Mineral reserve (ore reserve)]
A \cle{mineral reserve} is the \textbf{economically mineable part} of a measured or indicated mineral resource. It is demonstrated by a \emph{feasibility study} which takes into account mining, metallurgical, economic, marketing, legal, environmental and governmental factors. A reserve is what a company can legally and profitably extract \emph{today}.
\end{definition}

\begin{attention}[Resource $\neq$ reserve]
In GCE answers, never use the two words as synonyms. A \textbf{reserve is always a resource}, but a resource is \textbf{not always a reserve}: it becomes one only when economic extraction is demonstrated. Writing ``Cameroon has huge iron reserves at Mbalam'' when the deposit is still at the resource stage is a factual error that examiners penalise.
\end{attention}

\courssub{Classification by Degree of Geological Confidence}

As exploration progresses (outcrop mapping $\rightarrow$ geophysics $\rightarrow$ trenching $\rightarrow$ borehole drilling $\rightarrow$ bulk sampling), the geologist's confidence in the tonnage and grade increases. Three categories of resource are distinguished:

\begin{definition}[Inferred, indicated and measured resources]
\begin{itemize}
\item \cle{Inferred resource}: estimated from geological evidence and limited sampling; continuity of the orebody is \emph{assumed} but not verified. Low confidence.
\item \cle{Indicated resource}: sampling (boreholes, trenches) is dense enough to assume continuity between observation points; tonnage and grade are estimated with \emph{reasonable} confidence.
\item \cle{Measured resource}: sampling is so dense (close-spaced drilling, underground workings) that tonnage, grade, shape and mineral content are \emph{well established}. Highest confidence.
\end{itemize}
\end{definition}

\begin{popfigure}
\begin{center}
\begin{tikzpicture}[scale=1.0]
% pyramid layers
\fill[popblueL] (-4.5,0) -- (4.5,0) -- (3.2,1.4) -- (-3.2,1.4) -- cycle;
\fill[popgreenL] (-3.2,1.4) -- (3.2,1.4) -- (2.0,2.8) -- (-2.0,2.8) -- cycle;
\fill[popgold!40] (-2.0,2.8) -- (2.0,2.8) -- (0.9,4.2) -- (-0.9,4.2) -- cycle;
\fill[poporange!60] (-0.9,4.2) -- (0.9,4.2) -- (0,5.4) -- cycle;
\node at (0,0.7) {\textbf{Inferred resource}};
\node at (0,2.1) {\textbf{Indicated resource}};
\node at (0,3.5) {\textbf{Measured resource}};
\node[text=white] at (0,4.8) {\textbf{Probable reserve}};
\node[text=white] at (0,4.5) {\textbf{Proved reserve}};
% side annotations
\draw[->, thick, popdark] (5.2,0.2) -- (5.2,5.2) node[midway, right, align=center, text width=3.4cm] {Increasing geological confidence and decreasing tonnage};
\draw[->, thick, popdark] (-5.2,5.2) -- (-5.2,0.2) node[midway, left, align=center, text width=3.2cm] {Increasing tonnage, decreasing certainty};
\end{tikzpicture}
\end{center}
\legende{The resource--reserve pyramid: the largest tonnages are known with the least confidence (inferred), while the economically mineable, best-known part (proved and probable reserves) is the smallest.}
\end{popfigure}

\courssub{Classification by Economic Viability}

Geological confidence alone is not enough: the ore must be mineable \emph{at a profit}. This gives the second axis of classification:

\begin{definition}[Proved and probable reserves]
\begin{itemize}
\item \cle{Proved (proven) reserve}: the economically mineable part of a \textbf{measured} resource; highest confidence in both geology and economics.
\item \cle{Probable reserve}: the economically mineable part of an \textbf{indicated} (and sometimes measured) resource; confidence is lower than for a proved reserve but extraction is still judged profitable.
\end{itemize}
Resources that are \emph{almost} economic are called \cle{marginal} (borderline profitable); those that cannot be mined profitably under present conditions are \cle{sub-economic} resources.
\end{definition}

\begin{popfigure}
\begin{center}
\begin{tikzpicture}[scale=1.05]
% axes
\draw[->, thick] (0,0) -- (8.6,0) node[below left=2pt and -30pt] {Increasing economic viability};
\draw[->, thick] (0,0) -- (0,5.6) node[rotate=90, above left=14pt and -58pt] {Increasing geological certainty};
% grid cells (3 columns x 2 rows)
\draw[thick] (0,0) rectangle (8.4,5.2);
\draw[thick] (2.8,0) -- (2.8,5.2);
\draw[thick] (5.6,0) -- (5.6,5.2);
\draw[thick] (0,2.6) -- (8.4,2.6);
% cells: top row (measured/indicated)
\fill[poporange!55] (5.6,2.6) rectangle (8.4,5.2);
\fill[popgold!45] (2.8,2.6) rectangle (5.6,5.2);
\fill[popblueL] (0,2.6) rectangle (2.8,5.2);
% cells: bottom row (inferred)
\fill[popblueL] (5.6,0) rectangle (8.4,2.6);
\fill[popblueL] (2.8,0) rectangle (5.6,2.6);
\fill[popblueL] (0,0) rectangle (2.8,2.6);
% labels
\node[align=center] at (7.0,4.4) {\textbf{PROVED /}\\ \textbf{PROBABLE}\\ \textbf{RESERVES}};
\node[align=center] at (4.2,3.9) {Marginal\\ reserves};
\node[align=center] at (1.4,3.9) {Measured /\\ indicated\\ resources};
\node[align=center] at (7.0,1.3) {Marginal /\\ sub-economic\\ resources};
\node[align=center] at (4.2,1.3) {Sub-economic\\ resources};
\node[align=center] at (1.4,1.3) {Inferred\\ resources};
% row/column captions
\node[rotate=90] at (-0.55,3.9) {\small demonstrated};
\node[rotate=90] at (-0.55,1.3) {\small inferred};
\node at (1.4,-0.45) {\small sub-economic};
\node at (4.2,-0.45) {\small marginal};
\node at (7.0,-0.45) {\small economic};
\end{tikzpicture}
\end{center}
\legende{The ``McKelvey box'': resources and reserves classified simultaneously by geological certainty (vertical axis) and economic viability (horizontal axis). Only the upper-right zone constitutes true reserves.}
\end{popfigure}

\begin{aretenir}[The McKelvey box]
The \cle{McKelvey diagram} combines the two classifications: as one moves \emph{upwards}, geological certainty increases (inferred $\rightarrow$ indicated/measured); as one moves \emph{rightwards}, economic viability increases (sub-economic $\rightarrow$ marginal $\rightarrow$ economic). \textbf{Reserves occupy the upper-right corner.} A deposit can migrate inside the box without any new drilling, simply because prices or technology change.
\end{aretenir}

\courssub{From Resource to Reserve: the Conversion Factors}

What decides whether a resource crosses the boundary into the reserve zone? Five groups of factors:

\begin{enumerate}
\item \textbf{Metal and mineral prices.} If the world price of iron or gold rises, lower-grade ore that was previously worthless becomes profitable: the cut-off grade falls and the reserve grows.
\item \textbf{Mining and processing technology.} Cheaper open-pit machinery, better ore-dressing (e.g.\ magnetic separation of magnetite) or new hydrometallurgical methods lower costs and enlarge reserves.
\item \textbf{Infrastructure.} Railways, roads, deep-sea ports, water and electricity supply often decide everything for bulky, low-value commodities such as iron ore.
\item \textbf{Energy costs.} Smelting and crushing are energy-hungry; cheap hydroelectricity (a Cameroonian asset) improves viability.
\item \textbf{Legislation and environmental constraints.} Mining codes, taxes, land rights and environmental protection can either encourage or block conversion of a resource into a reserve.
\end{enumerate}

\begin{attention}[A price rise can ``create'' a reserve]
A rise in the world price of a metal can transform a sub-economic resource into a reserve \textbf{without any new discovery and without moving a single rock}. Reserves are therefore an \emph{economic} stock, not a fixed physical quantity. Conversely, a price collapse can wipe out reserves overnight.
\end{attention}

\courssub{The Dynamic Nature of Reserves and Reserve Life}

\begin{propriete}[Reserves are a dynamic stock]
Reserves \textbf{shrink} through extraction (depletion by mining) and \textbf{grow} through new discoveries, rises in commodity prices and technological progress. The ``end of a mineral'' is therefore rarely a physical exhaustion; it is usually an economic limit that moves with time. (Statement examinable; no mathematical proof required.)
\end{propriete}

A practical consequence is the notion of \cle{reserve life}: the number of years a deposit (or a country, or the world) can continue producing at the current rate.

\begin{methode}[Calculating reserve life]
\begin{enumerate}
\item Identify the \textbf{reserve tonnage} $R$ (in tonnes, t) from the data.
\item Identify the \textbf{annual production} $P$ (in t/yr), making sure the units match.
\item Compute: $\text{Reserve life} = \dfrac{R}{P}$ (in years).
\item Comment critically: the result assumes constant production, constant prices and no new discovery --- a simplification.
\end{enumerate}
\end{methode}

\begin{exoresolu}[Reserve life of an iron deposit]
Statement: An iron deposit in the South Region of Cameroon has a proved reserve of $2.0 \times 10^{8}\ \mathrm{t}$ of saleable ore. The planned mine will produce $1.0 \times 10^{7}\ \mathrm{t}$ per year. (a) Calculate the reserve life. (b) Explain two reasons why the true mine life may differ from this figure.
\tcblower
\textbf{Solution.}
(a) Applying the method:
\[
\text{Reserve life} = \frac{R}{P} = \frac{2.0 \times 10^{8}\ \mathrm{t}}{1.0 \times 10^{7}\ \mathrm{t\,yr^{-1}}} = 20\ \text{years}.
\]
(b) The true life may be \emph{longer} because further drilling may convert inferred resources into reserves, or a rise in iron prices may make lower-grade ore economic; it may be \emph{shorter} if production is increased above the planned rate or if part of the reserve proves impossible to recover (sterile zones, environmental restrictions). Reserve life is therefore an estimate valid only under stated conditions.
\end{exoresolu}

\courssub{Case Study: the Mbalam--Nabeba Iron Project}

\begin{experience}[Guided case study --- Mbalam--Nabeba (East Region, Cameroon)]
The Mbalam--Nabeba iron deposit, straddling the Cameroon--Congo border, contains large tonnages of high-grade hematite and itabirite mineralisation hosted in Precambrian rocks of the Congo craton margin. Geological exploration established a \textbf{resource} of several hundred million tonnes of iron ore. Yet for years this resource could not be declared a full \textbf{reserve}, for one decisive reason: \textbf{infrastructure}. Iron ore is a bulky, low-value-per-tonne commodity; it cannot be exported profitably by lorry over hundreds of kilometres of forest roads. The project requires:
\begin{itemize}
\item a heavy-haul \textbf{railway} of roughly 500 km from the mine to the coast;
\item a \textbf{deep-sea port} (mineral terminal near Kribi) capable of loading large bulk carriers;
\item power and water supplies for the processing plant.
\end{itemize}
Only when financing and agreements for the railway and port are secured does the feasibility study demonstrate economic extraction, allowing the resource to be reclassified as a reserve. \textbf{Conclusion:} in Cameroon, as elsewhere, the passage from resource to reserve depends as much on engineers, bankers and ministers as on geologists.
\end{experience}

\begin{savaistu}[Did you know?]
The Mbalam--Nabeba project illustrates a general African paradox: the continent is extraordinarily rich in mineral \emph{resources} (iron, bauxite, cobalt, gold) but converts them into \emph{reserves} --- and into revenue --- slowly, because railways, ports and energy are lacking. Cameroon's bauxite at Minim-Martap and Ngaoundal, among the largest untapped resources in the world, faces exactly the same constraint.
\end{savaistu}

\begin{exercice}[Resources, reserves and the McKelvey box]
\begin{enumerate}
\item Define \emph{mineral resource} and \emph{mineral reserve}, and state the essential difference between them.
\item A gold prospect is known from only three boreholes 400 m apart. To which confidence category does its resource belong? Justify.
\item The world price of cobalt doubles. Using the McKelvey box, explain what happens to a marginal cobalt resource.
\item A mine holds a probable reserve of $5.0 \times 10^{7}\ \mathrm{t}$ and produces $2.5 \times 10^{6}\ \mathrm{t\,yr^{-1}}$. Calculate its reserve life and state one assumption behind your answer.
\end{enumerate}
\end{exercice}

\begin{corrige}[Exercise 1]
\begin{enumerate}
\item A mineral resource is a concentration of mineral material with reasonable prospects of eventual economic extraction; a mineral reserve is the economically mineable part of a measured or indicated resource, demonstrated by a feasibility study. The difference: a reserve is proven profitable \emph{now}, a resource is only \emph{potentially} economic.
\item Inferred resource: sampling is too sparse to confirm continuity of the orebody between boreholes; continuity is assumed, not verified.
\item The resource moves rightwards in the McKelvey box: higher price lowers the effective cut-off grade and raises revenue per tonne, so a marginal resource can become economic and be reclassified as a probable/proved reserve --- with no new discovery.
\item Reserve life $= \dfrac{5.0 \times 10^{7}}{2.5 \times 10^{6}} = 20$ years. Assumption: production rate, prices and the reserve itself remain constant (no new discoveries, no reclassification).
\end{enumerate}
\end{corrige}

\begin{aretenir}[Key points of the section]
\begin{itemize}
\item \textbf{Resource} = potentially economic concentration; \textbf{reserve} = demonstrated economically mineable part of a measured/indicated resource.
\item Confidence increases: inferred $\rightarrow$ indicated $\rightarrow$ measured; viability increases: sub-economic $\rightarrow$ marginal $\rightarrow$ proved/probable reserve.
\item Prices, technology, infrastructure, energy and legislation convert resources into reserves --- or the reverse.
\item Reserves are dynamic: depleted by mining, renewed by discovery, price rises and technology.
\item Reserve life $=$ reserve tonnage $\div$ annual production; the Mbalam--Nabeba case shows how railway and port condition the resource-to-reserve transition in Cameroon.
\end{itemize}
\end{aretenir}

\coursec{Classification of Mineral Deposits}

In the previous section we learned how to quantify what we find: \cle{resources} tell us what exists in the ground, while \cle{reserves} tell us what we can mine profitably today. But before we can even estimate tonnage and grade, the exploration geologist must answer a more fundamental question: \emph{what kind of deposit are we looking at?} The answer dictates the exploration strategy, the mining method, the processing route and ultimately the economic viability. This section builds the classification toolkit that every geologist — and every GCE candidate — must master. It corresponds to \textbf{Lessons 91 and 92} of the official syllabus.

\courssub{Why Classify Mineral Deposits?}

Imagine you are handed a geological map showing a mineralised zone. Without a classification framework you would be facing a chaotic list of observations: \emph{``quartz veins in granite'', ``banded iron layers in schist'', ``gold grains in river sand''}. Classification imposes order. It groups deposits that share \textbf{geometry}, \textbf{genetic origin}, \textbf{relationship to the host rock} or \textbf{commodity type}. Each criterion answers a different practical question:

\begin{itemize}
  \item \textbf{Morphology (shape)} $\rightarrow$ determines the \emph{mining method} (open pit vs underground, selective vs bulk mining).
  \item \textbf{Host-rock relation (syngenetic vs epigenetic)} $\rightarrow$ guides \emph{exploration targeting} (look along stratigraphic horizons vs cross-cutting structures).
  \item \textbf{Genetic process} $\rightarrow$ predicts \emph{associated minerals, alteration halos, geochemical signatures} and \emph{down-dip continuity}.
  \item \textbf{Commodity group} $\rightarrow$ links to \emph{market dynamics, processing technology, strategic importance}.
\end{itemize}

\begin{aretenir}[The Golden Rule of Classification]
A single deposit can (and must) be described using \textbf{all four criteria} simultaneously. For example, the \emph{Mbalam iron deposit} (Cameroon) is: \textbf{morphologically} stratiform; \textbf{genetically} sedimentary (BIF); \textbf{host-relation} syngenetic; \textbf{commodity} ferrous metallic. Never reduce a deposit to just one label.
\end{aretenir}

\courssub{Classification by Morphology (Shape and Geometry)}

Morphology is the most immediately observable characteristic in the field or on a cross-section. It controls the \cle{ore-body continuity} and the \cle{stripping ratio}.

\begin{definition}[Stratiform Deposit]
A tabular ore body that \emph{parallels the bedding or layering} of the host rock. Thickness is small compared to lateral extent. Typical of sedimentary exhalative (SEDEX), volcanogenic massive sulphide (VMS) and banded iron formations (BIF).
\end{definition}

\begin{definition}[Stratabound Deposit]
An ore body \emph{confined to a specific stratigraphic unit} but not necessarily parallel to bedding; it may be discordant within that unit. Example: Mississippi Valley Type (MVT) lead-zinc in a specific carbonate horizon.
\end{definition}

\begin{definition}[Vein and Lode]
A \cle{vein} is a sheet-like body filling a fracture or fault; a \cle{lode} is a mineralised fault zone (often multiple veins + altered wall rock). Both are \cle{discordant} — they cut across host-rock layering.
\end{definition}

\begin{definition}[Stockwork and Dissemination]
A \cle{stockwork} is a network of closely spaced, small veins; a \cle{dissemination} consists of fine-grained ore minerals scattered through the host rock. Both imply \emph{bulk-mining} methods (porphyry Cu, epithermal Au).
\end{definition}

\begin{definition}[Pipe, Pod and Lens]
\cle{Pipe}: cylindrical, near-vertical body (kimberlite, breccia pipe). \cle{Pod}: small, irregular, blob-like. \cle{Lens}: lenticular, pinching out in all directions.
\end{definition}

\begin{definition}[Placer Deposit]
Concentration of heavy, resistant minerals (gold, cassiterite, diamond, ilmenite) in \emph{unconsolidated sediments} — river beds (alluvial), beaches, palaeo-channels. Morphology follows paleo-drainage.
\end{definition}

\begin{popfigure}
\begin{tikzpicture}[scale=0.9, every node/.style={font=\small}]
% Stratiform
\begin{scope}[xshift=0cm]
\draw[thick] (0,0) -- (3,0) node[midway,above] {Surface};
\foreach \y in {0.5,1.0,1.5,2.0} {
  \draw[thin,gray] (0,\y) -- (3,\y);
}
\draw[poporange,thick,fill=poporange!30] (0.3,1.0) rectangle (2.7,1.2);
\node[poporange,align=center] at (1.5,1.6) {Stratiform\\ore layer};
\node[align=center] at (1.5,-0.5) {(a) Stratiform};
\end{scope}

% Vein cutting strata
\begin{scope}[xshift=4cm]
\draw[thick] (0,0) -- (3,0) node[midway,above] {Surface};
\foreach \y in {0.5,1.0,1.5,2.0} {
  \draw[thin,gray] (0,\y) -- (3,\y);
}
\draw[popblue,thick,fill=popblue!20] (1.4,0.2) -- (1.6,0.2) -- (1.3,2.2) -- (1.1,2.2) -- cycle;
\node[popblue,align=center] at (1.5,2.5) {Quartz vein\\cross-cutting};
\node[align=center] at (1.5,-0.5) {(b) Vein / Lode};
\end{scope}

% Stockwork
\begin{scope}[xshift=8cm]
\draw[thick] (0,0) -- (3,0) node[midway,above] {Surface};
\foreach \y in {0.5,1.0,1.5,2.0} {
  \draw[thin,gray] (0,\y) -- (3,\y);
}
\foreach \i in {1,...,30} {
  \pgfmathsetmacro{\xi}{0.2+rand*2.6}
  \pgfmathsetmacro{\yi}{0.2+rand*1.8}
  \pgfmathsetmacro{\ang}{rand*180}
  \pgfmathsetmacro{\len}{0.1+rand*0.3}
  \draw[popgreen!80!black,thin] (\xi,\yi) -- ++(\ang:\len);
}
\node[popgreen!80!black,align=center] at (1.5,2.3) {Stockwork\\veinlets};
\node[align=center] at (1.5,-0.5) {(c) Stockwork};
\end{scope}

% Placer in river bed
\begin{scope}[xshift=0cm, yshift=-3.5cm]
\draw[popblue!50!black,thick] (0,0.5) .. controls (1,0.2) and (2,0.3) .. (3,0.4);
\draw[popblue!50!black,thick] (0,0) .. controls (1,-0.1) and (2,0) .. (3,0.1);
\foreach \i in {1,...,40} {
  \pgfmathsetmacro{\xi}{0.3+rand*2.4}
  \pgfmathsetmacro{\yi}{0.05+rand*0.35}
  \fill[popgold] (\xi,\yi) circle (0.04);
}
\node[popgold,align=center] at (1.5,1.0) {Heavy minerals\\(gold, cassiterite)};
\node[align=center] at (1.5,-0.5) {(d) Alluvial Placer};
\end{scope}

% Pipe
\begin{scope}[xshift=4cm, yshift=-3.5cm]
\draw[thick] (0,0) -- (3,0) node[midway,above] {Surface};
\foreach \y in {0.5,1.0,1.5,2.0} {
  \draw[thin,gray] (0,\y) -- (3,\y);
}
\draw[poppurple,thick,fill=poppurple!20] (1.2,0) -- (1.2,2.2) -- (1.8,2.2) -- (1.8,0) -- cycle;
\draw[poppurple,thick] (1.2,2.2) -- (1.0,2.5) -- (2.0,2.5) -- (1.8,2.2);
\node[poppurple,align=center] at (1.5,2.7) {Kimberlite/\\(breccia) pipe};
\node[align=center] at (1.5,-0.5) {(e) Pipe};
\end{scope}
\end{tikzpicture}
\legende{Main morphological types of ore bodies in cross-section. (a) Stratiform layer parallel to bedding. (b) Discordant vein/lode cutting strata. (c) Stockwork of veinlets. (d) Placer concentrate in river channel. (e) Vertical pipe.}
\end{popfigure}

\begin{attention}[Morphology $\neq$ Genesis]
A \emph{vein} is a shape; a \emph{hydrothermal deposit} is an origin. A hydrothermal deposit can be a vein, a stockwork, a replacement body or a stratiform exhalative layer. In exam questions, if asked \emph{``Classify the deposit by morphology''}, answer \emph{vein, stratiform, stockwork, etc.} If asked \emph{``Classify by genesis''}, answer \emph{magmatic, hydrothermal, sedimentary, etc.} Mixing the two criteria loses marks.
\end{attention}

\courssub{Classification by Relation to Host Rock: Syngenetic vs Epigenetic}

This criterion asks: \emph{Did the ore form at the same time as the host rock, or later?} It is the single most powerful predictor of where to drill next.

\begin{definition}[Syngenetic Deposit]
Ore minerals crystallised \emph{simultaneously} with the host rock. The ore is an integral part of the stratigraphy. Examples: Banded Iron Formations (BIF), sedimentary manganese, VMS lenses, magmatic chromite layers in layered intrusions.
\end{definition}

\begin{definition}[Epigenetic Deposit]
Ore minerals were introduced \emph{after} the host rock had formed and usually after it had lithified. The ore cuts across, replaces or fills voids in the pre-existing rock. Examples: Hydrothermal veins, porphyry copper, skarns, Mississippi Valley Type (MVT), placer deposits (the host sediment is younger than the source rock).
\end{definition}

\begin{methode}[Recognising Syngenetic vs Epigenetic on a Cross-Section]
\begin{enumerate}
\item \textbf{Step 1:} Identify the \emph{primary layering/bedding} of the host rock (sedimentary beds, volcanic flow bands, igneous layering).
\item \textbf{Step 2:} Trace the \emph{ore-body boundaries}.
\item \textbf{Step 3:} \textbf{Concordant} (parallel to layering) $\rightarrow$ \emph{suggests syngenetic} (but check for replacement!). \textbf{Discordant} (cross-cutting) $\rightarrow$ \emph{epigenetic}.
\item \textbf{Step 4:} Look for \emph{alteration halos} (epigenetic) vs \emph{primary sedimentary structures} preserved in ore (syngenetic).
\item \textbf{Step 5:} Check \emph{age relationships}: cross-cutting veins, contact metamorphism, unconformities.
\end{enumerate}
\end{methode}

\begin{exemplebox}[Syngenetic vs Epigenetic in the Field]
\textbf{Case A:} In the \emph{Ntem Complex} (South Cameroon), a layer of massive chromitite follows the magmatic layering of a dunite-peridotite sequence. No cross-cutting veins, no alteration halo. $\rightarrow$ \textbf{Syngenetic magmatic}.

\textbf{Case B:} At \emph{Batouri} (East Cameroon), gold-bearing quartz veins cut across metasedimentary schists and are surrounded by sericite-pyrite alteration halos. $\rightarrow$ \textbf{Epigenetic hydrothermal}.
\end{exemplebox}

\courssub{Classification by Genetic Process (Overview)}

Genesis explains \emph{how} the ore formed. The GCE syllabus recognises five main families. Detailed mechanisms are covered in Sections 4 and 5; here we give the \emph{recognition criteria} for each family.

\begin{center}
\begin{tabularx}{\linewidth}{|l|X|X|}\hline
\rowcolor{popblueL}
\textbf{Genetic Family} & \textbf{Key Processes} & \textbf{Typical Examples (Cameroon / World)} \\ \hline
\textbf{Magmatic} & Crystallisation from silicate melt: early crystallisation (chromite, PGE), immiscible sulphide liquids (Ni-Cu-PGE), residual liquids (pegmatites). & \emph{Nkamouna} (Co-Ni-Mn laterite over ultramafic), \emph{Bushveld} (PGE), \emph{Sudbury} (Ni-Cu). \\ \hline
\textbf{Hydrothermal} & Precipitation from hot aqueous fluids: veins, stockworks, replacement, VMS, porphyry, epithermal, skarn. & \emph{Batouri} (Au), \emph{Mayo Darle} (W-Sn), \emph{Ok Tedi} (Cu-Au porphyry), \emph{Kidd Creek} (VMS). \\ \hline
\textbf{Sedimentary} & Chemical/biochemical precipitation from water: BIF, manganese, phosphorite, evaporites, placer (mechanical). & \emph{Mbalam/Nabeba} (BIF Fe), \emph{Bafoussam} (bauxite — residual), \emph{Witwatersrand} (Au-U paleoplacer). \\ \hline
\textbf{Residual / Weathering} & In-situ chemical weathering concentrates insoluble residues: lateritic Ni-Co, bauxite, supergene enrichment. & \emph{Nkamouna} (Co-Ni laterite), \emph{Minim-Martap} (bauxite), \emph{New Caledonia} (Ni laterite). \\ \hline
\textbf{Metamorphic} & Recrystallisation / remobilisation during metamorphism: graphite, kyanite, some Au deposits. & \emph{Mfou} (graphite), \emph{Lom} (kyanite). \\ \hline
\end{tabularx}
\end{center}

\begin{savaistu}[Cameroon's Genetic Diversity]
Cameroon hosts \textbf{all five genetic families} along the \cle{Cameroon Volcanic Line (CVL)} and the \cle{Congo Craton} margin: magmatic (chromite in ultramafics), hydrothermal (Au-W-Sn in Pan-African granites), sedimentary (BIF iron in Mbalam), residual (Co-Ni laterites at Nkamouna, bauxite on the Adamawa Plateau), metamorphic (graphite in Yaoundé Group). This makes Cameroon a natural laboratory for economic geology.
\end{savaistu}

\courssub{Classification by Commodity (Economic Grouping)}

This classification speaks the language of economists and metallurgists.

\begin{definition}[Metallic Deposits]
Ores from which metals are extracted. Sub-groups:
\begin{itemize}
  \item \textbf{Ferrous:} Fe, Mn, Cr (steel industry).
  \item \textbf{Base metals:} Cu, Pb, Zn, Ni, Co (industrial bulk).
  \item \textbf{Precious metals:} Au, Ag, PGE (high value, low tonnage).
  \item \textbf{Strategic / Critical:} W, Sn, Ta, Nb, Li, REE, Co (high-tech, supply risk).
\end{itemize}
\end{definition}

\begin{definition}[Non-Metallic (Industrial) Deposits]
Valued for physical/chemical properties, not metal content: limestone, phosphate, potash, salt, gypsum, barite, fluorite, graphite, kaolin, dimension stone, aggregates.
\end{definition}

\begin{definition}[Energy Minerals]
Uranium, thorium, coal, oil shale, geothermal — sometimes treated separately.
\end{definition}

\begin{attention}[Commodity $\neq$ Genesis]
A \emph{base-metal} deposit can be hydrothermal (VMS), sedimentary (SEDEX) or magmatic (Ni-Cu sulphide). Always state the genetic family \emph{and} the commodity group when describing a deposit fully.
\end{attention}

\courssub{Concordant vs Discordant Ore Bodies: Mining Consequences}

This geometric distinction, rooted in the syngenetic/epigenetic divide, has direct engineering impact.

\begin{propriete}[Concordant Ore Bodies]
Follow host-rock layering (stratiform, some stratabound, some replacement bodies).
\begin{itemize}
  \item \textbf{Mining:} Amenable to \emph{open-pit} (if near surface) or \emph{room-and-pillar / longwall} underground methods. High continuity $\rightarrow$ predictable reserves.
  \item \textbf{Dilution:} Low if contacts are sharp; high if gradational.
  \item \textbf{Example:} Mbalam BIF iron ore — open pit, low stripping ratio.
\end{itemize}
\end{propriete}

\begin{propriete}[Discordant Ore Bodies]
Cut across layering (veins, lodes, pipes, stockworks, some breccias).
\begin{itemize}
  \item \textbf{Mining:} Often \emph{selective underground} (cut-and-fill, shrinkage stoping, sublevel caving for stockworks). Irregular geometry $\rightarrow$ higher exploration density needed.
  \item \textbf{Dilution:} High — wall rock must be broken with ore.
  \item \textbf{Example:} Batouri gold veins — narrow, high-grade, underground.
\end{itemize}
\end{propriete}

\begin{aretenir}[Exam Checklist: Describe a Deposit Fully]
When a question says \emph{``Classify the following deposit''}, give \textbf{four labels}:
\begin{enumerate}
  \item Morphology (vein, stratiform, stockwork, placer, pipe\ldots)
  \item Host-rock relation (syngenetic / epigenetic)
  \item Genetic family (magmatic, hydrothermal, sedimentary, residual, metamorphic)
  \item Commodity group (ferrous, base, precious, strategic, industrial, energy)
\end{enumerate}
\end{aretenir}

\courssub{Summary Table: One Example per Morphological \& Genetic Type}

\begin{center}
\begin{tabularx}{\linewidth}{|l|l|l|X|}\hline
\rowcolor{popblueL}
\textbf{Morphology} & \textbf{Genetic Family} & \textbf{Deposit Example} & \textbf{Key Features} \\ \hline
Stratiform & Sedimentary (BIF) & \textbf{Mbalam / Nabeba} (Cameroon–Congo border) & Hematite-magnetite layers in Palaeoproterozoic metasediments; syngenetic; open-pit Fe. \\ \hline
Stratiform & Magmatic (layered intrusion) & \textbf{Bushveld Complex} (South Africa) & Chromitite / PGE reefs in pyroxenite; syngenetic; underground. \\ \hline
Vein / Lode & Hydrothermal (orogenic Au) & \textbf{Batouri} (East Cameroon) & Quartz-carbonate veins in shear zones; epigenetic; high-grade Au. \\ \hline
Stockwork & Hydrothermal (porphyry Cu) & \textbf{Ok Tedi} (Papua New Guinea) / \textbf{Mayo Darle} (Cameroon, W-Sn) & Quartz veinlets in granodiorite; epigenetic; bulk mining. \\ \hline
Pipe & Magmatic (kimberlite) & \textbf{Kimberley} (South Africa) / \textbf{Mobilong} (Cameroon, diamondiferous kimberlite?) & Vertical carrot-shaped; epigenetic; open pit then underground. \\ \hline
Pod / Lens & Hydrothermal (skarn) & \textbf{Mina Justa} (Peru) / \textbf{Bibemi} (Cameroon, Fe skarn?) & Irregular replacement at carbonate-intrusion contact. \\ \hline
Placer (alluvial) & Sedimentary (mechanical) & \textbf{Batouri / Kette} (Cameroon, Au) / \textbf{Witwatersrand} (paleoplacer) & Heavy minerals in river gravels; epigenetic (host sediment younger). \\ \hline
Residual blanket & Weathering (laterite) & \textbf{Nkamouna} (Cameroon, Co-Ni) / \emph{Minim-Martap} (bauxite) & In-situ enrichment over ultramafic/basalt; syngenetic w.r.t. weathering profile. \\ \hline
Disseminated & Metamorphic (graphite) & \textbf{Mfou} (Cameroon) / \emph{Bogala} (Sri Lanka) & Flaky graphite in gneiss; syngenetic (carbon-rich sediment) + metamorphic remobilisation. \\ \hline
\end{tabularx}
\end{center}

\courssub{Classification Tree: A Visual Synthesis}

The following diagram organises the four classification criteria into a single decision tree. Use it as a mental checklist during field mapping or exam revision.

\begin{popfigure}
\begin{tikzpicture}[
  scale=0.75,
  every node/.style={font=\small, align=center},
  level 1/.style={sibling distance=5.5cm},
  level 2/.style={sibling distance=3.2cm},
  level 3/.style={sibling distance=2.2cm},
  edge from parent/.style={draw, popline, thick},
  grow'=down
]
\node[popblue, rounded corners, fill=popblueL, draw=popblue, thick, text width=3.5cm]
  {\textbf{MINERAL DEPOSIT}\\\textit{Classify by 4 criteria}}
  child { node[poporange, rounded corners, fill=poporange!20, draw=poporange, text width=3cm]
    {\textbf{1. MORPHOLOGY}\\(Shape)}
    child { node[poporange!80!black, text width=2.5cm] {Stratiform / Stratabound} }
    child { node[poporange!80!black, text width=2.5cm] {Vein / Lode} }
    child { node[poporange!80!black, text width=2.5cm] {Stockwork / Disseminated} }
    child { node[poporange!80!black, text width=2.5cm] {Pipe / Pod / Lens} }
    child { node[poporange!80!black, text width=2.5cm] {Placer (alluvial, beach)} }
  }
  child { node[popgreen, rounded corners, fill=popgreenL, draw=popgreen, text width=3cm]
    {\textbf{2. HOST-ROCK RELATION}\\(Timing)}
    child { node[popgreen!70!black, text width=2.5cm] {Syngenetic\\(with host rock)} }
    child { node[popgreen!70!black, text width=2.5cm] {Epigenetic\\(post-dates host)} }
  }
  child { node[poppurple, rounded corners, fill=poppurple!20, draw=poppurple, text width=3.5cm]
    {\textbf{3. GENETIC FAMILY}\\(Origin)}
    child { node[poppurple!70!black, text width=2.8cm] {Magmatic} }
    child { node[poppurple!70!black, text width=2.8cm] {Hydrothermal} }
    child { node[poppurple!70!black, text width=2.8cm] {Sedimentary} }
    child { node[poppurple!70!black, text width=2.8cm] {Residual / Weathering} }
    child { node[poppurple!70!black, text width=2.8cm] {Metamorphic} }
  }
  child { node[poppink, rounded corners, fill=poppink!20, draw=poppink, text width=3cm]
    {\textbf{4. COMMODITY}\\(Economic value)}
    child { node[poppink!70!black, text width=2.5cm] {Metallic: Ferrous, Base, Precious, Strategic} }
    child { node[poppink!70!black, text width=2.5cm] {Non-metallic (Industrial)} }
    child { node[poppink!70!black, text width=2.5cm] {Energy} }
  };
\end{tikzpicture}
\legende{Classification decision tree. Every deposit receives one label from each of the four branches.}
\end{popfigure}

\courssub{Worked Example: Full Classification of a Deposit}

\begin{exoresolu}[Classifying the Nkamouna Cobalt-Nickel Deposit (East Cameroon)]
\textbf{Statement:} The Nkamouna deposit consists of a 10--30 m thick blanket of cobalt-nickel-manganese oxides (asbolane, lithiophorite) overlying weathered ultramafic rocks (serpentinised peridotite) of the \emph{Ntem Complex}. The ore zone follows the paleo-topography and shows a clear vertical zonation from Fe-rich duricrust at surface to Co-Ni-Mn enrichment at depth. Classify this deposit using all four criteria.

\textbf{Detailed solution:}

\begin{enumerate}
  \item \textbf{Morphology:} \cle{Residual blanket / lateritic profile}. The ore forms a horizontal to gently undulating layer parallel to the paleo-surface, not to bedrock structure. It is a \emph{weathering mantle}, not a primary igneous layer.
  \item \textbf{Host-rock relation:} \cle{Epigenetic} with respect to the \emph{bedrock} (the ultramafic host is Archaean, the ore formed during Cenozoic weathering). However, it is \cle{syngenetic} with respect to the \emph{lateritic profile} itself (the Co-Ni enrichment formed concurrently with the laterite horizons). \emph{Always specify the reference frame!}
  \item \textbf{Genetic family:} \cle{Residual / Weathering (supergene lateritic)}. Intense tropical weathering leached Mg, Si, Fe (partly) and concentrated immobile Co, Ni, Mn as oxide/hydroxide minerals in the saprolite and mottled zone.
  \item \textbf{Commodity group:} \cle{Strategic / Critical metals} (Co for batteries, Ni for stainless steel, Mn for alloys). Also base metals by traditional classification.
\end{enumerate}

\textbf{Final classification statement:} \emph{``Nkamouna is a residual blanket (morphology), epigenetic-to-bedrock / syngenetic-to-profile (host relation), lateritic supergene (genesis), strategic Co-Ni-Mn (commodity) deposit.''}
\end{exoresolu}

\courssub{Exercises}

\begin{exercice}[Classification Practice]
For each of the following descriptions, give the \textbf{four classification labels} (Morphology, Host-relation, Genesis, Commodity).

\begin{enumerate}
  \item A 2 m thick layer of magnetite-quartzite intercalated with amphibolites in the \emph{Yaoundé Group} (Pan-African belt). The layer is folded with the host rocks.
  \item A swarm of 5--20 cm wide quartz-tourmaline-cassiterite veins cutting a two-mica granite near \emph{Mayo Darle} (Adamawa). Veins are surrounded by greisen alteration.
  \item Rounded grains of gold and diamond concentrated in the gravels of the \emph{Sangha River} palaeo-channels (East Cameroon).
  \item A lenticular body of massive chromitite (5 m $\times$ 200 m) within a dunite layer of the \emph{Ntem Complex} ultramafic suite. Contacts are sharp and parallel to magmatic layering.
  \item A 500 m $\times$ 200 m $\times$ 50 m zone of disseminated chalcopyrite and molybdenite in a granodiorite stock, with a surrounding halo of propylitic alteration (chlorite-epidote).
\end{enumerate}
\end{exercice}

\begin{corrige}[Exercise 1]
\begin{enumerate}
  \item \textbf{Morphology:} Stratiform. \textbf{Host-relation:} Syngenetic (sedimentary/exhalative). \textbf{Genesis:} Sedimentary (BIF-type) + metamorphic overprint. \textbf{Commodity:} Ferrous (Fe).
  \item \textbf{Morphology:} Vein / Lode (swarm). \textbf{Host-relation:} Epigenetic. \textbf{Genesis:} Hydrothermal (granite-related, Sn-W). \textbf{Commodity:} Strategic (Sn, W).
  \item \textbf{Morphology:} Placer (alluvial / palaeo-channel). \textbf{Host-relation:} Epigenetic (host sediment younger than source). \textbf{Genesis:} Sedimentary (mechanical concentration). \textbf{Commodity:} Precious (Au) + Gem (diamond).
  \item \textbf{Morphology:} Lens / Pod (within layer). \textbf{Host-relation:} Syngenetic (magmatic). \textbf{Genesis:} Magmatic (early crystallisation / immiscible liquid). \textbf{Commodity:} Strategic (Cr, PGE).
  \item \textbf{Morphology:} Disseminated / Stockwork. \textbf{Host-relation:} Epigenetic. \textbf{Genesis:} Hydrothermal (porphyry Cu-Mo). \textbf{Commodity:} Base metals (Cu, Mo).
\end{enumerate}
\end{corrige}

\begin{attention}[Common Exam Traps]
\begin{itemize}
  \item \textbf{Trap 1:} Calling a BIF ``metamorphic'' because it is in gneiss. \emph{Correction:} The \emph{protolith} is sedimentary; metamorphism recrystallised it. Classify by \emph{primary genesis}: Sedimentary (BIF).
  \item \textbf{Trap 2:} Calling a placer ``syngenetic'' because the gold grains are in the sediment. \emph{Correction:} The \emph{sediment} is the host; the gold was derived from an older source. $\rightarrow$ Epigenetic.
  \item \textbf{Trap 3:} Confusing \emph{stockwork} (many small veins) with \emph{dissemination} (individual grains). In a stockwork you see veinlets; in a dissemination you need a hand lens.
  \item \textbf{Trap 4:} Forgetting to state the \emph{reference frame} for syngenetic/epigenetic (bedrock vs weathering profile vs sediment).
\end{itemize}
\end{attention}

\begin{savaistu}[From Classification to Exploration: The Cameroon Example]
The \emph{Minim-Martap} bauxite deposit (Adamawa Plateau) was discovered by \textbf{morphological reasoning}: lateritic plateaux (residual blanket morphology) overlying Cretaceous basalts (host rock) in a tropical climate $\rightarrow$ target for \emph{residual/weathering} genesis $\rightarrow$ bauxite (non-metallic/industrial commodity). The same logic guides exploration for \emph{Nkamouna-style} Co-Ni laterites: look for ultramafic bedrock + deep weathering profile + plateau topography. Classification is not academic — it is the explorationist's compass.
\end{savaistu}

\courssub{Key Takeaways}

\begin{aretenir}[What You Must Remember]
\begin{itemize}
  \item Classification uses \textbf{four independent criteria}: Morphology, Host-rock relation (Syngenetic/Epigenetic), Genesis, Commodity. Always give all four.
  \item \textbf{Morphology} controls mining method: stratiform $\rightarrow$ open pit / longwall; veins/pipes $\rightarrow$ selective underground; stockworks $\rightarrow$ bulk underground.
  \item \textbf{Syngenetic} = ore formed \emph{with} host rock (concordant, stratigraphic control). \textbf{Epigenetic} = ore introduced \emph{later} (discordant, structural control).
  \item \textbf{Five genetic families}: Magmatic, Hydrothermal, Sedimentary, Residual/Weathering, Metamorphic. Cameroon has examples of all five.
  \item \textbf{Commodity groups}: Ferrous, Base, Precious, Strategic (metallic); Industrial (non-metallic); Energy.
  \item \textbf{Concordant} ore bodies are stratigraphically controlled; \textbf{discordant} bodies are structurally controlled.
  \item \textbf{Exam tip:} When asked to ``classify'', explicitly list the four labels. Never mix morphology terms (vein) with genesis terms (hydrothermal) as if they were the same level of classification.
\end{itemize}
\end{aretenir}

The next two sections will dive deep into the \textbf{formation processes} of each genetic family — first magmatic and hydrothermal (Section 4), then sedimentary, weathering and metamorphic (Section 5). Master this classification framework now; it is the scaffold on which all those genetic models will hang.

\coursec{Formation of Ore Deposits I: Magmatic and Hydrothermal Processes}

\courssub{From Classification to Genesis}
In the previous section we learned how mineral deposits are \emph{classified} according to their morphology, host rock and genetic affiliation.  Classification tells us \emph{what} a deposit looks like; genesis explains \emph{how} it formed.  Understanding the processes that concentrate metals from trace levels in the Earth's crust into mineable ore bodies is the core of economic geology.  Two great families of processes dominate the formation of the world's major metal deposits: \textbf{magmatic processes} that operate while a melt is still largely liquid, and \textbf{hydrothermal processes} that involve hot, aqueous fluids moving through the crust.  Both families are intimately linked — magmas supply the heat, the metals and often the fluids themselves.  In this section we follow the journey of metals from a cooling magma chamber to the veins and stockworks that make up many of Cameroon's gold and base‑metal prospects.

\courssub{Magmatic Processes: Segregation, Immiscibility and the Pegmatitic Stage}

\textbf{Magmatic segregation} is the physical separation of crystals from a melt by gravity.  As a mafic magma cools, dense oxide and sulphide minerals crystallise early and, because they are heavier than the surrounding silicate liquid, they sink toward the floor of the intrusion.  This produces \emph{layered intrusions} with rhythmic bands of chromite, magnetite, ilmenite and platinum‑group minerals — classic examples are the Bushveld Complex (South Africa) and the \emph{Ntem Complex} in southern Cameroon.

\begin{definition}[Magmatic Segregation]
The gravitational settling of early‑formed, dense crystals (oxides, sulphides, platinum‑group minerals) from a convecting magma, producing layered concentrations at the base of an intrusion.
\end{definition}

\begin{popfigure}
\begin{tikzpicture}[scale=0.9, every node/.style={font=\small, text width=2.2cm, align=center}]
% Country rock
\fill[popmuted!30] (-3,3) rectangle (3,4);
\node at (0,3.5) {Country rock};
% Magma chamber
\fill[poporange!20] (-3,-2) rectangle (3,3);
% Settled layers at base
\foreach \y/\col/\lab in {-1.8/popink/Chromite, -1.4/poppurple/Magnetite, -1.0/popblue/Ilmenite, -0.6/popgold/Platinum} {
  \fill[\col!60] (-2.5,\y) rectangle (2.5,\y+0.3);
  \node at (0,\y+0.15) {\lab\ layer};
}
% Sulphide liquid pool
\fill[popdark!80] (-2,-2) rectangle (2,-1.8);
\node at (0,-1.9) {Immiscible sulphide liquid (Ni‑Cu)};
% Arrows for settling
\draw[->, thick, popdark] (0,2.5) -- (0,1.5) node[midway, right] {Crystal settling};
\draw[->, thick, popdark] (0,-1.5) -- (0,-1.85) node[midway, right] {Sulphide droplet sinking};
% Labels
\node at (-3.5,1) {Cooling intrusion};
\node at (3.5,1) {Convection currents};
\end{tikzpicture}
\legende{Cross‑section of a cooling mafic intrusion showing crystal settling (chromite, magnetite, ilmenite, platinum layers) and an immiscible sulphide liquid accumulating at the base.}
\end{popfigure}

A second magmatic mechanism is \textbf{magmatic immiscibility}.  When a mafic magma becomes saturated in sulphur, a separate sulphide liquid phase splits off from the silicate melt — just as oil separates from water.  This sulphide liquid is an extremely efficient scavenger of chalcophile metals (Ni, Cu, Co, PGE).  Because it is denser than the silicate melt, it also sinks and pools at the base of the intrusion, forming massive sulphide ore bodies (e.g. the Sudbury nickel‑copper deposit in Canada).  In Cameroon, the \emph{Ntem Complex} shows evidence of such immiscible sulphide segregation in its gabbroic units.

\begin{definition}[Immiscible Sulphide Liquid]
A dense, sulphur‑rich liquid phase that separates from a silicate magma upon sulphur saturation, strongly partitioning chalcophile metals (Ni, Cu, Co, PGE) and sinking to form magmatic sulphide deposits.
\end{definition}

The final magmatic stage is the \textbf{pegmatitic stage}.  As crystallisation proceeds, the residual melt becomes enriched in water, fluxing components (B, F, P) and incompatible elements (Li, Be, Sn, Ta, Nb, rare earths).  This highly fluid, low‑viscosity melt intrudes fractures and crystallises as coarse‑grained \emph{pegmatites}.  Pegmatites are the primary source of tin (cassiterite), tantalum (columbite‑tantalite), lithium (spodumene, lepidolite) and gemstones (beryl, tourmaline).  The \emph{Central African tin belt}, which extends through the Dja River region of East Cameroon, owes its mineralisation to late‑stage pegmatitic fluids derived from the Pan‑African granites.

\begin{definition}[Pegmatite]
An exceptionally coarse‑grained igneous rock formed from the final, volatile‑rich residual melt of a granitic magma, typically enriched in rare elements (Li, Be, Sn, Ta, Nb, gemstones).
\end{definition}

\begin{savaistu}[Cameroon's Tin and Rare‑Metal Potential]
The \emph{Mbalam} and \emph{Bipindi} areas in the East Region host alluvial cassiterite derived from weathered pegmatites.  Exploration targeting the parent granites and their pegmatitic apophyses could reveal primary Sn‑Ta‑Nb deposits similar to those of the Manono‑Kitotolo belt in the DRC.
\end{savaistu}

\courssub{Hydrothermal Processes: Fluids, Transport and Deposition}

When magmas crystallise, they release large volumes of hot, metal‑bearing aqueous fluids — \textbf{hydrothermal solutions}.  These fluids may also originate from \emph{meteoric water} (rainwater that circulates deep into the crust) or \emph{connate water} (trapped sedimentary brine).  Regardless of origin, hydrothermal fluids are the principal agents that transport metals over kilometres and deposit them in fractures, pore spaces and reactive host rocks.

\begin{definition}[Hydrothermal Solution]
A hot (typically \SI{100}{\ensuremath{{}^{\circ}\mathrm{C}}}–\SI{600}{\ensuremath{{}^{\circ}\mathrm{C}}}), pressurised aqueous fluid, often saline and acidic, capable of dissolving, transporting and precipitating metallic ions in the Earth's crust.
\end{definition}

\begin{propriete}[Temperature‑Pressure Control on Metal Solubility]
The solubility of most ore metals (Cu, Pb, Zn, Au, Ag) in hydrothermal fluids \emph{decreases} as temperature and pressure fall.  Consequently, cooling, decompression, boiling, fluid mixing or reaction with wall rocks triggers metal precipitation.
\end{propriete}

Hydrothermal deposition occurs in three main geometric styles:
\begin{itemize}
\item \textbf{Vein filling} — fluids precipitate minerals in open fractures (e.g. quartz‑gold veins).
\item \textbf{Replacement} — fluids dissolve the host rock and simultaneously precipitate ore minerals, preserving the original rock texture (e.g. Mississippi Valley‑type Pb‑Zn in limestone).
\item \textbf{Dissemination} — fine‑grained ore minerals precipitate throughout a permeable rock volume (e.g. porphyry copper stockworks).
\end{itemize}

\begin{definition}[Replacement (Metasomatism)]
A hydrothermal process in which a mineral assemblage is chemically replaced by a new assemblage, the fluid simultaneously removing the original components and supplying the new ones, often preserving the original rock fabric.
\end{definition}

\courssub{Metal Zoning Around Intrusions: Porphyry and Epithermal Systems}

A classic illustration of hydrothermal zoning is the \textbf{porphyry copper‑gold system}.  A small, porphyritic granodiorite to diorite intrusion (the \emph{porphyry stock}) emplaces at shallow crustal levels (\SI{1}{–}\SI{3}{\kilo\meter}).  It releases a magmatic‑hydrothermal fluid that creates a dense network of quartz‑veinlets — a \emph{stockwork} — in the intrusion and the surrounding country rock.  Metals precipitate in concentric shells reflecting decreasing temperature and changing fluid chemistry:

\begin{popfigure}
\begin{tikzpicture}[scale=0.85, every node/.style={font=\small, text width=2.5cm, align=center}]
% Central intrusion
\fill[poporange!30] (-1.5,-1.5) rectangle (1.5,1.5);
\node at (0,0) {Porphyry stock (granodiorite)};
% Stockwork veins
\foreach \ang in {0,30,60,90,120,150,180,210,240,270,300,330} {
  \draw[popdark, thin] (0,0) -- (\ang:2.2);
}
% Zones
\draw[popblue, thick] (0,0) circle (2.2);
\node[popblue] at (0,2.5) {Cu‑Mo core (> \SI{350}{\ensuremath{{}^{\circ}\mathrm{C}}})};
\draw[popgreen, thick] (0,0) circle (3.0);
\node[popgreen] at (0,3.3) {Zn‑Pb shell (\SI{250}{\ensuremath{{}^{\circ}\mathrm{C}}}–\SI{350}{\ensuremath{{}^{\circ}\mathrm{C}}})};
\draw[poppurple, thick] (0,0) circle (3.8);
\node[poppurple] at (0,4.1) {Ag‑Au halo (< \SI{250}{\ensuremath{{}^{\circ}\mathrm{C}}})};
% Epithermal veins above
\draw[poppink, thick, decorate, decoration={snake, amplitude=2pt, segment length=6pt}] (-3,4.5) -- (3,4.5);
\node[poppink] at (0,5) {Epithermal Au‑Ag veins (low‑sulphidation)};
% Country rock
\fill[popmuted!20] (-5,-2) rectangle (5,5);
\node at (-4.5,4.5) {Country rock};
\end{tikzpicture}
\legende{Schematic cross‑section of a porphyry Cu‑Au system showing the central stock, stockwork veinlets, and concentric metal zoning (Cu‑Mo core → Zn‑Pb shell → Ag‑Au halo).  Epithermal veins form above the porphyry environment.}
\end{popfigure}

\begin{exemplebox}[Porphyry‑Epithermal Link in East Cameroon]
The \emph{Batouri} and \emph{Bertoua} greenstone belts host numerous artisanal gold workings.  Geological mapping reveals that many quartz‑gold veins are spatially associated with small, porphyritic granodiorite intrusions — the roots of ancient porphyry systems.  The high‑temperature Cu‑Mo core is often eroded away, leaving the lower‑temperature Au‑Ag epithermal veins exposed at surface.  Recognising this genetic link guides exploration: soil geochemistry for Cu‑Mo pathfinders can vector toward buried porphyry centres beneath the gold veins.
\end{exemplebox}

\courssub{The Critical Role of Structures}

Faults, fractures, shear zones and lithological contacts act as \emph{high‑permeability pathways} that channel hydrothermal fluids from their deep source to the depositional site.  They also serve as \emph{traps}: a change in rock competency, a fault bend or a lithological contact can cause a pressure drop, fluid boiling or chemical reaction — all triggering precipitation.  In the \emph{Adamawa Plateau}, the \emph{Sanaga Fault} and its splays control the localisation of both magmatic (Ni‑Cu) and hydrothermal (Au) mineralisation.

\begin{methode}[Reconstructing Mineralisation Sequence from Cross‑Cutting Veins]
\begin{enumerate}
\item Map all vein sets, noting orientation, mineralogy and alteration halos.
\item Identify cross‑cutting relationships: a vein that cuts another is younger.
\item Record the mineral assemblage of each vein generation (e.g. early quartz‑pyrite‑chalcopyrite → late quartz‑gold‑sericite).
\item Correlate vein generations with known regional deformation events (e.g. D1, D2).
\item Build a paragenetic sequence table linking structural event, fluid temperature (from fluid inclusions) and metal assemblage.
\end{enumerate}
\end{methode}

\begin{attention}[Epigenetic Nature of Hydrothermal Deposits]
Even when hydrothermal veins occur \emph{inside} the intrusion that supplied the fluid, they are \textbf{epigenetic} — they formed after the host rock solidified.  Do not classify them as magmatic (syngenetic) simply because of their spatial association.
\end{attention}

\begin{attention}[Gold: Primary vs. Placer]
Most gold in the crust is deposited by hydrothermal fluids (primary gold).  Placer gold is \emph{secondary}, derived from the erosion of primary hydrothermal veins.  In Cameroon, the alluvial gold of the \emph{Sangha} and \emph{Boumba} rivers originates from the primary quartz‑gold veins of the East Region greenstone belts.
\end{attention}

\courssub{Worked Example: Interpreting a Porphyry Cross‑Section}

\begin{exoresolu}[Porphyry Zoning Interpretation]
\textbf{Statement:} A geological cross‑section through a porphyry system shows (from centre outward): a stockwork of quartz‑chalcopyrite‑molybdenite veinlets in a granodiorite; a surrounding zone of quartz‑sphalerite‑galena veins in the contact aureole; and an outer halo of quartz‑arsenopyrite‑gold veins in the country rock.  Fluid inclusion data give homogenisation temperatures of \SI{420}{\ensuremath{{}^{\circ}\mathrm{C}}} (core), \SI{300}{\ensuremath{{}^{\circ}\mathrm{C}}} (intermediate) and \SI{200}{\ensuremath{{}^{\circ}\mathrm{C}}} (outer).  Explain the metal zoning and suggest an exploration vector for the buried Cu‑Mo core.

\textbf{Detailed solution:}
\begin{enumerate}
\item \textbf{Identify the three zones:} The mineral assemblages and temperatures match the classic porphyry zoning: Cu‑Mo (chalcopyrite + molybdenite) at highest T, Zn‑Pb (sphalerite + galena) at intermediate T, Au‑Ag‑As (arsenopyrite + gold) at lowest T.
\item \textbf{Explain the zoning:} As the magmatic‑hydrothermal fluid migrates outward, it cools and depressurises.  Metal solubility drops sequentially: Cu and Mo precipitate first (high T, high fS$_2$), then Zn and Pb, finally Au and Ag (low T, lower fS$_2$, possible boiling).  Wall‑rock reaction (e.g. with carbonate) can also shift pH and cause deposition.
\item \textbf{Exploration vector:} The Cu‑Mo core is centred on the porphyry stock.  Soil/stream sediment surveys for Cu, Mo and pathfinder elements (Re, Se, Te) should be focused \emph{up‑gradient} (toward the centre of the stockwork) from the known Zn‑Pb and Au anomalies.  Geophysical IP/resistivity surveys can delineate the stockwork chargeability high that marks the core.
\item \textbf{Key takeaway:} Metal zoning is a direct consequence of the temperature‑solubility relationship (Property above).  Mapping the zoning allows the geologist to \emph{vector} toward the highest‑value part of the system.
\end{enumerate}
\end{exoresolu}

\courssub{Key Points to Retain}

\begin{aretenir}[Formation of Ore Deposits I — Essentials]
\begin{itemize}
\item Magmatic segregation (crystal settling) forms layered oxide/sulphide deposits in mafic intrusions (chromite, Pt).
\item Magmatic immiscibility produces dense sulphide liquids that scavenge Ni, Cu, PGE and settle as massive sulphides.
\item Pegmatites concentrate Li, Be, Sn, Ta, Nb and gemstones from the final volatile‑rich melt.
\item Hydrothermal fluids (magmatic, meteoric, connate) transport metals in solution; solubility falls with T and P, causing precipitation.
\item Deposition styles: vein filling, replacement, dissemination.
\item Porphyry Cu‑Au systems show concentric metal zoning (Cu‑Mo core → Zn‑Pb shell → Ag‑Au halo) linked to cooling fluid.
\item Epithermal veins (low‑sulphidation Au‑Ag) form above porphyry environments; they are the exposed expression of deeper systems.
\item Structures (faults, fractures, contacts) are fluid pathways and traps — essential for localising ore.
\item Hydrothermal deposits are epigenetic; primary gold is hydrothermal, placers are secondary.
\end{itemize}
\end{aretenir}

\coursec{Formation of Ore Deposits II: Sedimentary, Weathering and Metamorphic Processes}

In the previous section we followed ore-forming fluids \emph{upwards} from magmas and hot hydrothermal systems deep inside the crust. But not all ores are born in fire. Walk along the Wouri estuary, or across the red plateaux of the Adamaoua, and you are standing on ore deposits that formed \emph{at or near the Earth's surface}, at ordinary temperatures, through the patient work of water, wind, organisms and time. Rivers sort and concentrate heavy minerals; tropical rain leaches away silica and leaves aluminium behind; seawater evaporates and abandons its salts; and when rocks are later buried and heated, metamorphism recrystallises old materials into new economic products. This section studies these \cle{exogenous} (surface) and \cle{metamorphic} processes, which together account for most of the world's iron, aluminium, salt, phosphate, placer gold and diamonds.

\courssub{Sedimentary (Syngenetic) Deposits}

\textbf{Observation.} When a shallow coastal lagoon or an enclosed salt lake evaporates under a hot sun, white crusts of salt appear at its margins. When certain ancient seas stagnated, iron precipitated layer upon layer on their floors. In both cases, the ore was deposited \emph{at the same time as the enclosing sediment}: such deposits are called \cle{syngenetic}.

\begin{definition}[Syngenetic sedimentary deposit]
A \cle{syngenetic sedimentary deposit} is an ore body formed by the direct precipitation or accumulation of useful material from surface waters (sea, lake, lagoon) at the same time as the sediments that enclose it. It typically forms extensive, regular \emph{beds} or \emph{seams} interstratified with normal sedimentary rocks.
\end{definition}

\textbf{Chemical precipitates.} Three great families are distinguished.

\begin{enumerate}
\item \textbf{Banded iron formations (BIF).} In the early history of the Earth (roughly before 1.8 billion years ago), the oceans contained abundant dissolved iron and the atmosphere little oxygen. As photosynthetic organisms released oxygen, the iron oxidised and precipitated, producing alternations of iron-rich (haematite, magnetite) and silica-rich (chert) bands. These \cle{banded iron formations} supply the greater part of the world's iron ore (e.g.\ the enormous deposits of Australia, Brazil, and the Kribi--Mbalam iron province of southern Cameroon).
\item \textbf{Evaporites.} When an arm of the sea or a salt lake is cut off and evaporation exceeds inflow, dissolved salts crystallise in a predictable order of increasing solubility: first carbonates, then \cle{gypsum} (\ce{CaSO4.2H2O}), then rock salt (\cle{halite}, \ce{NaCl}), and finally the highly soluble \cle{potash} salts (sylvite, \ce{KCl}). Evaporite beds are mined for salt, fertiliser and chemical industries.
\item \textbf{Biochemical accumulations: phosphates.} In warm, shallow seas with abundant life, phosphorus from organic remains (bones, shells, guano) accumulates on the sea floor and is converted to \cle{phosphate} rock (apatite), the essential raw material of fertilisers.
\end{enumerate}

\begin{attention}[Syngenetic versus epigenetic]
Do not confuse \emph{syngenetic} (ore formed \textbf{with} the host sediment, as a bed) with \emph{epigenetic} (ore introduced \textbf{later} into pre-existing rocks, as veins). A salt bed is syngenetic; a gold-bearing quartz vein cutting across schists is epigenetic. GCE examiners regularly test this distinction.
\end{attention}

\courssub{Placer Deposits: Concentration by Moving Water}

\textbf{Observation.} A gold panner on the River Lom or the Vina swirls gravel in a pan: the light sand washes away while dense, shiny grains settle at the bottom. The river has already done most of this sorting naturally. Minerals that are both \emph{dense} and \emph{chemically resistant} survive weathering and transport, and accumulate where the current slackens.

\begin{definition}[Placer deposit]
A \cle{placer} is a deposit formed by the mechanical concentration of dense, durable minerals (gold, diamond, cassiterite \ce{SnO2}, ilmenite \ce{FeTiO3}, platinum, monazite) by moving water, after the weathering and erosion of a primary source rock. Placers are classified by their setting:
\begin{itemize}
\item \cle{eluvial} placers: concentrated by gravity and creep on slopes, just downhill from the source;
\item \cle{alluvial} placers: concentrated in river gravels, at the base of the gravel bed, in potholes and behind obstacles;
\item \cle{beach} placers: concentrated by wave action along shorelines (e.g.\ ilmenite and rutile sands).
\end{itemize}
\end{definition}

The concentrating agent is \emph{hydraulic sorting}: a current that can carry quartz (density 2.65) may be unable to move gold (density 19.3) or cassiterite (density 6.9). The heavy grains therefore sink through the moving gravel and collect at its base, on or near the bedrock, often in cracks, potholes and the lee of boulders or bars.

\begin{popfigure}
\begin{tikzpicture}[scale=1.0, font=\small]
% bedrock
\fill[popmuted!40] (0,0) -- (0,0.9) .. controls (2,0.4) and (3,1.1) .. (5,0.6) .. controls (6.5,0.2) and (7,1.0) .. (9,0.5) -- (9,0) -- cycle;
\node at (4.5,0.15) {\textbf{Bedrock}};
% gravel layer
\fill[popgold!25] (0,0.9) .. controls (2,0.4) and (3,1.1) .. (5,0.6) .. controls (6.5,0.2) and (7,1.0) .. (9,0.5) -- (9,1.5) .. controls (6,1.9) and (3,1.6) .. (0,1.8) -- cycle;
\node[popdark] at (1.6,1.45) {Gravels (quartz, pebbles)};
% sand layer
\fill[popblueL] (0,1.8) .. controls (3,1.6) and (6,1.9) .. (9,1.5) -- (9,2.5) .. controls (5,2.9) and (3,2.5) .. (0,2.7) -- cycle;
\node[popblue] at (7.6,2.2) {Sand};
% water
\fill[popblue,opacity=0.25] (0,2.7) .. controls (3,2.5) and (5,2.9) .. (9,2.5) -- (9,3.4) -- (0,3.4) -- cycle;
\node[popblue] at (1.2,3.1) {River water};
\draw[->, popblue, thick] (3.2,3.05) -- (4.6,3.05);
\node[popblue] at (3.9,3.3) {current};
% gold grains at base
\foreach \x/\y in {1.1/0.95, 2.2/0.62, 2.6/0.7, 3.1/1.0, 5.2/0.62, 5.6/0.6, 6.7/0.35, 7.0/0.42, 7.3/0.55}{
  \fill[popgold] (\x,\y) circle (0.07);}
\node[popdark, align=center, text width=3.4cm] at (7.2,1.35) {Gold grains concentrated at base of gravels};
\draw[->, popdark] (6.9,1.15) -- (6.7,0.55);
% obstacle
\fill[popmuted!70] (4.6,0.6) -- (4.9,1.35) -- (5.2,0.6) -- cycle;
\node[popdark, align=center, text width=2.2cm] at (4.9,2.0) {obstacle (boulder)};
\draw[->, popdark] (4.9,1.85) -- (4.9,1.45);
\end{tikzpicture}
\legende{Placer concentration in a river: dense gold grains sink through the gravels and accumulate at the base, in bedrock hollows and behind obstacles.}
\end{popfigure}

\begin{methode}[Recognising a placer versus a primary vein]
\begin{enumerate}
\item \textbf{Grain shape:} placer grains are \emph{rounded and worn} by transport; vein minerals are \emph{angular, interlocking crystals}.
\item \textbf{Association:} placers contain a mixture of \emph{dense, resistant} minerals (gold + cassiterite + ilmenite + magnetite); veins contain quartz, sulphides and gangue minerals.
\item \textbf{Geometry:} placers follow \emph{sedimentary layers} (gravel beds, terraces, beaches); veins \emph{cut across} the rock structure.
\item \textbf{Context:} a placer lies in an unconsolidated sedimentary deposit near a drainage system; a vein lies within hard, often altered bedrock.
\end{enumerate}
\end{methode}

\begin{savaistu}[Placers of Cameroon]
Alluvial gold is washed by local miners in the rivers of the \textbf{East Region} (B\'etar\'e-Oya, Batouri) and the \textbf{Adamaoua}, and diamonds have long been recovered from placers in the Mobilong area near Yokadouma, close to the Central African border. These placers point to primary sources in the surrounding Precambrian basement rocks.
\end{savaistu}

\courssub{Residual (Weathering) Deposits: Laterites and Bauxite}

\textbf{Observation.} Laterite road cuttings around Yaound\'e show a striking vertical zonation: dark fresh rock at depth passes upward into soft, blotched, clay-rich material, then into a red, iron-rich horizon, sometimes capped by a hard crust. This is a \cle{weathering profile}, and it is an ore factory.

Under a \cle{humid tropical climate} --- high temperature, abundant rainfall, dense vegetation --- rainwater charged with \ce{CO2} and organic acids attacks the bedrock. The most soluble elements (\ce{Ca}, \ce{Na}, \ce{K}, and even \ce{SiO2}) are dissolved and carried away, while the least mobile elements (\ce{Al}, \ce{Fe}, \ce{Ti}) remain and are relatively enriched. This is \emph{residual concentration}: nothing is added, but the removal of the soluble fraction concentrates what is left.

\begin{definition}[Laterite and bauxite]
A \cle{laterite} is a red, residual soil and weathering mantle rich in iron and aluminium oxides and hydroxides, formed by intense tropical weathering. \cle{Bauxite} is the aluminium-rich end-product of this process: a \emph{rock} composed essentially of aluminium hydroxide minerals (gibbsite \ce{Al(OH)3}, boehmite, diaspore) with iron oxides, clay and residual quartz. A \cle{residual deposit} is any ore formed by such in-situ enrichment during weathering.
\end{definition}

\begin{propriete}[Condition for residual concentration]
Residual deposits form where soluble elements (\ce{Si}, \ce{Ca}, \ce{Na}, \ce{K}) are leached away faster than immobile elements (\ce{Al}, \ce{Fe}) are removed. They therefore require: (i) a warm, humid climate with strong chemical weathering; (ii) good drainage; (iii) a stable land surface (low erosion) preserved for a long time; and (iv) a suitable parent rock (aluminosilicate rocks for bauxite, ultrabasic rocks for nickel laterite).
\end{propriete}

\begin{popfigure}
\begin{tikzpicture}[scale=1.0, font=\small]
% layers (bottom to top)
\fill[popmuted!50] (0,0) rectangle (9,1.4);
\fill[poppink!25] (0,1.4) rectangle (9,3.0);
\fill[poporange!30] (0,3.0) rectangle (9,4.2);
\fill[popink!60] (0,4.2) rectangle (9,4.9);
% wavy boundaries
\draw[popdark] (0,1.4) .. controls (3,1.2) and (6,1.6) .. (9,1.4);
\draw[popdark] (0,3.0) .. controls (3,2.8) and (6,3.2) .. (9,3.0);
\draw[popdark] (0,4.2) .. controls (3,4.0) and (6,4.4) .. (9,4.2);
% labels
\node[align=center] at (4.5,0.7) {\textbf{Fresh rock} (unaltered parent rock)};
\node[align=center] at (4.5,2.2) {\textbf{Saprolite}: soft, clay-rich,\\ soluble elements leached};
\node[align=center] at (4.5,3.6) {\textbf{Ferruginous / bauxitic horizon}:\\ residual enrichment in Al and Fe};
\node[align=center, text=white] at (4.5,4.55) {\textbf{Cuirass (hard ferruginous cap)}};
% arrows showing leaching and enrichment
\draw[->, popblue, very thick] (10.0,4.6) -- (10.0,1.0);
\node[popblue, align=center, text width=3.2cm] at (10.0,0.4) {rainwater leaches Si, Ca, Na, K downward and out};
\draw[->, poporange, very thick] (11.6,1.2) -- (11.6,4.4);
\node[poporange, align=center, text width=2.6cm] at (11.6,5.15) {Al, Fe relatively enriched upward};
\end{tikzpicture}
\legende{Ideal lateritic weathering profile. Downward-percolating water removes soluble elements; immobile Al and Fe accumulate residually in the upper horizons.}
\end{popfigure}

\textbf{Main residual deposits.}
\begin{itemize}
\item \textbf{Bauxite (aluminium ore):} forms by lateritic weathering of aluminosilicate rocks (granites, gneisses, syenites) under humid tropical climates. The great Cameroonian example is the \textbf{Minim--Martap / Ngaoundal bauxite} on the Adamaoua Plateau: a thick lateritic mantle developed over syenitic and granitic basement, preserved on a high, stable, well-drained peneplain.
\item \textbf{Lateritic nickel:} weathering of ultrabasic rocks (peridotites) leaves a mantle enriched in nickel (e.g.\ New Caledonia; nickel--cobalt laterites occur in the Lomi\'e region of East Cameroon).
\item \textbf{Residual iron and manganese:} iron and manganese can be concentrated in lateritic caps and weathered zones above lower-grade protores.
\end{itemize}

\begin{attention}[Bauxite is a rock, not a mineral]
A very frequent GCE error: \textbf{bauxite is a rock} --- a mixture of aluminium hydroxide minerals (gibbsite, boehmite, diaspore) plus impurities --- not a single mineral with a fixed formula. Similarly, laterite is a weathering product (a soil/rock type), not a mineral species.
\end{attention}

\courssub{Supergene Enrichment}

\textbf{Observation.} Many copper sulphide deposits are too poor to mine in their primary state, yet their upper parts have been naturally \emph{upgraded} by weathering. How? Descending rainwater dissolves copper from the oxidised top of the deposit and re-deposits it lower down, just below the water table, where it is safe from further oxidation.

\begin{definition}[Supergene enrichment]
\cle{Supergene enrichment} is the process by which a metal (especially copper) is leached from the oxidised upper part of a sulphide deposit, transported downward in solution, and re-precipitated as rich secondary sulphides (e.g.\ chalcocite \ce{Cu2S}, covellite \ce{CuS}) just below the \emph{water table}, thereby upgrading a low-grade primary ore into an economic one.
\end{definition}

\begin{popfigure}
\begin{tikzpicture}[scale=1.0, font=\small]
% zones
\fill[popmuted!30] (0,0) rectangle (8,1.6);
\fill[popgreen!25] (0,1.6) rectangle (8,2.8);
\fill[popgold!30] (0,2.8) rectangle (8,3.8);
\fill[popmuted!60] (0,3.8) rectangle (8,4.6);
\node at (4,0.8) {\textbf{Primary (hypogene) ore}: low-grade sulphides (chalcopyrite)};
\node[align=center] at (4,2.2) {\textbf{Enriched sulphide zone}:\\ secondary \ce{Cu2S}, \ce{CuS} --- high grade};
\node[align=center] at (4,3.3) {\textbf{Oxidised zone}: oxides, carbonates\\ (malachite, azurite), iron oxides};
\node[align=center, text=white] at (4,4.2) {\textbf{Leached zone (gossan)}: iron hat, metals removed};
% water table
\draw[popblue, very thick, dashed] (0,2.8) -- (8,2.8);
\node[popblue, anchor=west] at (8.1,2.8) {water table};
% arrows
\draw[->, popblue, very thick] (1.2,4.4) -- (1.2,2.0);
\node[popblue, align=center, text width=2.6cm] at (1.2,1.55) {Cu leached downward};
\draw[->, popgreen, very thick] (6.8,2.0) -- (6.8,1.2);
\node[popgreen, align=center, text width=2.8cm] at (6.8,3.0) {Cu re-deposited below water table};
\end{tikzpicture}
\legende{Supergene enrichment of a copper sulphide deposit: copper leached from the oxidised zone is re-precipitated as rich secondary sulphides just below the water table.}
\end{popfigure}

The process works because copper is soluble in oxidising, acidic surface waters but precipitates in the reducing environment below the water table. Without this natural upgrading, many of the world's great copper camps (e.g.\ in Zambia and the Congolese Copperbelt) would not have been economic.

\courssub{Metamorphic and Metamorphosed Deposits}

Finally, some deposits owe their economic value to \cle{metamorphism} --- the transformation of rocks by heat and pressure.

\begin{itemize}
\item \textbf{Metamorphosed deposits} are pre-existing ores modified by metamorphism: recrystallisation may improve quality (e.g.\ banded iron formations recrystallised into coarse, easily concentrated haematite ores).
\item \textbf{Metamorphic deposits} are new economic materials created by metamorphism itself: \cle{graphite} from carbon-rich sediments, \cle{asbestos} and \cle{talc} from ultrabasic rocks, \cle{marble} from limestone, garnet and kyanite from aluminous sediments.
\item \textbf{Skarn (contact-metasomatic) deposits} form where a hot intrusion invades carbonate rocks (limestone, dolomite): fluids from the magma react with the carbonate wall-rock, producing calcium--silicate minerals (garnet, pyroxene) together with ores of iron, copper, tungsten, lead and zinc.
\end{itemize}

\begin{definition}[Skarn]
A \cle{skarn} is a deposit formed at the contact between an igneous intrusion and a carbonate host rock, where hot magmatic fluids chemically replace the limestone or dolomite with calcium--magnesium silicate minerals and associated ores (magnetite, chalcopyrite, scheelite, sphalerite, galena).
\end{definition}

\courssub{Synthesis and Worked Exercise}

\begin{center}
\begin{tabularx}{\linewidth}{|l|X|X|X|}
\hline
\rowcolor{popblueL} \textbf{Genetic process} & \textbf{Typical products} & \textbf{Deposit morphology} & \textbf{Example} \\
\hline
Chemical precipitation (sedimentary) & Iron (BIF), salt, gypsum, potash & Extensive regular beds & BIF of the Kribi--Mbalam iron province \\
\hline
Biochemical accumulation & Phosphates & Beds and nodules in marine sediments & Phosphate basins of Morocco \\
\hline
Mechanical concentration (placer) & Gold, diamond, cassiterite, ilmenite & Gravel beds, terraces, beach sands & Gold placers of B\'etar\'e-Oya (East Cameroon) \\
\hline
Residual weathering (lateritic) & Bauxite, lateritic Ni, Fe, Mn & Blanket-like weathering mantle on plateaux & Minim--Martap / Ngaoundal bauxite (Adamaoua) \\
\hline
Supergene enrichment & Copper (secondary sulphides) & Horizontal enriched horizon below water table & Zambian Copperbelt \\
\hline
Metamorphic / skarn & Graphite, talc, marble; Fe, Cu, W ores & Irregular bodies at intrusion contacts & Skarns of the Urals (Russia) \\
\hline
\end{tabularx}
\end{center}

\begin{exoresolu}[Explaining the Adamaoua bauxite]
\textbf{Statement.} Using your knowledge of residual deposits, explain why economic bauxite occurs on the Adamaoua Plateau (Minim--Martap--Ngaoundal), and state two other products of tropical residual concentration.
\tcblower
\textbf{Solution.}
\begin{enumerate}
\item \textbf{Parent rock:} the plateau is underlain by aluminosilicate rocks (granites, syenites, gneisses) rich in aluminium-bearing feldspars.
\item \textbf{Climate:} the region experiences a humid tropical climate with high rainfall and temperature, causing intense chemical weathering by \ce{CO2}-charged rainwater.
\item \textbf{Process:} soluble elements (\ce{SiO2}, \ce{Ca}, \ce{Na}, \ce{K}) are leached downward and removed, while immobile aluminium accumulates as hydroxides (gibbsite, boehmite), producing a thick bauxitic laterite mantle.
\item \textbf{Preservation:} the plateau is a high, flat, well-drained and long-stable surface, so the weathering mantle was neither waterlogged nor eroded away, allowing a thick, high-grade bauxite blanket to accumulate.
\item \textbf{Other residual products:} lateritic nickel (from ultrabasic rocks) and residual iron or manganese.
\end{enumerate}
\end{exoresolu}

\begin{exercice}[Placer or vein?]
A prospector finds gold as (a) angular grains locked inside a quartz vein cutting across schists, and (b) rounded flakes in the basal gravel of a river terrace. For each occurrence, state the type of deposit, the process that formed it, and one feature that justifies your answer.
\end{exercice}

\begin{corrige}[Exercise 1]
(a) \textbf{Primary (hypogene) vein deposit:} the gold crystallised from hydrothermal fluids within a fracture; justification --- angular grains enclosed in vein quartz that cuts across the country rock (epigenetic geometry).
(b) \textbf{Alluvial placer deposit:} the gold was liberated by weathering of a primary source and mechanically concentrated by river action; justification --- rounded, water-worn flakes in a sedimentary gravel bed, at the base of the terrace where dense minerals settle.
\end{corrige}

\begin{aretenir}[Key points]
\begin{itemize}
\item Syngenetic sedimentary deposits (BIF, evaporites, phosphates) form as beds at the same time as their host sediments.
\item Placers concentrate dense, resistant minerals (Au, diamond, cassiterite) by hydraulic sorting in rivers and on beaches.
\item Lateritic weathering in humid tropical climates leaches \ce{Si}, \ce{Ca}, \ce{Na}, \ce{K} and residually enriches \ce{Al} and \ce{Fe}: this yields bauxite (Minim--Martap) and lateritic nickel.
\item Supergene enrichment upgrades copper ores by leaching above and re-deposition below the water table.
\item Metamorphism creates graphite, talc, marble and skarn ores at intrusion--limestone contacts.
\item \textbf{Exam tip:} explicitly linking Cameroon's humid tropical climate to lateritic bauxite (Adamaoua) and placer gold (East Region) earns marks in questions on national mineral wealth.
\end{itemize}
\end{aretenir}

\coursec{Occurrence of Mineral Deposits: From the World to Cameroon}

In the previous two sections we learned \emph{how} ore deposits form: by magmatic segregation, by hydrothermal circulation, by sedimentary concentration, by tropical weathering and by metamorphic remobilisation. A natural question now arises: \emph{where} do these processes actually operate on the Earth's surface, and can we predict where to look for a given metal? A prospector in the East Region of Cameroon does not dig at random: he pans for gold in streams draining ancient greenstone belts, because experience and theory tell him that this is where gold occurs. This final section answers the question of \cle{occurrence}: the geographical and geological distribution of mineral deposits, from the scale of the planet down to the mineral map of Cameroon.

\courssub{Geological Controls on the Occurrence of Deposits}

\textbf{Observation.} If you compare a world map of copper mines with a world map of gold mines, the two patterns do not coincide. Copper is concentrated along the western margins of the Americas and in central Africa; much of the world's gold comes from very old rocks in South Africa, Canada and Australia. This is not chance: each metal is tied to a particular \cle{geological setting}.

\begin{definition}[Mineral occurrence and prospect]
A \cle{mineral occurrence} is any place where a useful mineral substance is present in anomalous concentration, whether or not it can be mined profitably. A \cle{prospect} is an occurrence that has shown enough encouraging signs (good grades in samples, favourable geology) to justify detailed exploration. Only after full economic evaluation does a prospect become an \cle{ore deposit} (a mineable reserve).
\end{definition}

The occurrence of deposits is controlled by three great geological factors:

\begin{itemize}
\item \textbf{The type of crust and tectonic setting.} \cle{Cratons} (old, stable continental nuclei such as the Congo craton) host Archaean gold, banded iron formations and diamonds. \cle{Greenstone belts} within cratons concentrate gold and base metals. \cle{Orogenic belts} (mountain chains) host hydrothermal veins and metamorphic deposits. \cle{Rifts} and \cle{sedimentary basins} host evaporites, phosphates, petroleum and stratabound copper--lead--zinc.
\item \textbf{The age of the rocks.} Some deposits are linked to particular periods of Earth history: the great banded iron formations are essentially older than about $1.8$ billion years, because their deposition required an ocean--atmosphere system still poor in free oxygen; Witwatersrand-type gold is late Archaean; most porphyry copper deposits are geologically young (Mesozoic to Cenozoic).
\item \textbf{The tectonic and magmatic history.} A region that has experienced subduction, intrusion of granites, or rifting has been ``prepared'' by those events to receive the corresponding deposits: the heat, the fluids and the plumbing (faults, shear zones) are already in place.
\end{itemize}

\begin{propriete}[Fundamental principle of metallogeny]
The distribution of mineral deposits is controlled by \cle{geology} (rock type, age, tectonic history), \textbf{not by political boundaries}. The copper belt crosses the Zambia--DRC frontier; the gold-bearing greenstones continue from Cameroon into Gabon and the Central African Republic. (Not examinable as a proof, but must be understood and illustrated with examples.)
\end{propriete}

\courssub{Metallogenic Provinces and Metallogenic Epochs}

\textbf{Observation.} Why does a single region, the Zambian--Congolese copper belt, supply a large share of the world's copper and cobalt, while the Witwatersrand basin of South Africa has supplied a huge share of its gold? Because certain regions, endowed with the right rocks and the right history, concentrate certain metals.

\begin{definition}[Metallogenic province and epoch]
A \cle{metallogenic province} is a region within which mineral deposits of a particular type, or family of types, are unusually abundant because of a shared geological evolution. A \cle{metallogenic epoch} is a period of geological time during which a particular family of deposits formed preferentially worldwide.
\end{definition}

\begin{exemplebox}[Classic metallogenic provinces]
\begin{itemize}
\item \textbf{Copper belt of Zambia and the DRC:} stratabound copper--cobalt in Neoproterozoic sedimentary rocks of the Katanga basin.
\item \textbf{Witwatersrand (South Africa):} gold in late Archaean conglomerates (fossil placers), the richest gold province ever known.
\item \textbf{Banded iron formations:} vast iron provinces in the Hamersley basin (Australia), the Transvaal (South Africa), Minas Gerais (Brazil) and Krivoy Rog (Ukraine), all of comparable great (Palaeoproterozoic) age --- a fine example of a metallogenic \emph{epoch}.
\end{itemize}
\end{exemplebox}

\courssub{Plate Tectonics and Ore Deposits}

Plate tectonics provides the unifying explanation: each tectonic setting creates the conditions for a characteristic family of deposits.

\begin{itemize}
\item \textbf{Subduction zones:} rising magmas release metal-rich fluids, forming \cle{porphyry copper} (with Mo and Au) and epithermal gold veins, as along the Andes.
\item \textbf{Mid-ocean ridges and ophiolites:} slices of oceanic lithosphere obducted onto continents carry \cle{chromite} (podiform deposits) and Cyprus-type massive sulphide copper.
\item \textbf{Rift basins:} thick evaporite sequences (salt, potash), stratabound base metals, and the organic-rich sediments that generate \cle{petroleum}.
\item \textbf{Cratons and greenstone belts:} Archaean gold, banded iron, komatiite-hosted nickel, diamonds in kimberlite pipes.
\end{itemize}

\begin{popfigure}
\begin{tikzpicture}[font=\small]
% Ocean ridge
\fill[popblueL] (-0.5,0) rectangle (3.5,2.2);
\draw[popink] (-0.5,2.2) -- (0.8,1.2) -- (1.5,2.2) -- (2.2,1.2) -- (3.5,2.2);
\draw[->,thick,popdark] (0.9,1.0) -- (0.2,1.0);
\draw[->,thick,popdark] (2.1,1.0) -- (2.8,1.0);
\node[align=center] at (1.5,0.4) {\textbf{Ridge}\\ chromite, Cu};
% Subduction
\fill[poporange!25] (4.5,0) rectangle (8.0,2.2);
\draw[popink] (4.5,2.2) -- (6.2,2.2) -- (7.2,1.4) -- (8.0,0.6);
\draw[->,thick,popdark] (5.2,1.8) -- (6.4,1.1);
\node[align=center] at (6.2,0.35) {\textbf{Subduction}\\ porphyry Cu--Au};
% Rift
\fill[popgreenL] (9.0,0) rectangle (12.0,2.2);
\draw[popink] (9.0,2.2) -- (10.0,1.5) -- (11.0,1.5) -- (12.0,2.2);
\draw[->,thick,popdark] (10.2,1.3) -- (9.5,1.3);
\draw[->,thick,popdark] (10.8,1.3) -- (11.5,1.3);
\node[align=center] at (10.5,0.4) {\textbf{Rift basin}\\ evaporites, oil};
% Craton
\fill[popgold!30] (13.0,0) rectangle (16.0,2.2);
\draw[popink] (13.0,2.2) -- (16.0,2.2);
\node[align=center] at (14.5,1.1) {\textbf{Craton}\\ Au, Fe (BIF), Ni};
\end{tikzpicture}
\legende{Plate-tectonic settings and their characteristic deposit types.}
\end{popfigure}

\courssub{The Mineral Potential of Cameroon}

Cameroon lies at the junction of three great geological domains: the \cle{Congo craton} in the south (ancient, stable crust with greenstone belts), the \cle{Pan-African mobile belts} in the centre and north (deformed and intruded about $600$ million years ago), and the \cle{Cameroon Volcanic Line} in the west (a much younger volcanic axis running from the Atlantic Ocean towards the Adamaoua plateau). Each domain carries its own mineral family, exactly as the principle of metallogeny predicts.

\begin{popfigure}
\begin{tikzpicture}[font=\scriptsize,scale=0.95]
% Simplified outline of Cameroon
\draw[thick,popink] (2.0,0.5) -- (4.6,0.8) -- (5.4,2.2) -- (5.2,4.2) -- (4.4,5.6) -- (4.8,7.0) -- (3.6,8.6) -- (2.6,7.4) -- (1.6,6.2) -- (1.2,4.6) -- (0.6,3.2) -- (1.0,1.6) -- cycle;
% Domains
\fill[popgold!25] (1.0,0.6) -- (4.6,0.85) -- (5.3,2.6) -- (1.2,2.8) -- (0.65,3.1) -- (1.0,1.6) -- cycle;
\fill[popgreen!15] (1.2,2.8) -- (5.3,2.6) -- (5.2,4.2) -- (4.4,5.6) -- (2.2,5.4) -- (1.2,4.6) -- cycle;
\fill[popblue!12] (2.2,5.4) -- (4.4,5.6) -- (4.8,7.0) -- (3.6,8.6) -- (2.6,7.4) -- (1.6,6.2) -- cycle;
\node at (3.0,1.5) {Congo craton};
\node at (3.3,4.0) {Pan-African belts};
\node at (3.4,6.8) {Northern basins};
% Volcanic line
\draw[line width=2.5pt,poporange,opacity=0.7] (0.5,0.9) -- (1.3,2.2) -- (2.0,3.6) -- (2.6,5.0);
\node[poporange,align=center] at (0.4,3.4) {Volcanic\\ Line};
% Mineral symbols
\node[popdark] at (4.3,1.4) {$\blacktriangle$};
\node[align=left,anchor=west] at (5.6,1.4) {$\blacktriangle$ Au (East greenstones)};
\node[popdark] at (2.2,1.0) {$\blacksquare$};
\node[align=left,anchor=west] at (5.6,0.9) {$\blacksquare$ Fe: Mbalam, Nkout};
\node[popdark] at (3.4,4.6) {$\bigcirc$};
\node[align=left,anchor=west] at (5.6,4.6) {$\bigcirc$ Al (bauxite, Adamaoua)};
\node[popdark] at (3.0,2.4) {$\star$};
\node[align=left,anchor=west] at (5.6,2.4) {$\star$ Ni--Co: Nkamouna};
\node[popdark] at (2.6,3.0) {$\diamond$};
\node[align=left,anchor=west] at (5.6,3.3) {$\diamond$ Rutile: Akonolinga};
\node[popdark] at (4.2,5.4) {$\bullet$};
\node[align=left,anchor=west] at (5.6,5.4) {$\bullet$ Limestone: Figuil};
\node[popdark] at (1.0,1.9) {$\triangledown$};
\node[align=left,anchor=west] at (5.6,1.9) {$\triangledown$ Pozzolana (volcanic line)};
\node[align=left,anchor=west] at (5.6,6.2) {$\otimes$ Oil: Rio del Rey, Logone};
\node[popdark] at (1.0,0.9) {$\otimes$};
\node[popdark] at (3.9,7.4) {$\otimes$};
\end{tikzpicture}
\legende{Schematic map of Cameroon: main mineral occurrences plotted on simplified geological domains.}
\end{popfigure}

\begin{center}
\begin{tabularx}{\linewidth}{|l|l|X|}
\hline
\rowcolor{popblueL} \textbf{Substance} & \textbf{Locality} & \textbf{Genetic type} \\
\hline
Gold (Au) & East Region (B\'etar\'e-Oya, Batouri) & Hydrothermal veins in greenstone belts; placers \\
\hline
Iron (Fe) & Mbalam, Nkout (South) & Banded iron formation (itabirite), enriched \\
\hline
Bauxite (Al) & Adamaoua plateau (Ngaoundal, Minim-Martap) & Tropical lateritic weathering \\
\hline
Nickel--cobalt & Nkamouna (East) & Lateritic weathering of ultramafic rocks \\
\hline
Rutile (\ce{TiO2}) & Akonolinga (Centre) & Residual and alluvial concentration \\
\hline
Limestone & Figuil (North) & Sedimentary (marine carbonate) \\
\hline
Pozzolana & Volcanic line (West) & Volcanic scoria and ash (industrial mineral) \\
\hline
Petroleum & Rio del Rey, Logone basins & Sedimentary basin fill (fuel mineral) \\
\hline
\end{tabularx}
\end{center}

\begin{savaistu}[Fuel minerals of Cameroon]
Petroleum, exploited offshore in the \cle{Rio del Rey} basin near Limbe and onshore in the \cle{Logone} basin of the Far North, belongs to the family of \cle{fuel minerals} (fossil energy resources), studied separately from metallic ores. It has long been the backbone of Cameroon's export earnings.
\end{savaistu}

\courssub{Reading a Mineral Occurrence Map}

A mineral occurrence map plots symbols (each representing a substance) on a geological background. The skill tested at the GCE is to \emph{relate} the symbols to the geology.

\begin{attention}[Do not read symbols in isolation]
Never interpret a mineral symbol without relating it to the underlying geology. A gold symbol means little alone; a gold symbol \emph{inside a greenstone belt cut by granitic intrusions} tells a coherent genetic story and points to prospective ground.
\end{attention}

\begin{methode}[Interpreting a mineral occurrence map]
\begin{enumerate}
\item Read the legend: identify each symbol and the substance it represents.
\item Study the geological background: rock types, ages, structures (faults, fold belts).
\item Note clusters of occurrences and ask which geological unit they follow.
\item Link each cluster to a genetic process (e.g.\ bauxite on a tropical plateau = lateritisation).
\item Identify \cle{prospective zones}: areas with the same favourable geology but few recorded occurrences.
\end{enumerate}
\end{methode}

\courssub{From Occurrence to Exploitation}

\begin{methode}[Stages of mineral exploration]
\begin{enumerate}
\item \textbf{Regional survey:} study of maps, satellite images and stream-sediment sampling over a large area.
\item \textbf{Prospecting:} fieldwork, trenching and pitting on anomalies; the occurrence becomes a prospect.
\item \textbf{Geophysics and geochemistry:} magnetic, electrical and soil-geochemical surveys to locate hidden ore.
\item \textbf{Drilling:} boreholes give samples at depth; grades and thicknesses are measured.
\item \textbf{Reserve estimation and feasibility:} tonnage and grade are computed; only then is the prospect declared an exploitable \cle{deposit}.
\end{enumerate}
\end{methode}

\begin{attention}[An occurrence is not a deposit]
A mineral showing, however spectacular, is not yet a mine. Economic evaluation --- \cle{grade}, \cle{tonnage}, accessibility, market price, feasibility --- must come first. Confusing ``occurrence'' with ``reserve'' is a classic GCE error.
\end{attention}

Mining itself uses two great families of methods: \cle{open-pit} mining for shallow, large or low-grade bodies (lateritic bauxite and nickel, iron of Mbalam) and \cle{underground} mining for deep, high-grade bodies (vein gold). Finally, exploitation raises environmental and socio-economic stakes: uncontrolled \cle{artisanal mining} (as in the gold fields of the East Region) degrades rivers and soils, while responsible projects plan \cle{mine rehabilitation} (re-vegetation, backfilling) from the start.

\begin{exoresolu}[Interpreting a prospecting report]
A prospecting team reports: (i) gold values up to $4\ \mathrm{g/t}$ in quartz veins cutting a greenstone belt near Batouri; (ii) a bauxite cap averaging $45\ \%$ \ce{Al2O3} over a wide plateau near Ngaoundal. Classify each result (occurrence, prospect or reserve) and justify; state the genetic type of each.
\tcblower
\textbf{Solution.}
\begin{itemize}
\item \textbf{Batouri gold:} encouraging grades in veins within a favourable greenstone setting, but no drilling or tonnage estimate is mentioned: it is a \cle{prospect} (an occurrence under active exploration). Genetic type: hydrothermal (mesothermal) gold veins, typical of greenstone belts.
\item \textbf{Ngaoundal bauxite:} high grade ($45\ \%$ \ce{Al2O3}) over a wide plateau suggests large tonnage, but without a full feasibility study it remains a well-advanced \cle{prospect}, not yet a proven reserve. Genetic type: residual deposit from tropical lateritic weathering of aluminous rocks.
\item \textbf{Key point:} neither result is yet a \emph{reserve}, because reserves require demonstrated tonnage, grade and economic feasibility.
\end{itemize}
\end{exoresolu}

\begin{exercice}[Cameroon's mineral wealth]
Name three Cameroonian mineral deposits, giving for each the substance, the locality and the genetic type. Then explain, in five lines, why the East Region is prospective for gold.
\end{exercice}

\begin{corrige}[Exercise 1]
Examples: gold at B\'etar\'e-Oya (East) --- hydrothermal veins in greenstone belts; iron at Mbalam (South) --- enriched banded iron formation; bauxite at Minim-Martap (Adamaoua) --- lateritic weathering; nickel--cobalt at Nkamouna --- laterite on ultramafics; limestone at Figuil --- marine sedimentary. The East Region is prospective for gold because it contains Archaean-to-Proterozoic greenstone belts cut by granitic intrusions and shear zones: the intrusions supplied heat and fluids, the shears channelled them, and gold precipitated in quartz veins, exactly the setting of the world's great gold provinces.
\end{corrige}

\begin{aretenir}[Exam tip]
Memorise at least \textbf{three Cameroonian deposits} with their substance, location and genetic type (e.g.\ Au--B\'etar\'e-Oya--hydrothermal; Fe--Mbalam--BIF; Al--Adamaoua--lateritic; Ni--Co--Nkamouna--lateritic; limestone--Figuil--sedimentary). This is a classic GCE question. Remember also: deposits follow \emph{geology}, not borders; an occurrence becomes a reserve only after economic evaluation; and always read a mineral symbol against the geological background.
\end{aretenir}

\newpage

\coursec{Integration activity}

\courssub{Aims of the activity}

This integration activity closes the chapter by asking you to act as a \cle{junior exploration geologist}. Working in teams, you will combine every notion studied in this chapter to answer one practical question: \emph{is a real iron deposit in southern Cameroon worth mining today?} By the end of the activity you should be able to:

\begin{itemize}
\item recognise the \cle{genetic type} and \cle{morphology} of a deposit from map, cross-section and grade data;
\item compute the \cle{contained metal} in a deposit and the \cle{reserve life} for a given production rate;
\item distinguish \cle{inferred} and \cle{indicated resources} from a \cle{proved reserve}, and state what additional work is needed to move from one category to the next;
\item apply the notion of \cle{economic mineral} by weighing grade, world market price, mining cost, technology and infrastructure;
\item defend a reasoned recommendation orally, including one \cle{environmental safeguard}, before a mock ``Ministry of Mines'' panel.
\end{itemize}

\begin{attention}[Skills being assessed]
This activity mobilises all six sections of the chapter: the notion of economic mineral, resources and reserves, classification of deposits, magmatic/hydrothermal and sedimentary/weathering genesis, and the occurrence of deposits in Cameroon. Marks are awarded for correct \emph{reasoning}, not for a single ``right'' verdict: two teams may reach opposite conclusions and both score highly if their arguments are rigorous. Every numerical answer must carry its \emph{unit}, and every conclusion must be tied to a named notion from the course.
\end{attention}

\courssub{Situation}

\textbf{Dossier GEO05--NK: The Nkout Iron Project (South Region, Cameroon)}\par

You are employed by a small exploration company, \emph{Mvila Minerals SARL}. The company holds an exploration permit near Nkout, in the South Region, about $\SI{290}{km}$ inland from the Atlantic coast (Kribi). The Ministry of Mines, Industry and Technological Development has asked your company for an independent technical note before deciding whether to encourage a full feasibility study. Your team must write that note.

The dossier in your hands contains the following simplified data.

\textbf{1. Geological setting.} The permit lies within the \cle{Congo craton}, on Archaean to Palaeoproterozoic basement (gneisses and greenstone belts, older than about 2.5~Ga). The geological map shows a ridge of \cle{banded iron formation} (BIF) --- alternating millimetric bands of quartz and iron oxides --- traceable for about $\SI{8}{km}$ along strike. Where the BIF crops out at the surface, deep tropical weathering has removed the silica and oxidised the iron minerals, producing a soft, reddish, enriched cap of \cle{hematite} and \cle{goethite} (a \emph{supergene enrichment} developed on the BIF protore). Fresh BIF at depth contains only about $30\ \%$ Fe; the enriched cap contains $55$--$60\ \%$ Fe.

\textbf{2. Morphology and dimensions (from 40 drill holes).} The enriched orebody forms a \cle{blanket} (a sub-horizontal \emph{stratabound} layer) capping the ridge:
\begin{itemize}
\item average thickness: $\SI{40}{m}$;
\item lateral extent: $\SI{8000}{m}$ (strike) $\times$ $\SI{500}{m}$ (width);
\item average dry density of the ore: $\rho = 3.2\ \mathrm{t\,m^{-3}}$;
\item average grade of the enriched cap: $45\ \%$ Fe (drill-hole grades range from $30\ \%$ Fe at the edges to $60\ \%$ Fe at the hilltop).
\end{itemize}

\textbf{3. Drilling pattern.} On the central $\SI{4}{km}$ of the ridge, holes are spaced $\SI{200}{m}$ apart (close enough to correlate the ore layer confidently from hole to hole). On the outer $\SI{4}{km}$, only widely spaced holes ($\geq \SI{800}{m}$) exist, so continuity is assumed but not demonstrated. Roughly half of the total tonnage lies in each zone.

\textbf{4. Economic data.}
\begin{itemize}
\item world price of iron ore (62\,\% Fe benchmark): about $100$ US dollars per tonne of ore;
\item estimated mining and processing cost (open pit, crushing and upgrading to a saleable concentrate): $45$ US dollars per tonne of ore;
\item there is \textbf{no railway}; the ore would have to be trucked $\SI{290}{km}$ to Kribi, at an estimated transport cost of $35$ US dollars per tonne;
\item planned production rate if mined: $5$ million tonnes of ore per year;
\item the deep-water port of Kribi is operational and can handle bulk carriers.
\end{itemize}

\textbf{5. Environmental and social notes.} The ridge lies in the Congo Basin rainforest, near the Dja Faunal Reserve (a UNESCO World Heritage site). A river crossing the permit supplies water to three villages downstream.

\begin{popfigure}
\begin{center}
\begin{tikzpicture}[x=1cm,y=1cm,scale=0.9]
% topography
\draw[thick,popink] (0,0) -- (2,0) .. controls (3.5,0) and (4,1.6) .. (5.5,1.8) .. controls (7,1.6) and (7.5,0) .. (9,0) -- (11,0);
% enriched cap
\draw[fill=poporange,draw=popdark,thick] (3.1,0.35) .. controls (4,1.75) and (7,1.75) .. (7.9,0.35) .. controls (7,1.15) and (4,1.15) .. (3.1,0.35) -- cycle;
% BIF below
\draw[fill=popblueL,draw=popblue,thick] (3.1,0.35) .. controls (4,1.15) and (7,1.15) .. (7.9,0.35) -- (7.9,-0.9) .. controls (7,0.1) and (4,0.1) .. (3.1,-0.9) -- cycle;
% basement
\draw[fill=gray!15] (0,0) -- (3.1,0.35) -- (3.1,-0.9) .. controls (4,0.1) and (7,0.1) .. (7.9,-0.9) -- (7.9,0.35) -- (11,0) -- (11,-2.2) -- (0,-2.2) -- cycle;
% drill holes
\draw[thick,popdark] (4.6,1.55) -- (4.6,-1.4);
\draw[thick,popdark] (5.5,1.8) -- (5.5,-1.4);
\draw[thick,popdark] (6.4,1.55) -- (6.4,-1.4);
\draw[thick,popmuted,dashed] (1.6,0.05) -- (1.6,-1.4);
\draw[thick,popmuted,dashed] (9.6,0.05) -- (9.6,-1.4);
% labels
\node[align=center,font=\small] at (5.5,1.15) {enriched cap\\ $45$--$60\ \%$ Fe};
\node[align=center,font=\small] at (5.5,-0.45) {fresh BIF\\ $\approx 30\ \%$ Fe};
\node[font=\small] at (2,-1.7) {Archaean basement (Congo craton)};
\node[font=\small,align=center] at (5.5,-1.7) {close-spaced drilling ($\SI{200}{m}$): indicated zone};
\node[font=\small,align=center] at (1.6,-2.6) {wide spacing:\\ inferred zone};
\node[font=\small,align=center] at (9.6,-2.6) {wide spacing:\\ inferred zone};
% scale arrow
\draw[<->,thick] (3.1,2.3) -- (7.9,2.3) node[midway,above,font=\small] {strike $\approx \SI{8}{km}$ (schematic)};
\end{tikzpicture}
\end{center}
\legende{Schematic cross-section of the Nkout ridge: a supergene-enriched hematite--goethite blanket (orange) caps the banded iron formation (blue), which rests in Archaean basement. Solid drill holes: close spacing in the central indicated zone; dashed holes: wide spacing in the inferred zones.}
\end{popfigure}

\courssub{Tasks}

Answer the following questions in order. Show all calculations and justify every conclusion with a notion from the chapter.

\begin{enumerate}
\item \textbf{Genesis and classification.} Using the dossier, identify the genetic type of the Nkout deposit. Name (a) the original sedimentary process that formed the BIF, and (b) the later process that raised the grade from $30\ \%$ to $45$--$60\ \%$ Fe. Classify the deposit by its \emph{morphology} and state whether it is syngenetic or epigenetic with respect to the host BIF.
\item \textbf{Tonnage and contained metal.} Calculate the total tonnage of enriched ore, then the total tonnage of contained iron metal.
\item \textbf{Resource classification.} Using the drilling pattern, divide the tonnage between \emph{indicated resource} and \emph{inferred resource}. Explain, in two or three sentences, what additional work would be required before any of this material could be declared a \emph{proved reserve}.
\item \textbf{Reserve life.} Assuming the whole tonnage were mined at $5$~Mt of ore per year, calculate the life of the mine.
\item \textbf{Economic mineral?} Compute the total cost per tonne of ore delivered to the ship (mining $+$ processing $+$ transport) and compare it with the world price. Is the ore an \emph{economic mineral} under present conditions? Discuss how the answer could change with (a) the world iron price, (b) construction of a railway, and (c) cut-off grade policy.
\item \textbf{Recommendation.} Write a one-page note to the Ministry panel: your verdict (mine now / keep exploring / wait), the two strongest arguments for it, and \textbf{one concrete environmental safeguard} your company would impose.
\item \textbf{Oral defence.} Present your note in five minutes and answer the panel's questions. Be ready to defend your resource classification and your economic threshold.
\end{enumerate}

\begin{methode}[How to attack a tonnage--grade problem]
\begin{enumerate}
\item Compute the \emph{volume} of the orebody from its geometry: $V = \text{length} \times \text{width} \times \text{thickness}$.
\item Convert volume to \emph{tonnage} using the density: $M_{\text{ore}} = V \times \rho$. Never skip this step.
\item Compute the \emph{contained metal}: $M_{\text{metal}} = M_{\text{ore}} \times \text{grade}$ (grade as a decimal fraction).
\item Classify the tonnage by \emph{confidence} (drilling density) before calling anything a reserve.
\item Only then discuss economics: compare \emph{total delivered cost} with the \emph{sale price}.
\end{enumerate}
\end{methode}

\courssub{Model answer}

\textbf{1. Genesis and classification.} The deposit is an \cle{enriched banded iron formation} (BIF-type iron deposit). (a) The original BIF is a \emph{chemical sedimentary} deposit: iron and silica were precipitated from seawater in alternating bands on an Archaean--Palaeoproterozoic sea floor, so the primary iron concentration is of \cle{sedimentary (chemical precipitation)} origin. (b) The high-grade cap results from \cle{supergene enrichment by tropical weathering}: lateritic weathering dissolved and leached the quartz bands (residual concentration) and oxidised the iron minerals to hematite and goethite, roughly doubling the grade. Morphologically the orebody is a \emph{stratabound blanket} (a horizontal layer) capping the ridge. The primary BIF is \cle{syngenetic} (formed at the same time as the enclosing sediments); the enriched cap is \cle{epigenetic} with respect to the BIF, since the enrichment post-dates the rock.

\textbf{2. Tonnage and contained metal.}
\begin{align*}
V &= 8000 \times 500 \times 40 = 1.6\times 10^{8}\ \mathrm{m^{3}}\\
M_{\text{ore}} &= V \times \rho = 1.6\times 10^{8} \times 3.2 = 5.12\times 10^{8}\ \mathrm{t} \approx 512\ \mathrm{Mt}\\
M_{\text{Fe}} &= 0.45 \times 512 \approx 230\ \mathrm{Mt\ of\ contained\ iron}
\end{align*}

\textbf{3. Resource classification.} The central $\SI{4}{km}$, drilled at $\SI{200}{m}$ spacing where continuity of the layer between holes is reasonably demonstrated, constitutes an \cle{indicated resource}: about half the tonnage, i.e.\ $\approx 256\ \mathrm{Mt}$. The outer $\SI{4}{km}$, where continuity is only assumed from widely spaced holes, is an \cle{inferred resource}: $\approx 256\ \mathrm{Mt}$. To declare a \cle{proved reserve}, the company would need: (i) closer-spaced drilling (e.g.\ $50$--$\SI{100}{m}$) to demonstrate continuity in three dimensions; (ii) metallurgical tests proving the ore can be processed at the assumed cost; (iii) a full feasibility study including mining plan, infrastructure, market and environmental studies --- because a \emph{reserve} is the part of a resource that is \emph{economically mineable}, not merely present.

\begin{center}
\begin{tabularx}{\linewidth}{|l|X|X|X|}
\hline
\rowcolor{popblueL} \textbf{Zone} & \textbf{Drilling} & \textbf{Category} & \textbf{Tonnage} \\
\hline
Central $\SI{4}{km}$ & $\SI{200}{m}$ spacing, continuity demonstrated & Indicated resource & $\approx 256\ \mathrm{Mt}$ \\
\hline
Outer $\SI{4}{km}$ & $\geq \SI{800}{m}$ spacing, continuity assumed & Inferred resource & $\approx 256\ \mathrm{Mt}$ \\
\hline
Whole ridge & after feasibility study only & (Potential) proved reserve & to be confirmed \\
\hline
\end{tabularx}
\end{center}

\textbf{4. Reserve life.}
\[
\text{Life} = \dfrac{512\ \mathrm{Mt}}{5\ \mathrm{Mt\,yr^{-1}}} \approx 102\ \text{years}.
\]
Even using only the indicated tonnage ($256\ \mathrm{Mt}$), the life exceeds $50$ years --- comfortably long for a mining project.

\textbf{5. Economic mineral?} Total delivered cost per tonne of ore:
\[
45 + 35 = 80\ \text{USD/t} \quad\text{against a sale price of } 100\ \text{USD/t}.
\]
The gross margin is only about $20\ \text{USD/t}$, and this ignores capital costs (building the mine, the plant and the road), taxes and the fact that a $45\ \%$ Fe ore must first be upgraded to near $62\ \%$ Fe to fetch the benchmark price. The deposit is therefore \emph{marginally economic at best} under present conditions: it is a \cle{resource}, but not yet clearly an \cle{economic mineral}. (a) If the world price rose above roughly $120$--$130\ \text{USD/t}$, the margin would become comfortable. (b) A railway (as long discussed for the Mbalam--Nabeba--Kribi corridor) could cut transport from $35$ to perhaps $10\ \text{USD/t}$, transforming the economics. (c) Raising the \cle{cut-off grade} to mine only the $55$--$60\ \%$ Fe hilltop material would reduce tonnage but increase margin per tonne.

\textbf{6. Recommendation (model).} \emph{Verdict: keep exploring --- do not mine yet.} Arguments: (1) the tonnage and grade are real and large, but half of it is only an inferred resource and the economic margin ($\approx 20\ \text{USD/t}$ before capital costs) is too thin to justify investment without a railway or a higher iron price; (2) the indicated $256\ \mathrm{Mt}$ alone justifies further drilling and metallurgical testing to convert resources into proved reserves. \emph{Environmental safeguard:} impose a protected buffer zone of at least $\SI{100}{m}$ of untouched forest along the river, with continuous water-quality monitoring downstream, so that the three villages' water supply is never contaminated by runoff from the open pit and waste dumps.

\begin{attention}[Marking pointers]
Frequent errors: forgetting the density when converting volume to tonnage; calling the whole $512\ \mathrm{Mt}$ a ``reserve'' (it is a \emph{resource} --- no feasibility study exists); confusing the syngenetic BIF with the epigenetic enrichment; and concluding ``economic'' from grade alone while ignoring transport and infrastructure. A verdict of ``wait'' or ``mine'' is acceptable --- only the quality of the argument is marked.
\end{attention}

\begin{savaistu}[Iron in Cameroon: a real national stake]
The Nkout deposit is not a fiction: it belongs to a belt of BIF-hosted iron occurrences (Mbalam, Nkout, Djadom and others) along the northern edge of the Congo craton in the South and East Regions of Cameroon. Their development has been discussed for decades, and the missing link has always been the same one you identified in this activity: a railway to the deep-water port of Kribi. This is a perfect illustration of the chapter's central message --- a \emph{mineral deposit} becomes an \emph{economic mineral} only when geology, technology, infrastructure and market price all align.
\end{savaistu}

\newpage

\coursec{Methods and worked examples}

\coursec{Economic Geology: Economic Minerals}

\courssub{Skill 1: Distinguishing ore, gangue, grade and economic mineral}

\begin{definition}[Economic mineral, ore mineral, gangue, grade]
An \cle{economic mineral} is a mineral from which a valuable metal or substance can be extracted \emph{profitably} under current economic, technological and legal conditions. The \cle{ore mineral} is the specific mineral species that carries the metal of value (e.g.\ cassiterite \ce{SnO2} for tin, chalcopyrite \ce{CuFeS2} for copper). The \cle{gangue} is the assemblage of valueless minerals (quartz, calcite, feldspars, micas\ldots) that are extracted together with the ore mineral and must be separated during processing. The \cle{grade} (or tenor) expresses the concentration of the valuable substance in the rock, usually in percent (\%), grams per tonne (g/t) or parts per million (ppm); $1\ \mathrm{g/t} = 1\ \mathrm{ppm}$.
\end{definition}

\begin{methode}[Identifying the components of an ore and computing a grade]
When faced with a description of a mined rock, proceed in order:
\begin{enumerate}
\item \textbf{Identify the useful mineral(s)}: the mineral species that contains the metal or substance of value (e.g.\ cassiterite \ce{SnO2} for tin).
\item \textbf{Identify the gangue}: the valueless minerals extracted with the ore (quartz, calcite, feldspar\ldots) which must be removed during processing.
\item \textbf{Decide whether the mineral is ``economic''}: ask the three questions — is it \emph{concentrated enough} (grade), is it \emph{present in sufficient quantity} (tonnage), and can it be \emph{extracted profitably} with existing technology and market prices?
\item \textbf{Compute the grade} if data are given:
\[
\text{grade (\%)} = \frac{\text{mass of valuable substance}}{\text{total mass of rock}} \times 100
\]
\item \textbf{Convert between mineral grade and metal grade} using the chemical formula and molar masses when required.
\end{enumerate}
\textbf{Reflexes:} always state the unit of grade (\%, g/t, ppm); remember that $1\ \mathrm{g/t} = 1\ \mathrm{ppm}$; an ``economic mineral'' is defined by \emph{economics}, not only by chemistry — the same mineral may be ore in one place and waste in another.
\end{methode}

\begin{attention}[Common mistakes in grade calculations]
\begin{itemize}
\item Forgetting to convert \% to a decimal fraction before multiplying by tonnage (e.g.\ using $5$ instead of $0.05$).
\item Confusing \emph{mineral grade} (e.g.\ \% \ce{SnO2}) with \emph{metal grade} (e.g.\ \% Sn) — always specify which you are quoting.
\item Using the wrong molar mass or forgetting that the metal fraction in the mineral is $M_{\text{metal}} / M_{\text{mineral}}$.
\item Omitting units in the final answer (grade must carry \%, g/t or ppm).
\end{itemize}
\end{attention}

\begin{exoresolu}[Cassiterite ore from the Mayo Darl\'e tin field]
A prospector in the Mayo Darl\'e area (Adamawa Region, Cameroon) collects a $500\ \mathrm{kg}$ sample of mineralised granite. After concentration, the sample yields $8\ \mathrm{kg}$ of cassiterite (\ce{SnO2}).
\begin{enumerate}
\item[a)] Define the terms \emph{ore mineral} and \emph{gangue}, and identify them in this sample.
\item[b)] Calculate the grade of the rock in cassiterite.
\item[c)] Given $M(\mathrm{Sn}) = 118.7\ \mathrm{g/mol}$ and $M(\mathrm{O}) = 16.0\ \mathrm{g/mol}$, calculate the grade of the rock in metallic tin.
\end{enumerate}
\tcblower
\textbf{a)} An \cle{ore mineral} is a mineral from which a metal or useful substance can be extracted \emph{profitably}. The \cle{gangue} is the assemblage of valueless minerals (here quartz, feldspars and micas of the granite) extracted together with the ore mineral and discarded during processing. In this sample the ore mineral is \textbf{cassiterite \ce{SnO2}}; the gangue is the granitic material (quartz, feldspar, mica).

\textbf{b)} Grade in cassiterite:
\[
\text{grade}_{\ce{SnO2}} = \frac{8\ \mathrm{kg}}{500\ \mathrm{kg}} \times 100 = 1.6\ \%
\]
The rock contains $1.6\ \%$ cassiterite, i.e.\ $16\ \mathrm{kg}$ of \ce{SnO2} per tonne of rock.

\textbf{c)} Molar mass of cassiterite:
\[
M(\ce{SnO2}) = 118.7 + 2 \times 16.0 = 150.7\ \mathrm{g/mol}
\]
Mass fraction of tin in cassiterite:
\[
w_{\mathrm{Sn}} = \frac{118.7}{150.7} = 0.7876 \approx 78.8\ \%
\]
Mass of tin in the sample:
\[
m_{\mathrm{Sn}} = 8\ \mathrm{kg} \times 0.7876 = 6.30\ \mathrm{kg}
\]
Grade in metallic tin:
\[
\text{grade}_{\mathrm{Sn}} = \frac{6.30\ \mathrm{kg}}{500\ \mathrm{kg}} \times 100 = 1.26\ \%
\]
\textbf{Conclusion:} the rock grades $1.6\ \%$ \ce{SnO2}, equivalent to about $1.26\ \%$ Sn. Since typical economic tin ores grade around $0.5$--$1\ \%$ Sn, this sample is of economic interest.
\end{exoresolu}

\begin{aretenir}[Key points: ore, gangue, grade]
\begin{itemize}
\item Economic mineral = profitable extraction (geology + economics + technology).
\item Ore mineral = carrier of the metal; gangue = waste minerals.
\item Grade = concentration; units: \%, g/t, ppm ($1\ \mathrm{g/t}=1\ \mathrm{ppm}$).
\item Metal grade = mineral grade $\times$ (metal molar mass / mineral molar mass).
\end{itemize}
\end{aretenir}

\courssub{Skill 2: Classifying a deposit as resource or reserve}

\begin{definition}[Mineral resource and mineral reserve (CRIRSCO/JORC framework)]
A \cle{mineral resource} is a concentration of solid material of economic interest in or on the Earth's crust in such form, grade and quantity that there are reasonable prospects for eventual economic extraction. It is subdivided by geological confidence into \emph{Inferred}, \emph{Indicated} and \emph{Measured} resources (increasing confidence). A \cle{mineral reserve} is the economically mineable part of a Measured and/or Indicated Mineral Resource. It includes diluting materials and allowances for losses, and is subdivided into \emph{Probable} and \emph{Proved} reserves. \textbf{All reserves are resources, but not all resources are reserves.} The boundary is drawn by economics and technology, not by geology alone.
\end{definition}

\begin{methode}[Applying the resource/reserve terminology]
\begin{enumerate}
\item \textbf{Read the level of geological knowledge}: is the deposit only \emph{inferred} (few data, wide spacing of samples), \emph{indicated} (reasonably well sampled) or \emph{measured} (densely sampled, well known)?
\item \textbf{Ask the economic question}: can the material be extracted \emph{legally and profitably at present} with existing technology?
\item \textbf{Place the material}:
\begin{itemize}
\item \textbf{Resource}: concentration of economic interest, known or inferred, \emph{regardless of present profitability}. Subdivided into inferred, indicated, measured resources by increasing confidence.
\item \textbf{Reserve}: the part of a measured or indicated resource that is \emph{economically mineable now}. Subdivided into \emph{probable} and \emph{proved} reserves.
\end{itemize}
\item \textbf{Watch for changes of status}: a resource becomes a reserve if the metal price rises, technology improves, or infrastructure arrives; a reserve can revert to a resource if prices collapse.
\end{enumerate}
\textbf{Key idea to remember:} \emph{All reserves are resources, but not all resources are reserves.} The boundary is drawn by economics and technology, not by geology alone.
\end{methode}

\begin{attention}[Resource vs Reserve — examiner's traps]
\begin{itemize}
\item Never call a deposit a ``reserve'' just because it is large or well known. Profitability \emph{now} is mandatory.
\item ``Measured resource'' $\neq$ ``Proved reserve''. A measured resource becomes a proved reserve \emph{only after} a feasibility study demonstrates economic viability.
\item Inferred resources can \emph{never} be converted directly to reserves; they must first become indicated/measured.
\end{itemize}
\end{attention}

\begin{exoresolu}[The bauxite of Minim-Martap and Ngaoundal]
The bauxite plateaux of Minim-Martap and Ngaoundal (Adamawa, Cameroon) are estimated to contain very large tonnages of aluminium-rich laterite. However, the deposits lie far from the coast, the alumina grade is moderate, and no railway currently serves the plateaux, so no mining has started.
\begin{enumerate}
\item[a)] Define \emph{mineral resource} and \emph{mineral reserve}.
\item[b)] Should the Minim-Martap bauxite be reported as a reserve? Justify.
\item[c)] State two changes that could convert this material into a reserve.
\end{enumerate}
\tcblower
\textbf{a)} A \cle{mineral resource} is a natural concentration of solid material of economic interest in or on the Earth's crust, whose grade and tonnage give reasonable prospects of eventual economic extraction; it is classified as inferred, indicated or measured according to geological confidence. A \cle{mineral reserve} is the portion of a measured or indicated resource that is \emph{economically, legally and technically mineable at the time of reporting}; it is classified as probable or proved.

\textbf{b)} No. Although the tonnage is enormous and the geology is reasonably well known (so it is at least an \emph{indicated resource}), the deposit cannot presently be mined profitably: transport costs to a port would exceed the value of the ore, and the moderate grade increases processing costs. Therefore it fails the economic test and must be reported as a \textbf{resource}, not a reserve.

\textbf{c)} Any two of:
\begin{itemize}
\item construction of a railway or improved road linking the plateaux to the deep-water port (e.g.\ Kribi), cutting transport costs;
\item a sustained rise in the world price of aluminium/alumina;
\item improved, cheaper beneficiation technology raising the recoverable grade;
\item construction of a local alumina refinery or cheap energy source (e.g.\ hydroelectricity from the Lom Pangar dam area) reducing processing costs.
\end{itemize}
\textbf{Examiner's remark:} this example shows that ``reserve'' is a \emph{dynamic, economic} category: the geology of Minim-Martap does not change, but its status can.
\end{exoresolu}

\begin{aretenir}[Resource / Reserve hierarchy]
\begin{center}
\begin{tabular}{|l|l|l|}
\hline
\textbf{Category} & \textbf{Geological confidence} & \textbf{Economic viability} \\
\hline
Inferred Resource & Low & Not required \\
Indicated Resource & Moderate & Not required \\
Measured Resource & High & Not required \\
\hline
Probable Reserve & High (Indicated) & Demonstrated \\
Proved Reserve & Very high (Measured) & Demonstrated \\
\hline
\end{tabular}
\end{center}
\end{aretenir}

\courssub{Skill 3: Computing tonnage, contained metal and cut-off grade}

\begin{methode}[Deposit evaluation calculations]
\begin{enumerate}
\item \textbf{Compute the volume} of the ore body from its geometry (thickness $\times$ surface area for a tabular body; $\pi r^2 h$ for a cylinder; etc.).
\item \textbf{Compute the tonnage}: $T = V \times \rho$ where $\rho$ is the bulk density of the ore (in $\mathrm{t/m^3}$).
\item \textbf{Compute contained metal}: $M = T \times g$ where $g$ is the grade expressed as a decimal fraction (for g/t grades, $M$ in grams comes out directly: $M(\mathrm{g}) = T(\mathrm{t}) \times g(\mathrm{g/t})$).
\item \textbf{Apply a cut-off grade}: the \cle{cut-off grade} is the minimum grade at which a block of rock pays for its extraction and processing. Material below cut-off is waste. Compare each block's grade with the cut-off before summing reserves.
\item \textbf{Check profitability} with a simple balance: value of recovered metal $=$ contained metal $\times$ recovery $\times$ price; compare with extraction + processing costs.
\end{enumerate}
\textbf{Reflexes:} convert \% to decimals before multiplying; keep units consistent ($\mathrm{m^3} \times \mathrm{t/m^3} = \mathrm{t}$); never confuse \emph{contained metal} with \emph{recoverable metal} (multiply by the recovery factor).
\end{methode}

\begin{attention}[Cut-off grade and recovery]
\begin{itemize}
\item The cut-off grade is \emph{not} a geological constant — it depends on costs, metal price, recovery and dilution.
\item ``Contained metal'' is the total metal in the ground; ``Recoverable metal'' = contained $\times$ metallurgical recovery (always $< 100\%$).
\item In reserve estimation, only blocks \emph{above} cut-off are summed; blocks below cut-off are waste even if they contain some metal.
\end{itemize}
\end{attention}

\begin{exoresolu}[Evaluating a small gold vein]
A quartz vein near B\'etar\'e Oya (East Region, Cameroon) is $200\ \mathrm{m}$ long, $2\ \mathrm{m}$ wide on average and is mined to a depth of $50\ \mathrm{m}$. The bulk density of the ore is $2.7\ \mathrm{t/m^3}$ and the average grade is $5\ \mathrm{g/t}$ Au. The cut-off grade is $2\ \mathrm{g/t}$.
\begin{enumerate}
\item[a)] Calculate the tonnage of ore in the vein.
\item[b)] Calculate the mass of contained gold.
\item[c)] Is the vein above cut-off grade? Comment.
\item[d)] With a recovery of $90\ \%$ and a gold price of $60\,000\ \mathrm{FCFA/g}$, calculate the gross value of the recoverable gold.
\end{enumerate}
\tcblower
\textbf{a)} Volume of the vein (tabular approximation):
\[
V = \text{length} \times \text{width} \times \text{depth} = 200 \times 2 \times 50 = 20\,000\ \mathrm{m^3}
\]
Tonnage:
\[
T = V \times \rho = 20\,000\ \mathrm{m^3} \times 2.7\ \mathrm{t/m^3} = 54\,000\ \mathrm{t}
\]

\textbf{b)} Contained gold (grade in g/t $\times$ tonnage in t gives grams):
\[
M_{\mathrm{Au}} = 54\,000\ \mathrm{t} \times 5\ \mathrm{g/t} = 270\,000\ \mathrm{g} = 270\ \mathrm{kg\ of\ gold}
\]

\textbf{c)} The average grade ($5\ \mathrm{g/t}$) is \emph{above} the cut-off grade ($2\ \mathrm{g/t}$): $5 > 2$. The whole vein therefore qualifies as ore and can be counted in the reserve, provided mining costs are covered.

\textbf{d)} Recoverable gold (metallurgical recovery $= 90\%$):
\[
M_{\mathrm{rec}} = 270\,000\ \mathrm{g} \times 0.90 = 243\,000\ \mathrm{g}
\]
Gross value:
\[
V_{\mathrm{gross}} = 243\,000\ \mathrm{g} \times 60\,000\ \mathrm{FCFA/g} = 1.458 \times 10^{10}\ \mathrm{FCFA} \approx 14.6\ \text{billion FCFA}
\]
\textbf{Conclusion:} the vein contains $270\ \mathrm{kg}$ Au, of which $243\ \mathrm{kg}$ is recoverable, with a gross value of about $14.6$ billion FCFA. Whether it is a \emph{reserve} still depends on mining and processing costs being lower than this value.
\end{exoresolu}

\begin{exercice}[Tonnage and grade practice]
A stratiform zinc-lead deposit in the Benue Trough has a horizontal extent of $1.2\ \mathrm{km} \times 0.8\ \mathrm{km}$ and an average thickness of $4\ \mathrm{m}$. The bulk density is $3.2\ \mathrm{t/m^3}$. The average grade is $6\%$ Zn and $2\%$ Pb. The cut-off grade is $3\%$ Zn equivalent.
\begin{enumerate}
\item Calculate the tonnage of the deposit.
\item Calculate the contained zinc and lead (in tonnes).
\item If the metallurgical recovery is $85\%$ for Zn and $80\%$ for Pb, and prices are $2\,500\ \mathrm{FCFA/kg}$ Zn and $2\,000\ \mathrm{FCFA/kg}$ Pb, estimate the gross in-situ value of the recoverable metals.
\end{enumerate}
\end{exercice}

\courssub{Skill 4: Classifying a mineral deposit from its description}

\begin{definition}[Morphological and genetic classification]
Mineral deposits are classified by \textbf{morphology} (form, geometry, spatial relationship to host rocks) and by \textbf{genesis} (the geological process that concentrated the metal). Examiners expect candidates to state \emph{both} classifications and justify each with observed features.
\end{definition}

\begin{methode}[Using the classification schemes]
\begin{enumerate}
\item \textbf{Extract the diagnostic facts} from the statement: host rock, geometry of the body, mineral assemblage, associated igneous/sedimentary/metamorphic features.
\item \textbf{Apply the morphological (form) classification}: is the body \emph{stratiform/stratabound} (layer-parallel), a \emph{vein/lode} (cross-cutting fissure filling), a \emph{stockwork} (network of veinlets), a \emph{dissemination}, a \emph{pipe}, or a \emph{placer} (surficial accumulation)?
\item \textbf{Apply the genetic classification}: which process concentrated the metal?
\begin{itemize}
\item \textbf{Magmatic}: crystallisation, segregation, immiscible liquids (e.g.\ Cr, Ni, Pt in ultramafics; Sn, W, rare metals in pegmatites).
\item \textbf{Hydrothermal}: hot aqueous fluids depositing veins, replacements, stockworks (e.g.\ Au, Cu, Zn, Pb, Sn, W).
\item \textbf{Sedimentary}: placers (mechanical), chemical precipitates (BIF, evaporites), residual (laterites).
\item \textbf{Weathering/Residual}: lateritic bauxite, nickel laterite, supergene enrichment.
\item \textbf{Metamorphic}: skarns, recrystallised deposits, metamorphosed massive sulphides.
\end{itemize}
\item \textbf{Cross-check with the metal association}: e.g.\ Cr--Ni--Pt point to ultramafic magmatic segregation; Sn--W--Mo point to granitic hydrothermal systems; Fe--Al laterites point to tropical weathering.
\item \textbf{State both classifications} in your answer: examiners reward candidates who give form \emph{and} genesis.
\end{enumerate}
\end{methode}

\begin{aretenir}[Common morphological terms]
\begin{itemize}
\item \textbf{Stratiform}: layer-like, parallel to bedding/banding.
\item \textbf{Stratabound}: confined to a specific stratigraphic unit but not necessarily layered.
\item \textbf{Vein / Lode}: tabular, discordant, fracture-filling.
\item \textbf{Stockwork}: network of small veinlets.
\item \textbf{Disseminated}: fine-grained ore minerals scattered through the rock.
\item \textbf{Pipe / Chimney}: vertical, cylindrical body.
\item \textbf{Placer}: unconsolidated, surficial, hydraulically sorted.
\end{itemize}
\end{aretenir}

\begin{exoresolu}[Classifying four deposits]
Classify each of the following deposits \emph{morphologically} and \emph{genetically}, justifying each answer:
\begin{enumerate}
\item[a)] Layers of chromite interbedded with peridotite at the base of a large ultramafic intrusion.
\item[b)] Gold-bearing quartz veins cutting across granitic rocks at B\'etar\'e Oya.
\item[c)] The bauxite capping the lateritic plateaux of Minim-Martap.
\item[d)] Cassiterite and gold grains concentrated in the gravels of a river bed in the Kadey valley.
\end{enumerate}
\tcblower
\textbf{a)} \emph{Morphology:} stratiform (the chromite forms layers parallel to the igneous layering). \emph{Genesis:} \textbf{magmatic segregation} — early-crystallising dense chromite crystals settled and accumulated at the base of the crystallising ultramafic magma (gravitational settling during fractional crystallisation). Typical association: Cr with ultramafic rocks.

\textbf{b)} \emph{Morphology:} veins/lodes (cross-cutting fissure fillings). \emph{Genesis:} \textbf{hydrothermal} — hot aqueous fluids released during the late stages of granitic magmatism circulated through fractures; cooling and pressure drop caused quartz and gold to precipitate in the fissures. This is a mesothermal/epithermal-type Au--quartz vein system.

\textbf{c)} \emph{Morphology:} blanket/stratiform residual cap (lateritic profile). \emph{Genesis:} \textbf{residual weathering (lateritisation)} — intense tropical weathering of aluminous parent rocks leached silica and alkalis, concentrating insoluble aluminium hydroxides (gibbsite, boehmite) near the surface. This is a lateritic bauxite deposit.

\textbf{d)} \emph{Morphology:} placer (surficial, unconsolidated accumulation in river gravels). \emph{Genesis:} \textbf{sedimentary/mechanical concentration (alluvial placer)} — weathering liberated dense, chemically resistant minerals (cassiterite, gold) from primary sources; running water sorted and concentrated them by density in traps of the river bed.

\textbf{Synthesis:} the four deposits illustrate the four great genetic families — magmatic, hydrothermal, residual, and sedimentary (placer) — and show that the same metal (e.g.\ cassiterite) can occur both as a primary hydrothermal deposit and as a secondary placer.
\end{exoresolu}

\courssub{Skill 5: Explaining ore-forming processes (magmatic and hydrothermal)}

\begin{methode}[Building a genetic explanation for magmatic/hydrothermal deposits]
\begin{enumerate}
\item \textbf{State the source} of the metal (magma, hydrothermal fluid, country rock leached by fluids).
\item \textbf{State the transport medium} (silicate melt, immiscible sulphide liquid, hot aqueous fluid).
\item \textbf{State the concentration mechanism}:
\begin{itemize}
\item \textbf{Magmatic}: crystal settling (Cr, Pt), immiscible sulphide liquid segregation (Ni--Cu), late-stage residual melt enrichment (rare metals in pegmatites, greisens).
\item \textbf{Hydrothermal}: dissolution at depth, transport as complexes in hot aqueous fluids, then precipitation caused by cooling, pressure drop, boiling, fluid mixing or reaction with wall rock.
\end{itemize}
\item \textbf{State the trap/host}: layered intrusion, fracture/vein, porous or reactive (e.g.\ carbonate) wall rock, contact zone (skarn).
\item \textbf{Link each observation to the process} in a cause--effect chain; examiners penalise bare lists of facts.
\end{enumerate}
\textbf{Key ideas:} hydrothermal fluids are mostly magmatic or heated meteoric water; deposits are zoned (high-temperature minerals near the source, low-temperature ones farther away); porphyry and skarn deposits form around intrusive contacts.
\end{methode}

\begin{experience}[Magmatic sulphide segregation — a physical model]
Fill a transparent container with water (silicate melt analogue) and add a dense, immiscible liquid (e.g.\ glycerol with food colouring, or simply oil) to represent sulphide liquid. Shake to form droplets; observe them coalesce and sink to the bottom. This demonstrates gravitational accumulation of an immiscible dense liquid — the key step in magmatic Ni--Cu--PGE deposit formation.
\end{experience}

\begin{exoresolu}[From magma to a magmatic Ni--Cu sulphide deposit]
Describe and explain the sequence of processes by which a nickel--copper sulphide deposit can form from a mafic magma. Illustrate your answer with a labelled cross-section.
\tcblower
\textbf{Step 1 — Source and melting.} A mantle-derived mafic/ultramafic magma rises into the crust. It carries Fe, Ni, Cu and sulphur dissolved in the silicate melt.

\textbf{Step 2 — Sulphur saturation and immiscibility.} As the magma cools (or assimilates sulphur from crustal rocks), it becomes saturated in sulphur. A dense \emph{immiscible sulphide liquid} separates from the silicate melt, like oil separating from water. Because Ni, Cu and the platinum-group elements have a strong chemical affinity for sulphide (they are \emph{chalcophile}), they partition preferentially into this sulphide liquid, scavenging the metals from a large volume of magma (partition coefficients $D_{\mathrm{Ni}}^{\mathrm{sul/sil}} \sim 100$--$1000$).

\textbf{Step 3 — Gravitational accumulation.} The dense sulphide droplets coalesce and sink, collecting at the base of the magma chamber or in structural traps (embayments, feeder conduits, troughs).

\textbf{Step 4 — Crystallisation.} On final cooling the sulphide liquid crystallises as pyrrhotite (\ce{FeS}), pentlandite ((\ce{Fe,Ni})$_{9}$\ce{S8}) and chalcopyrite (\ce{CuFeS2}), forming a massive or disseminated sulphide ore body at the base of the intrusion.

\begin{popfigure}
\begin{center}
\begin{tikzpicture}[scale=0.9]
\fill[poporange!20] (0,0) rectangle (8,3.5);
\fill[popdark!70] (1,0) -- (7,0) -- (6.2,0.7) -- (1.8,0.7) -- cycle;
\draw[thick] (0,0) rectangle (8,3.5);
\foreach \x/\y in {2.5/2.8,4/2.2,5.5/3.0,3.2/1.8,6/2.0,2.0/1.2,4.5/1.5,6.5/1.3} {\fill[popdark] (\x,\y) circle (0.05);}
\draw[->,thick,popdark] (4,3.2) -- (4,1.2);
\node[align=center, text width=2.5cm] at (7.0,2.8) {\small mafic magma\\[-2pt]\small (silicate melt)};
\node[align=center, text width=2cm] at (1.5,2.0) {\small sulphide\\[-2pt]\small droplets};
\node[align=center, text width=3cm] at (4,-0.5) {\small massive Ni--Cu sulphide ore\\[-2pt]\small (basal accumulation)};
\draw[->,thick] (4,-0.4) -- (4,0.2);
\end{tikzpicture}
\legende{Magmatic segregation: dense immiscible sulphide droplets sink and accumulate at the base of the magma chamber.}
\end{center}
\end{popfigure}

\textbf{Conclusion:} the deposit results from (i) sulphur saturation, (ii) immiscible sulphide liquid segregation which scavenges chalcophile metals, and (iii) gravitational settling — a classic \emph{magmatic segregation} process.
\end{exoresolu}

\begin{savaistu}[Cameroon context: magmatic deposits]
The ultramafic complexes of the Pan-African belt in southern Cameroon (e.g.\ Nkamouna, Lomi\'e) host Ni--Co--Mn laterites developed on serpentinised peridotites. The primary magmatic sulphide mineralisation (if any) has been largely overprinted by tropical weathering, but the ultramafic protolith is the essential source. The Cameroon Volcanic Line itself is not a major magmatic sulphide province, but its alkaline rocks host rare-metal mineralisation (e.g.\ Nb, REE at Lomi\'e, Akonolinga).
\end{savaistu}

\courssub{Skill 6: Explaining sedimentary, residual and supergene enrichment processes}

\begin{methode}[Explaining surface and near-surface ore formation]
\begin{enumerate}
\item \textbf{Identify the climate and landscape context}: tropical humid climate $\Rightarrow$ intense chemical weathering, laterites, bauxite, nickel laterite; arid basins $\Rightarrow$ evaporites; rivers and coasts $\Rightarrow$ placers.
\item \textbf{For residual deposits}: describe leaching of soluble components (\ce{SiO2}, Ca, Na, K, Mg) and residual concentration of insoluble ones (Al, Fe, Ti, Ni, Cr) in the weathering profile. Mention the lateritic profile: ferricrete / laterite duricrust $\rightarrow$ mottled zone $\rightarrow$ pallid zone $\rightarrow$ saprolite $\rightarrow$ bedrock.
\item \textbf{For placers}: apply the three conditions — (i) a primary source, (ii) liberation by weathering, (iii) hydraulic sorting concentrating \emph{dense and resistant} minerals in traps (bedrock irregularities, inside bends, base of gravels, riffles).
\item \textbf{For supergene enrichment}: describe the vertical zonation — leached cap (gossan) at the top, \emph{supergene enriched zone} below the water table where descending copper re-precipitates as chalcocite (\ce{Cu2S}), covellite (\ce{CuS}), native copper; primary (hypogene) ore at depth (chalcopyrite, bornite). Explain that enrichment can upgrade a sub-economic protore into an ore.
\item \textbf{For chemical sedimentary ores}: evaporation (evaporites: halite, gypsum, potash) or chemical/biochemical precipitation (banded iron formations — BIF).
\end{enumerate}
\end{methode}

\begin{attention}[Placer vs residual — don't confuse them]
\begin{itemize}
\item \textbf{Residual (lateritic)}: forms \emph{in situ} by chemical weathering; the ore stays where the parent rock was. No transport.
\item \textbf{Placer}: forms by \emph{mechanical} transport and hydraulic sorting; the ore moves from its source. Requires density contrast and durability.
\item Supergene enrichment occurs \emph{below} the water table in the saturated zone; the leached cap is \emph{above} it.
\end{itemize}
\end{attention}

\begin{exoresolu}[Why is a placer rich only in certain minerals?]
Alluvial workings in the East Region of Cameroon recover gold and cassiterite, but never feldspar or calcite, even though these are abundant in the source rocks.
\begin{enumerate}
\item[a)] State the two physical--chemical properties a mineral must possess to form a placer deposit, and name the property that controls hydraulic sorting.
\item[b)] Explain why feldspar and calcite are absent from the placers.
\item[c)] Draw a labelled sketch of a river section showing where placer gold accumulates.
\end{enumerate}
\tcblower
\textbf{a)} A placer mineral must be (i) \emph{dense} (high specific gravity) so that running water can separate it from lighter grains, and (ii) \emph{chemically resistant and mechanically durable} so that it survives weathering and transport. Hydraulic sorting is controlled essentially by \textbf{density} (together with grain size and shape): when water velocity drops, the densest grains settle first.

\textbf{b)} Feldspar and calcite fail both tests: they have low densities (about $2.6$--$2.7\ \mathrm{g/cm^3}$, close to quartz), so they are not separated from ordinary sand; and they are chemically unstable — feldspar alters to clay minerals (kaolinite) during tropical weathering, and calcite simply dissolves in slightly acidic surface waters. They are therefore destroyed or dispersed long before they could be concentrated. Gold ($19.3\ \mathrm{g/cm^3}$, inert) and cassiterite ($\approx 7.0\ \mathrm{g/cm^3}$, resistant) survive and are concentrated.

\textbf{c)}
\begin{popfigure}
\begin{center}
\begin{tikzpicture}[scale=1.0]
\draw[thick, popdark] (0,0) .. controls (1,0.4) and (1.6,-0.6) .. (2.6,-0.5) .. controls (3.6,-0.4) and (4.2,0.3) .. (5.2,0) .. controls (6.2,-0.3) and (6.8,0.2) .. (8,0);
\draw[popblue,thick] (0,0.55) .. controls (2,0.75) and (6,0.4) .. (8,0.6);
\fill[popgold!80] (1.9,-0.55) ellipse (0.55 and 0.16);
\fill[popgold!80] (5.0,-0.12) ellipse (0.45 and 0.12);
\node[align=center] at (1.9,-1.1) {\small Au trap:\\[-2pt]\small pothole in bedrock};
\node[align=center] at (5.0,-0.8) {\small Au at base\\[-2pt]\small of gravel bar};
\node at (6.9,0.85) {\small river flow};
\draw[->,popblue,thick] (6.1,0.85) -- (6.6,0.85);
\node[align=center] at (0.9,-0.5) {\small bedrock};
\draw[popline, dashed] (0,-0.5) -- (8,-0.5);
\node at (4,-0.9) {\small base of alluvial gravels};
\end{tikzpicture}
\legende{Placer gold accumulates where current velocity drops: bedrock potholes, irregularities and the base of gravel bars.}
\end{center}
\end{popfigure}

\textbf{Conclusion:} placers are natural density-concentration plants operated by running water; only dense, resistant minerals such as gold, cassiterite, diamond or platinum can form them.
\end{exoresolu}

\begin{experience}[Supergene enrichment profile — sketch and label]
Draw a vertical cross-section through a porphyry copper deposit showing: (1) leached cap (gossan) with iron oxides, (2) supergene enrichment zone (chalcocite blanket) at/near water table, (3) primary hypogene zone (chalcopyrite + pyrite) at depth. Label the water table and indicate the direction of metal migration (downward for Cu, upward for Fe). Explain why the chalcocite zone is richer than the primary ore.
\end{experience}

\courssub{Skill 7: Interpreting a geological cross-section to locate ore bodies}

\begin{methode}[Reading occurrence geometry on a map or section]
\begin{enumerate}
\item \textbf{Identify the host structures}: fractures and faults (veins), intrusive contacts (skarns, greisens), stratigraphic layers (stratiform ores), unconformities, palaeo-river channels (placers).
\item \textbf{Note the spatial relationship of ore to geology}: concordant (parallel to bedding $\Rightarrow$ syngenetic, formed with the rock) or discordant (cross-cutting $\Rightarrow$ epigenetic, formed later).
\item \textbf{Use metal zoning}: high-temperature metals (Sn, W, Mo, Bi) near the intrusion; intermediate (Cu, Zn, Pb, Fe) farther out; low-temperature (Ag, Hg, Sb, Au) at the periphery. Predict where each metal should be sought.
\item \textbf{Reason on depth and continuity}: a vein pinches and swells; a stratiform body is laterally continuous; a placer follows the channel.
\item \textbf{Conclude with an exploration recommendation}: where to drill or dig next, with justification.
\end{enumerate}
\end{methode}

\begin{definition}[Syngenetic vs Epigenetic]
A \cle{syngenetic} deposit forms \emph{at the same time} as its host rock (e.g.\ a sedimentary iron bed, a magmatic chromite layer); it is usually concordant. An \cle{epigenetic} deposit forms \emph{later} than its host rock, introduced into pre-existing rocks (veins, replacements, skarns, placers); it is usually discordant.
\end{definition}

\begin{exoresolu}[Where should the exploration company drill?]
The cross-section below shows a granite intruding schists and limestones. A tin-bearing greisen occurs at the granite roof (point A); skarn with copper occurs at the limestone contact (point B); quartz veins with silver occur in fractures far from the contact (point C).
\begin{enumerate}
\item[a)] Define \emph{epigenetic} and \emph{syngenetic} deposits and state which applies here.
\item[b)] Explain the observed metal zoning Sn $\rightarrow$ Cu $\rightarrow$ Ag away from the granite.
\item[c)] The company wants more tin. Where should it explore, and why?
\end{enumerate}
\tcblower
\textbf{a)} A \cle{syngenetic} deposit forms \emph{at the same time} as its host rock (e.g.\ a sedimentary iron bed); it is usually concordant. An \cle{epigenetic} deposit forms \emph{later} than its host rock, introduced into pre-existing rocks (veins, replacements); it is usually discordant. Here all three occurrences (greisen, skarn, veins) were emplaced by fluids related to the granite \emph{after} the schists and limestones existed: they are \textbf{epigenetic}.

\textbf{b)} The hydrothermal fluids were hottest at or within the granite and cooled as they migrated outward. Minerals precipitate in order of decreasing temperature of deposition:
\begin{itemize}
\item \textbf{Sn (cassiterite)} with tourmaline/topaz in the greisen: highest temperature ($400$--$500^\circ\mathrm{C}$), deposited at the granite roof (A);
\item \textbf{Cu (and Fe)} in the skarn: high-to-intermediate temperature ($300$--$400^\circ\mathrm{C}$), deposited where the fluids reacted with the reactive limestone at the contact (B);
\item \textbf{Ag} in distal quartz veins: lowest temperature ($<250^\circ\mathrm{C}$), deposited far from the heat source (C).
\end{itemize}
This is the classic \emph{temperature-controlled hydrothermal zoning} around an intrusion.

\textbf{c)} Tin is a high-temperature metal, so the company should explore \emph{close to and within the granite}, especially along the roof and cupolas of the intrusion and along contacts with the schists (greisen-bordered veins), i.e.\ around point A and at depth beneath it — \emph{not} in the distal fractures, which are the low-temperature silver zone. Drilling should target the granite margins and any unexposed apophyses (satellite intrusions) where greisenisation may have occurred.
\end{exoresolu}

\begin{popfigure}
\begin{center}
\begin{tikzpicture}[scale=0.85]
% Granite
\fill[poporange!30] (0,3) -- (8,3) -- (8,5) -- (0,5) -- cycle;
\draw[thick] (0,3) -- (8,3);
\draw[thick] (0,3) -- (0,5) -- (8,5) -- (8,3);
\node[align=center] at (4,4) {\textbf{GRANITE}};
% Schists
\fill[popblue!15] (0,0) -- (8,0) -- (8,3) -- (0,3) -- cycle;
\foreach \x in {0.5,1.5,...,7.5} \draw[popblue!60] (\x,0) -- (\x+0.3,3);
\node[align=center] at (4,1.5) {\textbf{SCHISTS}};
% Limestone
\fill[popgreen!15] (8,0) -- (12,0) -- (12,3) -- (8,3) -- cycle;
\draw[popgreen!60] (8,0) -- (8,3);
\foreach \y in {0.5,1.0,...,2.5} \draw[popgreen!60] (8,\y) -- (12,\y);
\node[align=center] at (10,1.5) {\textbf{LIMESTONE}};
% Greisen at A
\fill[poppink!50] (3.5,2.8) rectangle (4.5,3.2);
\node[poppink, align=center] at (4,3.4) {\textbf{A: Greisen (Sn)}};
% Skarn at B
\fill[poporange!50] (7.8,1.5) rectangle (8.2,2.5);
\node[poporange, align=center] at (8,2.8) {\textbf{B: Skarn (Cu)}};
% Veins at C
\draw[thick, popdark] (10,1.0) -- (10.5,2.5);
\draw[thick, popdark] (10.2,0.8) -- (10.7,2.3);
\node[popdark, align=center] at (10.5,2.9) {\textbf{C: Veins (Ag)}};
% Arrows for fluid flow
\draw[->, thick, popblue] (4,4) -- (4,3.2);
\draw[->, thick, popblue] (4,3.2) -- (8,2.0);
\draw[->, thick, popblue] (8,2.0) -- (10.5,1.5);
\node[popblue] at (2,3.8) {Hot fluids};
\node[popblue] at (6,2.5) {Cooling};
\node[popblue] at (9.5,1.0) {Cool};
% Water table
\draw[dashed, popblue] (0,0.5) -- (12,0.5);
\node[popblue] at (11,0.3) {Water table};
\end{tikzpicture}
\legende{Hydrothermal zoning around a granitic intrusion: Sn (greisen) at roof $\rightarrow$ Cu (skarn) at carbonate contact $\rightarrow$ Ag (veins) distally.}
\end{center}
\end{popfigure}

\courssub{Skill 8: Writing a synoptic essay on Cameroon's mineral occurrences}

\begin{methode}[Structuring a GCE essay on mineral occurrence]
\begin{enumerate}
\item \textbf{Define the key terms} in one or two sentences (economic mineral, deposit, occurrence, resource, reserve).
\item \textbf{Organise by genetic family or by commodity}, never as a random list. A strong plan: magmatic/hydrothermal metals $\rightarrow$ residual deposits $\rightarrow$ sedimentary/placer deposits $\rightarrow$ industrial minerals and hydrocarbons.
\item \textbf{Attach each commodity to a place and a process}: e.g.\ bauxite -- Minim-Martap/Ngaoundal -- lateritisation; iron -- Mbalam -- BIF; gold -- B\'etar\'e Oya -- hydrothermal veins and placers; tin -- Mayo Darl\'e -- granitic hydrothermal; rutile -- Akonolinga; limestone -- Figuil (cement); petroleum -- Logone Birni and Rio del Rey basins.
\item \textbf{Link occurrence to the geological framework}: Congo craton (Fe, Au), Pan-African belts (Au, Sn, rutile), Cameroon Volcanic Line (volcanic rocks, industrial minerals, some REE), sedimentary basins (oil, gas, salt, limestone).
\item \textbf{Conclude} with the economic significance and the resource/reserve caution (many Cameroonian occurrences are resources awaiting infrastructure, investment or favourable prices).
\end{enumerate}
\end{methode}

\begin{aretenir}[Cameroon's geological domains and their mineral signatures]
\begin{center}
\begin{tabularx}{\linewidth}{|l|X|X|}
\hline
\rowcolor{popblueL}
\textbf{Domain} & \textbf{Key commodities} & \textbf{Genetic process / Setting} \\
\hline
Congo Craton (South) & Iron (Mbalam, Nkout), Gold (primary \& placer) & BIF (chemical sedimentary), shear-zone hydrothermal Au \\
Pan-African Belts (Centre, East, Adamawa) & Gold (B\'etar\'e Oya, Colomine, Batouri), Tin (Mayo Darl\'e), Rutile (Akonolinga), Ni--Co laterites (Lomi\'e, Nkamouna) & Hydrothermal veins (Au, Sn), granitic greisens (Sn), lateritic weathering on ultramafics (Ni--Co), placer reworking \\
Cameroon Volcanic Line & Pozzolana, basalt, scoria, some REE/Nb (Lomi\'e, Akonolinga) & Alkaline volcanism, weathering of volcanic rocks \\
Sedimentary Basins (Logone Birni, Rio del Rey, Douala/Kribi-Campo) & Petroleum, natural gas, limestone (Figuil, Mfou), salt, glass sand & Sedimentary (organic/inorganic), evaporites \\
\hline
\end{tabularx}
\end{center}
\end{aretenir}

\begin{exoresolu}[Guided essay: the mineral endowment of Cameroon]
``Cameroon's geology gives it a varied but unevenly exploited mineral endowment.'' Discuss, with reference to the mode of occurrence and genesis of at least four mineral commodities.
\tcblower
\textbf{Introduction.} An economic mineral is one that can be extracted at a profit; its occurrence is controlled by geology, its status by economics. Cameroon's territory spans three great geological domains — the Congo craton in the south, the Pan-African mobile belts in the centre and north, and the Cameroon Volcanic Line crossing both — plus coastal and interior sedimentary basins. Each domain hosts a characteristic family of deposits.

\textbf{1. Residual deposits of the cratonic south and Adamawa.} Intense tropical weathering of aluminous rocks produced the great lateritic \textbf{bauxite} plateaux of Minim-Martap and Ngaoundal (Adamawa), among the largest known bauxite resources in Africa (several billion tonnes), and nickel--cobalt laterites near Lomi\'e and Nkamoula. Genesis: leaching of silica and alkalis, residual enrichment of Al, Fe, Ni in a lateritic profile (duricrust $\rightarrow$ mottled $\rightarrow$ pallid $\rightarrow$ saprolite). Status: still essentially \emph{resources}, for lack of a railway to the coast and moderate alumina grades.

\textbf{2. Banded iron formations of the craton.} The \textbf{iron} deposits of Mbalam and Nkout (South Region) are itabirite-type banded iron formations (BIF), chemical sediments of Palaeoproterozoic age, locally upgraded by metamorphism and supergene enrichment. Genesis: chemical/biochemical precipitation of iron in ancient oceans, later metamorphic recrystallisation and tropical weathering leaching silica. Status: Mbalam is a \emph{reserve} (feasibility completed, rail-port infrastructure planned); Nkout remains a \emph{resource}.

\textbf{3. Hydrothermal and placer metals of the Pan-African belts.} \textbf{Gold} occurs in quartz veins at B\'etar\'e Oya, Colomine and Batouri (East Region), deposited from hydrothermal fluids in Pan-African shear zones, and in derived alluvial placers worked artisanally. \textbf{Tin (cassiterite)} at Mayo Darl\'e is associated with granitic magmatism (greisens and veins) and its placers. \textbf{Rutile} at Akonolinga occurs in hydrothermal veins and pegmatites related to Pan-African granites. Genesis: magmatic-hydrothermal concentration in shear zones and granitic cupolas, then mechanical re-concentration by rivers.

\textbf{4. Sedimentary and industrial commodities.} The Logone Birni (Chad Basin) and Rio del Rey / Douala / Kribi-Campo basins contain \textbf{petroleum and natural gas} trapped in structural and stratigraphic traps; \textbf{limestone} at Figuil (North) and Mfou (Centre) supports the cement industry; \textbf{pozzolana} and basalt from the Cameroon Volcanic Line are quarried for construction aggregates and cement additives.

\textbf{Conclusion.} Cameroon's endowment is indeed varied — magmatic-hydrothermal metals (Au, Sn, W, rutile), residual bauxite and nickel laterites, sedimentary iron (BIF), hydrocarbons and industrial minerals — but unevenly exploited: geology has provided the occurrences, while economics (grade, infrastructure, prices, policy) decides which are reserves. The bauxite of Adamawa, a giant resource awaiting a railway, is the perfect illustration that an ``economic mineral'' is defined as much by economics as by geology.
\end{exoresolu}

\begin{exercice}[Essay practice]
``The classification of a mineral occurrence as a resource or a reserve depends on factors beyond geology.'' Discuss this statement with reference to at least two Cameroonian mineral occurrences.
\end{exercice}

\begin{corrige}[Essay practice]
\textbf{Key points to include:}
\begin{itemize}
\item Definitions: resource (geological confidence) vs reserve (economic viability now).
\item Example 1: Minim-Martap bauxite — huge indicated/measured resource, but no reserve because transport costs > value; needs railway, price rise, or local refinery.
\item Example 2: Mbalam iron — moved from resource to reserve after feasibility study, rail-port agreement (Kribi deep-water port), and offtake contracts.
\item Example 3: B\'etar\'e Oya gold — artisanal exploitation proves technical feasibility; reserve status depends on mechanised mining costs vs gold price.
\item Conclusion: geology sets the stage (grade, tonnage, mineralogy), but economics, technology, infrastructure, law and politics draw the resource/reserve line.
\end{itemize}
\end{corrige}

\newpage

\coursec{Exercises}

\courssub{I check my knowledge}

\begin{exercice}[MCQ: economic minerals and deposits]
For each question, choose the single correct answer and justify your choice in one sentence.
\begin{enumerate}
\item An \cle{ore} is:
\begin{enumerate}
\item any rock containing a metallic element;
\item a rock from which a useful substance can be extracted \textbf{at a profit};
\item the waste material surrounding a valuable mineral;
\item a mineral deposit that has not yet been discovered.
\end{enumerate}
\item The \cle{grade} of a deposit is:
\begin{enumerate}
\item the total tonnage of rock in the deposit;
\item the concentration of the valuable substance in the ore;
\item the depth of the deposit below the surface;
\item the market price of the metal.
\end{enumerate}
\item Which of the following is a \textbf{syngenetic} deposit?
\begin{enumerate}
\item a gold-bearing quartz vein cutting granite;
\item a bedded ironstone formed at the same time as the enclosing sediments;
\item a skarn developed at a granite--limestone contact;
\item a hydrothermal copper vein.
\end{enumerate}
\item The Minim--Martap deposit of the Adamawa Plateau is mainly a deposit of:
\begin{enumerate}
\item gold; \item iron; \item bauxite; \item limestone.
\end{enumerate}
\end{enumerate}
\end{exercice}

\begin{exercice}[True or false?]
State whether each assertion is \textbf{true} or \textbf{false}, and justify your answer in one or two sentences.
\begin{enumerate}
\item A mineral resource automatically becomes a reserve as soon as it is discovered.
\item A rise in the world price of a metal can convert a sub-economic resource into a reserve without any new discovery.
\item Hydrothermal deposits are always older than the rocks that enclose them.
\item Residual (lateritic) concentration of aluminium is favoured by a hot, humid tropical climate such as that of the Adamawa Plateau.
\end{enumerate}
\end{exercice}

\begin{exercice}[Key definitions]
Define each of the following terms in one or two precise sentences, giving one example where appropriate:
\begin{enumerate}
\item \cle{economic mineral};
\item \cle{gangue};
\item \cle{cut-off grade};
\item \cle{epigenetic deposit};
\item \cle{placer deposit}.
\end{enumerate}
\end{exercice}

\begin{exercice}[Direct calculation: contained metal]
A copper deposit contains $5$ million tonnes of ore grading $0.8\ \%$ Cu.
\begin{enumerate}
\item Calculate the mass of copper metal contained in the deposit.
\item The mine produces $12\,000$ tonnes of copper metal per year. Calculate the life of the reserve, assuming total recovery.
\item State one reason why the real mine life may be shorter than your calculated value.
\end{enumerate}
\end{exercice}

\courssub{I practise}

\begin{exercice}[Resource or reserve?]
The table below summarises exploration results for three deposits A, B and C.

\begin{center}
\begin{tabularx}{\linewidth}{|l|X|X|X|}
\hline
\rowcolor{popblueL} \textbf{Deposit} & \textbf{Geological knowledge} & \textbf{Economic study} & \textbf{Estimated tonnage} \\
\hline
A & Few boreholes; continuity assumed between widely spaced data points & No feasibility study & $50$ Mt \\
\hline
B & Dense drilling grid; continuity well established & Positive feasibility study at current prices & $120$ Mt \\
\hline
C & Moderate drilling; continuity reasonably assumed & Preliminary study only & $80$ Mt \\
\hline
\end{tabularx}
\end{center}

\begin{enumerate}
\item Classify each deposit as an \cle{inferred resource}, an \cle{indicated resource} or a \cle{proved reserve}, justifying each classification.
\item Explain what additional work would be needed to upgrade deposit C to reserve status.
\end{enumerate}
\end{exercice}

\begin{exercice}[Interpreting a weathering profile]
The diagram below is a schematic lateritic weathering profile developed on syenite at Minim--Martap (Adamawa, Cameroon).

\begin{popfigure}
\begin{center}
\begin{tikzpicture}[scale=1]
\fill[popgreenL] (0,3.2) rectangle (7,4.2);
\fill[popgold!60] (0,2.0) rectangle (7,3.2);
\fill[poporange!50] (0,0.8) rectangle (7,2.0);
\fill[gray!40] (0,0) rectangle (7,0.8);
\draw[thick] (0,0) rectangle (7,4.2);
\node at (3.5,3.7) {(a)};
\node at (3.5,2.6) {(b)};
\node at (3.5,1.4) {(c)};
\node at (3.5,0.4) {(d)};
\draw[->,thick,popblue] (7.4,0.4) -- (7.4,3.9) node[midway,right,align=left,text width=3.2cm] {increasing weathering upwards};
\end{tikzpicture}
\end{center}
\legende{Schematic lateritic profile, Minim--Martap.}
\end{popfigure}

\begin{enumerate}
\item Name the horizons labelled (a) to (d), from the surface downwards.
\item In which horizon is bauxite concentrated? Name the main aluminium minerals it contains.
\item Explain, in three or four sentences, how intense chemical weathering removes silica and alkalis while concentrating aluminium.
\item State two features of Cameroon's climate that favour this process.
\end{enumerate}
\end{exercice}

\begin{exercice}[Cross-cutting relations and genetic classification]
A cross-section shows a granite intruding schists. Quartz vein V1, carrying chalcopyrite, cuts the schists but is itself cut and offset by quartz vein V2, which carries gold. Both veins stop against a younger dolerite dyke.
\begin{enumerate}
\item Draw a simple labelled sketch cross-section of these relationships.
\item Arrange the granite, schists, veins V1 and V2, and the dyke in order from oldest to youngest, justifying each step with the principle of cross-cutting relationships.
\item Classify the copper and gold mineralisations as \cle{syngenetic} or \cle{epigenetic}, and as \cle{magmatic}, \cle{hydrothermal}, \cle{sedimentary} or \cle{residual}, giving reasons.
\end{enumerate}
\end{exercice}

\begin{exercice}[From resource to reserve]
An iron deposit near Mbalam (East Region) is currently classified as a sub-economic resource.
\begin{enumerate}
\item Distinguish clearly between a \cle{resource} and a \cle{reserve}.
\item Explain \textbf{two} distinct ways in which this resource could become a reserve without any new discovery (think of prices, technology, infrastructure).
\item Suggest one reason why a deposit may remain unexploited even when it is technically a reserve.
\end{enumerate}
\end{exercice}

\courssub{I prepare for the GCE}

\begin{exercice}[Structured question: ore, grade and reserves (20 marks)]
\begin{enumerate}
\item Define the following terms, giving one example of each: \emph{ore}, \emph{gangue}, \emph{grade}, \emph{cut-off grade}. \hfill (8 marks)
\item A copper deposit contains $2$ million tonnes of ore grading $1.2\ \%$ Cu. Calculate the mass of copper metal contained in the deposit. \hfill (3 marks)
\item The mine produces $20\,000$ tonnes of copper metal per year. Calculate the life of the reserve, stating any assumption you make. \hfill (3 marks)
\item Explain three factors that determine whether a given mineral concentration can be worked as an ore. \hfill (6 marks)
\end{enumerate}
\end{exercice}

\begin{exercice}[Map exercise: mineral occurrences of Cameroon (20 marks)]
You are provided with an outline map of Cameroon showing the regions and the main towns.
\begin{enumerate}
\item On the map, place and label five major mineral occurrences: gold (East Region), bauxite (Minim--Martap), iron (Mbalam), nickel--cobalt (Lomie area), limestone (Figuil). \hfill (5 marks)
\item For each occurrence, state the genetic type of the deposit (magmatic, hydrothermal, sedimentary, residual/weathering or metamorphic) and justify your answer in one sentence. \hfill (10 marks)
\item Choose one region that you consider prospective for gold exploration. Justify your choice using two geological criteria (nature of the basement rocks, known occurrences, structural features). \hfill (5 marks)
\end{enumerate}
\end{exercice}

\begin{exercice}[Essay: geology controls mineral distribution (20 marks)]
\emph{``The distribution of mineral deposits is controlled by geology, not by politics.''}

Discuss this statement with reference to:
\begin{enumerate}
\item at least \textbf{three} Cameroonian deposits (for example the bauxite of Minim--Martap, the iron of Mbalam, the gold of the East Region, the nickel--cobalt of Lomie, the limestone of Figuil), explaining in each case the geological process that localised the mineralisation;
\item at least \textbf{one} world metallogenic province (for example the copper belt of Central Africa or the gold provinces of cratonic shields);
\item the extent to which political boundaries, infrastructure and economic decisions nevertheless influence whether a deposit is actually mined.
\end{enumerate}
Your essay should include an introduction, a developed argument with labelled diagrams where useful, and a conclusion. \hfill (20 marks)
\end{exercice}

\newpage

\coursec{Solutions to the exercises}

\begin{corrige}[Exercise 1 — MCQ: economic minerals and deposits]
\begin{enumerate}
\item \textbf{Answer: (b).} An ore is defined by \emph{economic} extractability: a rock is only an ore if the valuable substance it contains can be extracted and sold \textbf{at a profit} under current technical and economic conditions. A rock merely containing a metal (a) is not necessarily an ore; (c) defines gangue; (d) describes an undiscovered resource.
\item \textbf{Answer: (b).} The grade is the \textbf{concentration} of the valuable substance in the ore, expressed as a percentage (e.g.\ $\%$ Cu) or in $\mathrm{g/t}$ for precious metals. Tonnage (a), depth (c) and price (d) are separate parameters of a deposit.
\item \textbf{Answer: (b).} A syngenetic deposit forms \textbf{at the same time} as the enclosing rock: a bedded ironstone precipitated during sedimentation is contemporaneous with its host sediments. Veins (a, d) and skarns (c) are introduced \emph{after} the host rock and are therefore epigenetic.
\item \textbf{Answer: (c).} Minim--Martap on the Adamawa Plateau is Cameroon's major \textbf{bauxite} (aluminium ore) deposit, formed by intense lateritic weathering of nepheline syenite under a tropical climate.
\end{enumerate}
\end{corrige}

\begin{corrige}[Exercise 2 — True or false?]
\begin{enumerate}
\item \textbf{False.} A resource is a concentration whose existence and approximate size are known or inferred; it becomes a \textbf{reserve} only after detailed exploration (dense drilling) \emph{and} a feasibility study proving that extraction is technically possible and economically profitable. Discovery alone is not sufficient.
\item \textbf{True.} Reserves are defined partly by economics: if the metal price rises, the cut-off grade falls and previously sub-economic material becomes profitable to mine, so tonnages can be reclassified from resource to reserve with no new discovery.
\item \textbf{False.} Hydrothermal fluids circulate \emph{through} pre-existing rocks and deposit minerals in fractures and pores, so hydrothermal deposits are \textbf{epigenetic} — younger than their host rocks (principle of cross-cutting relationships).
\item \textbf{True.} Hot temperatures accelerate chemical reactions and abundant rainfall provides the percolating water that leaches silica and alkalis; this intense lateritic weathering concentrates residual aluminium hydroxides (bauxite), as at Minim--Martap on the Adamawa Plateau.
\end{enumerate}
\end{corrige}

\begin{corrige}[Exercise 3 — Key definitions]
\begin{enumerate}
\item \textbf{Economic mineral:} any naturally occurring mineral substance (metallic or non-metallic) whose concentration in the Earth's crust is sufficient for it to be extracted and used \textbf{at a profit}; e.g.\ bauxite at Minim--Martap, gold in the East Region.
\item \textbf{Gangue:} the valueless minerals or rock material intimately mixed with the ore minerals in a deposit, which must be separated and discarded during processing; e.g.\ quartz accompanying gold in a vein.
\item \textbf{Cut-off grade:} the minimum grade above which material is economically worth mining and processing; material below the cut-off grade is treated as waste. It depends on metal price, mining cost and recovery.
\item \textbf{Epigenetic deposit:} a deposit formed \emph{after} the enclosing rock, introduced into it by fluids or melts; e.g.\ a gold-bearing quartz vein cutting schists.
\item \textbf{Placer deposit:} a deposit formed by the mechanical concentration of dense, chemically resistant minerals (gold, diamond, cassiterite) in river gravels or beach sands after weathering, transport and sorting; e.g.\ alluvial gold in the rivers of the East Region of Cameroon.
\end{enumerate}
\end{corrige}

\begin{corrige}[Exercise 4 — Direct calculation: contained metal]
\begin{enumerate}
\item Mass of copper = tonnage of ore $\times$ grade:
\[ m_{\mathrm{Cu}} = 5\times10^{6}\ \mathrm{t} \times \frac{0.8}{100} = 5\times10^{6}\times 0.008 = 40\,000\ \mathrm{t\ Cu}. \]
The deposit contains \textbf{40\,000 tonnes of copper metal}.
\item Life of the reserve:
\[ L = \frac{\text{contained metal}}{\text{annual production}} = \frac{40\,000}{12\,000} \approx 3.3\ \text{years}. \]
\item The real mine life may be shorter because recovery is never total: during mining and processing a fraction of the metal is lost (dilution by waste, metallurgical losses), so the recoverable metal is less than the contained metal. (Also acceptable: part of the ore may prove inaccessible or below cut-off grade.)
\end{enumerate}
\end{corrige}

\begin{corrige}[Exercise 5 — Resource or reserve?]
\begin{enumerate}
\item \textbf{Deposit A — inferred resource.} Geological knowledge is poor (few boreholes, continuity only assumed between widely spaced points) and no economic study exists; tonnage and grade can only be estimated with low confidence.
\item \textbf{Deposit B — proved reserve.} The dense drilling grid establishes continuity with high confidence (measured category) and a positive feasibility study demonstrates that extraction is technically feasible and economically viable at current prices.
\item \textbf{Deposit C — indicated resource.} Moderate drilling gives reasonable confidence in continuity (better than inferred, less than measured), but only a preliminary economic study exists, so reserve status cannot yet be claimed.
\item To upgrade C to a reserve one must: (i) carry out \textbf{additional drilling} (infill drilling) to raise geological confidence to the measured category; (ii) perform \textbf{metallurgical tests} to establish recovery; (iii) complete a full \textbf{feasibility study} covering mining method, costs, infrastructure, market and environmental impact, demonstrating profitability.
\end{enumerate}
\end{corrige}

\begin{corrige}[Exercise 6 — Interpreting a weathering profile]
\begin{enumerate}
\item From the surface downwards: (a) \textbf{ferruginous crust / iron cap (cuirasse)}, rich in iron oxides; (b) \textbf{bauxitic horizon}, rich in aluminium hydroxides; (c) \textbf{saprolite} (altered, clay-rich, friable rock retaining the parent structure); (d) \textbf{fresh parent rock} (syenite).
\item Bauxite is concentrated in horizon \textbf{(b)}, the bauxitic horizon. Its main aluminium minerals are the hydroxides \textbf{gibbsite} $\ce{Al(OH)3}$ (dominant), with subordinate \textbf{boehmite} and \textbf{diaspore} ($\ce{AlO(OH)}$).
\item Under a hot, humid climate, rainwater charged with dissolved $\ce{CO2}$ percolates through the rock and hydrolyses the silicate minerals (feldspars) of the syenite. Alkalis and alkaline earths ($\ce{Na+}$, $\ce{K+}$, $\ce{Ca2+}$) are removed in solution, and silica is leached away as soluble species. Aluminium and iron, being almost insoluble at near-neutral pH, remain behind and accumulate as residual hydroxides, so the residue is progressively enriched in Al — this is bauxitisation.
\item Favourable features of Cameroon's climate: (i) \textbf{high temperatures} all year round, which accelerate chemical weathering reactions; (ii) \textbf{abundant, well-distributed rainfall} (humid tropical climate of the Adamawa Plateau), which ensures continuous leaching and removal of soluble products.
\end{enumerate}
\end{corrige}

\begin{corrige}[Exercise 7 — Cross-cutting relations and genetic classification]
\begin{enumerate}
\item Sketch cross-section:
\begin{popfigure}
\begin{center}
\begin{tikzpicture}[scale=0.9]
\fill[gray!25] (0,0) rectangle (8,4);
\fill[poppink!40] (0,0) -- (0,2.2) -- (2.6,4) -- (5.4,4) -- (8,2.4) -- (8,0) -- cycle;
\draw[thick] (0,2.2) -- (2.6,4);
\draw[thick] (5.4,4) -- (8,2.4);
\node at (6.6,1.2) {schists};
\node at (2.2,1.0) {granite};
\draw[line width=2pt,poporange] (1.2,0) -- (3.4,4);
\node[poporange] at (1.6,3.6) {V1 (Cu)};
\draw[line width=2pt,popgold] (4.6,0) -- (2.6,4);
\node[popgold] at (4.9,3.5) {V2 (Au)};
\fill[popdark!70] (6.2,0) -- (6.9,0) -- (7.5,4) -- (6.8,4) -- cycle;
\node[white] at (7.15,2.0) {\rotatebox{76}{dyke}};
\end{tikzpicture}
\end{center}
\legende{Cross-section: granite in schists, veins V1 and V2, dolerite dyke.}
\end{popfigure}
\item Order from oldest to youngest, using the principle of cross-cutting relationships (a geological feature is younger than any feature it cuts):
\begin{enumerate}
\item \textbf{Schists} — oldest, they host everything else;
\item \textbf{Granite} — intrudes the schists, so it is younger;
\item \textbf{Vein V1 (chalcopyrite)} — cuts the schists, so post-dates them;
\item \textbf{Vein V2 (gold)} — cuts and offsets V1, so it is younger than V1;
\item \textbf{Dolerite dyke} — cuts and terminates both veins, so it is the youngest.
\end{enumerate}
\item Both mineralisations are \textbf{epigenetic}: they were introduced into pre-existing rocks after those rocks formed (the veins cut the schists). Genetically they are \textbf{hydrothermal} deposits: hot aqueous fluids circulated through fractures and precipitated quartz, chalcopyrite (Cu) and gold as the fluids cooled. They are neither syngenetic (not contemporaneous with the host), nor sedimentary/residual (no surface concentration process involved).
\end{enumerate}
\end{corrige}

\begin{corrige}[Exercise 8 — From resource to reserve]
\begin{enumerate}
\item A \textbf{resource} is a concentration of a useful substance in the crust whose location, approximate tonnage and grade are known or reasonably inferred, but whose profitable extraction has not been demonstrated. A \textbf{reserve} is the portion of a resource that has been proved by detailed exploration \emph{and} shown by a feasibility study to be extractable \textbf{economically and technically} under present conditions. Every reserve is a resource, but not every resource is a reserve.
\item Two ways the Mbalam resource could become a reserve without new discovery:
\begin{enumerate}
\item \textbf{Economic change:} a sustained rise in the world price of iron ore lowers the effective cut-off grade and makes the deposit profitable; equally, a fall in costs (energy, reagents, labour) has the same effect.
\item \textbf{Technical/infrastructural change:} construction of a railway and a deep-water port (e.g.\ towards Kribi) drastically reduces transport costs, and improved mining/processing technology raises recovery — both can turn a sub-economic resource into a viable reserve.
\end{enumerate}
\item A deposit may remain unexploited even as a reserve because of \textbf{non-geological obstacles}: lack of investment capital, political or administrative instability, land-use or environmental conflicts, or absence of the infrastructure needed to bring the product to market.
\end{enumerate}
\end{corrige}

\begin{corrige}[Exercise 9 — Structured question: ore, grade and reserves (20 marks)]
\textbf{(a) Definitions (8 marks — 2 per term: 1 for the definition, 1 for a valid example).}
\begin{itemize}
\item \textbf{Ore:} a rock or aggregate of minerals from which a useful substance can be extracted at a profit; e.g.\ bauxite as the ore of aluminium.
\item \textbf{Gangue:} the valueless minerals associated with the ore minerals, separated and discarded during processing; e.g.\ quartz in a gold vein.
\item \textbf{Grade:} the concentration of the valuable substance in the ore, in $\%$ or $\mathrm{g/t}$; e.g.\ a copper ore grading $1.2\ \%$ Cu.
\item \textbf{Cut-off grade:} the minimum grade at which material can be mined and processed profitably; e.g.\ material below $0.4\ \%$ Cu sent to waste in a given mine.
\end{itemize}
\textbf{(b) Contained metal (3 marks).}
\[ m_{\mathrm{Cu}} = 2\times10^{6}\ \mathrm{t}\times\frac{1.2}{100} = 24\,000\ \mathrm{t\ Cu}. \]
\emph{Marking: 1 mark formula, 1 mark correct substitution, 1 mark answer with unit.}
\textbf{(c) Life of the reserve (3 marks).}
\[ L = \frac{24\,000}{20\,000} = 1.2\ \text{years}. \]
Assumption stated: \textbf{total recovery} (all contained metal is mined and processed without loss). \emph{Marking: 1 mark method, 1 mark result, 1 mark the assumption.}
\textbf{(d) Three factors (6 marks — 2 per factor, any three).}
\begin{itemize}
\item \textbf{Grade and tonnage:} the concentration must be high enough and the volume large enough to repay investment.
\item \textbf{Market price and demand:} the value of the metal must exceed the total cost of extraction and processing.
\item \textbf{Technology and cost of extraction:} mining method (open pit vs underground), depth, metallurgical recoverability of the mineral.
\item \textbf{Infrastructure and location:} roads, railway, port, energy and water supply (e.g.\ Mbalam iron awaiting a railway).
\item \textbf{Environmental, legal and political conditions:} permits, stability, environmental costs.
\end{itemize}
\emph{Examiner expectation: precise vocabulary (ore, grade, cut-off grade), correct units in calculations, and factors clearly explained rather than listed.}
\end{corrige}

\begin{corrige}[Exercise 10 — Map exercise: mineral occurrences of Cameroon (20 marks)]
\textbf{(a) Placement (5 marks — 1 per correctly located and labelled occurrence).}
\begin{itemize}
\item \textbf{Gold — East Region:} in the Batouri--B\'etar\'e Oya--Kette area, eastern Cameroon, on the Archean--Proterozoic basement.
\item \textbf{Bauxite — Minim--Martap:} on the Adamawa Plateau, north of Ngaound\'er\'e.
\item \textbf{Iron — Mbalam:} East Region, near the Congo/Gabon border, south-east of Batouri.
\item \textbf{Nickel--cobalt — Lomie:} East Region, east of Yokadouma area, in the forest zone.
\item \textbf{Limestone — Figuil:} North Region, in the B\'enou\'e trough, south of Garoua.
\end{itemize}
\textbf{(b) Genetic types (10 marks — 2 per occurrence: 1 for the type, 1 for the justification).}
\begin{center}
\begin{tabularx}{\linewidth}{|l|l|X|}
\hline
\rowcolor{popblueL} \textbf{Occurrence} & \textbf{Genetic type} & \textbf{Justification} \\
\hline
Gold (East Region) & Hydrothermal (or placer) & Gold occurs in quartz veins cutting basement rocks (epigenetic hydrothermal), reworked into alluvial placers in rivers. \\
\hline
Bauxite (Minim--Martap) & Residual / weathering & Intense lateritic weathering of syenite leached silica and alkalis, concentrating residual Al hydroxides. \\
\hline
Iron (Mbalam) & Sedimentary (BIF), enriched by weathering & Banded iron formations deposited with ancient sediments, then enriched by leaching of silica. \\
\hline
Ni--Co (Lomie) & Residual / weathering (lateritic) & Deep weathering of ultrabasic rocks concentrates Ni and Co in the lateritic profile. \\
\hline
Limestone (Figuil) & Sedimentary & Marine chemical/biochemical precipitation of $\ce{CaCO3}$ in the Cretaceous B\'enou\'e trough. \\
\hline
\end{tabularx}
\end{center}
\textbf{(c) Prospective region (5 marks).} The \textbf{East Region} is the most prospective for gold: (i) it is underlain by \textbf{Archean to Proterozoic greenstone and schist belts} of the Congo craton margin, a rock association known worldwide to host gold; (ii) \textbf{known occurrences} (Batouri, B\'etar\'e Oya) and active artisanal workings prove the presence of a mineralised system; (iii) regional \textbf{faults and shear zones} provided the channels for gold-bearing hydrothermal fluids. \emph{Marking: 1 mark a coherent choice, 2 marks per justified geological criterion.}
\end{corrige}

\begin{corrige}[Exercise 11 — Essay: geology controls mineral distribution (20 marks)]
\textbf{Indicative mark scheme:} introduction and definition of terms (3); three Cameroonian deposits with correct genetic process (3 $\times$ 3 = 9); one world metallogenic province correctly explained (3); balanced discussion of political/economic factors (3); organisation, clarity, labelled diagram (2).

\textbf{Model answer (outline).}

\emph{Introduction.} Mineral deposits are not evenly distributed: they occur where specific geological processes have concentrated useful elements. The statement is therefore largely true at the scale of \emph{where} deposits exist, but must be qualified at the scale of \emph{whether} they are mined.

\emph{Development — Cameroonian examples.}
\begin{itemize}
\item \textbf{Bauxite of Minim--Martap (Adamawa):} residual deposit; intense lateritic weathering of syenite under a hot humid climate leached silica and alkalis and concentrated Al hydroxides on a stable, well-drained plateau surface.
\item \textbf{Iron of Mbalam (East Region):} sedimentary banded iron formations of the Precambrian craton, enriched by tropical weathering which removed silica and raised the iron grade.
\item \textbf{Gold of the East Region:} hydrothermal quartz veins localised by shear zones in Archean--Proterozoic greenstone/schist belts, reworked into placers.
\item \textbf{Nickel--cobalt of Lomie:} lateritic concentration over ultrabasic bedrock.
\item \textbf{Limestone of Figuil:} marine sedimentation in the Cretaceous B\'enou\'e trough.
\end{itemize}
In every case the \textbf{localising factor is geological}: the parent rock, the climate-driven weathering, the sedimentary basin or the structural channelway — none of which respects administrative boundaries.

\emph{World province.} The \textbf{Central African Copper Belt} (DRC--Zambia) illustrates this: stratiform copper mineralisation is confined to specific sedimentary horizons of the Katanga Supergroup, so the ore runs straight across the political border, following the geology. Similarly, gold provinces of the Kaapvaal, Canadian and West African cratons reflect the distribution of ancient greenstone belts.

\emph{Nuance.} Geology fixes the \emph{location} of deposits, but politics and economics decide their \emph{exploitation}: a reserve requires investment, stable legislation, and infrastructure — the Mbalam iron, though geologically proved, long awaited a railway and port; conversely the Figuil limestone is worked because a cement plant and market exist. Prices, technology and political decisions convert resources into reserves and mines.

\emph{Conclusion.} The distribution of mineral deposits is indeed written in the rocks and controlled by geology; but the map of \emph{producing mines} is drawn jointly by geology, economics and political choice.

\emph{Examiner expectations:} precise genetic vocabulary (syngenetic/epigenetic, residual, hydrothermal, BIF), correctly located examples, a clear two-sided argument, and a structured essay with introduction and conclusion.
\end{corrige}

\newpage

\coursec{Summary sheet}

\begin{popfigure}
\begin{center}\tcbox[colback=popink,colframe=popink,coltext=white,arc=9pt,boxsep=4pt,left=6pt,right=6pt]{\bfseries\small Economic Minerals (GEO05)}\end{center}
\vspace{-2pt}
\begin{tcbraster}[raster columns=3,raster equal height=rows,raster column skip=6pt,raster row skip=6pt,enhanced,blankest]
\begin{tcolorbox}[enhanced,colback=white,colframe=popblue,boxrule=0.9pt,arc=3pt,title={Economic mineral: notion},fonttitle=\bfseries\scriptsize,coltitle=white,colbacktitle=popblue,left=4pt,right=4pt,top=2pt,bottom=2pt,boxsep=1pt,toptitle=1.5pt,bottomtitle=1.5pt,halign=left]
\begin{itemize}[nosep,leftmargin=*,itemsep=1pt,font=\scriptsize]
\item Ore mineral, gangue, grade, tenor
\item Cut-off grade; value vs.\ cost
\end{itemize}
\end{tcolorbox}
\begin{tcolorbox}[enhanced,colback=white,colframe=popgreen,boxrule=0.9pt,arc=3pt,title={Resources \& Reserves},fonttitle=\bfseries\scriptsize,coltitle=white,colbacktitle=popgreen,left=4pt,right=4pt,top=2pt,bottom=2pt,boxsep=1pt,toptitle=1.5pt,bottomtitle=1.5pt,halign=left]
\begin{itemize}[nosep,leftmargin=*,itemsep=1pt,font=\scriptsize]
\item Inferred / indicated / measured
\item Probable / proven reserves
\end{itemize}
\end{tcolorbox}
\begin{tcolorbox}[enhanced,colback=white,colframe=poporange,boxrule=0.9pt,arc=3pt,title={Classification of deposits},fonttitle=\bfseries\scriptsize,coltitle=white,colbacktitle=poporange,left=4pt,right=4pt,top=2pt,bottom=2pt,boxsep=1pt,toptitle=1.5pt,bottomtitle=1.5pt,halign=left]
\begin{itemize}[nosep,leftmargin=*,itemsep=1pt,font=\scriptsize]
\item Syngenetic vs.\ epigenetic
\item Magmatic, hydrothermal, sedimentary, residual, metamorphic
\end{itemize}
\end{tcolorbox}
\begin{tcolorbox}[enhanced,colback=white,colframe=popgold,boxrule=0.9pt,arc=3pt,title={Formation I: magmatic \& hydrothermal},fonttitle=\bfseries\scriptsize,coltitle=white,colbacktitle=popgold,left=4pt,right=4pt,top=2pt,bottom=2pt,boxsep=1pt,toptitle=1.5pt,bottomtitle=1.5pt,halign=left]
\begin{itemize}[nosep,leftmargin=*,itemsep=1pt,font=\scriptsize]
\item Crystal settling, immiscibility, pegmatites
\item Veins, skarns, porphyry Cu
\end{itemize}
\end{tcolorbox}
\begin{tcolorbox}[enhanced,colback=white,colframe=poppurple,boxrule=0.9pt,arc=3pt,title={Formation II: sedimentary, weathering, metamorphic},fonttitle=\bfseries\scriptsize,coltitle=white,colbacktitle=poppurple,left=4pt,right=4pt,top=2pt,bottom=2pt,boxsep=1pt,toptitle=1.5pt,bottomtitle=1.5pt,halign=left]
\begin{itemize}[nosep,leftmargin=*,itemsep=1pt,font=\scriptsize]
\item Placers, BIF, evaporites
\item Bauxite, lateritic Ni; marble, graphite
\end{itemize}
\end{tcolorbox}
\begin{tcolorbox}[enhanced,colback=white,colframe=poppink,boxrule=0.9pt,arc=3pt,title={Occurrence: world to Cameroon},fonttitle=\bfseries\scriptsize,coltitle=white,colbacktitle=poppink,left=4pt,right=4pt,top=2pt,bottom=2pt,boxsep=1pt,toptitle=1.5pt,bottomtitle=1.5pt,halign=left]
\begin{itemize}[nosep,leftmargin=*,itemsep=1pt,font=\scriptsize]
\item Bauxite (Minim-Martap), Fe (Mbalam), Ni-Co (Lomi\'e)
\item Gold (East), rutile, limestone, oil
\end{itemize}
\end{tcolorbox}
\end{tcbraster}

\legende{Mind map of the chapter ``Economic Minerals'': six branches linking definitions, resource classification, deposit types, genetic processes and Cameroonian examples.}
\end{popfigure}

\begin{aretenir}[Summary Sheet 1 --- Key Definitions]
\begin{itemize}
\item \cle{Mineral}: a naturally occurring, inorganic solid with a definite chemical composition and an ordered internal (crystalline) structure.
\item \cle{Economic mineral}: any mineral (or rock) that can be extracted and sold at a \textbf{profit} under current technical and economic conditions. Profitability, not mere presence, is the deciding criterion.
\item \cle{Ore}: a rock containing one or more valuable minerals in sufficient concentration to be mined profitably. The \cle{ore mineral} is the valuable mineral; the \cle{gangue} is the worthless associated material (e.g.\ quartz, calcite) that must be separated during processing.
\item \cle{Grade (tenor)}: the concentration of the useful substance in the ore, expressed in \% for common metals (e.g.\ $55\ \%$ Fe) or in $\mathrm{g/t}$ (equivalent to ppm) for precious metals (e.g.\ $2\ \mathrm{g/t}$ Au).
\item \cle{Cut-off grade}: the minimum grade below which exploitation is not profitable. It is \emph{not fixed}: it varies with the world price of the metal, the available technology, energy and transport costs, and the location of the deposit.
\item Scale of knowledge: \cle{mineral occurrence} (anomalous concentration noted) $\rightarrow$ \cle{prospect} (showing investigated, encouraging results) $\rightarrow$ \cle{deposit} (body of ore of demonstrated economic interest).
\end{itemize}
\end{aretenir}

\begin{aretenir}[Summary Sheet 2 --- Resources and Reserves]
\begin{itemize}
\item \cle{Resource}: the total quantity of a commodity known or suspected to exist in the ground, whether or not it is currently economic to extract.
\item \cle{Reserve}: the part of the resource that is \textbf{economically extractable now} with existing technology and at current prices.
\item Increasing \emph{geological confidence} (how well the body is known): \textbf{inferred} $\rightarrow$ \textbf{indicated} $\rightarrow$ \textbf{measured} resource.
\item Increasing \emph{economic confidence} (how sure the extraction is): \textbf{probable} $\rightarrow$ \textbf{proven} reserve.
\item Reserves are \emph{dynamic}: they grow when metal prices rise or technology improves (former resources become reserves), and shrink when prices fall or costs increase.
\item \cle{Reserve life} $=$ reserves $\div$ annual production; it estimates how many years a mine (or a country) can continue producing at the present rate.
\end{itemize}
\end{aretenir}

\begin{aretenir}[Summary Sheet 3 --- Classification of Mineral Deposits]
\begin{itemize}
\item \cle{Syngenetic} deposits: formed \emph{at the same time} as the host rock (chromite layers in basic intrusions, placers, banded iron formations, evaporites). They are typically concordant (parallel to bedding or layering).
\item \cle{Epigenetic} deposits: formed \emph{after} the host rock, by introduction of material into it (hydrothermal veins, skarns, replacement bodies). They are typically discordant (cross-cutting).
\item Genetic classes to master: \textbf{magmatic}, \textbf{hydrothermal}, \textbf{sedimentary}, \textbf{residual/weathering}, \textbf{metamorphic}.
\end{itemize}
\end{aretenir}

\begin{aretenir}[Summary Sheet 4 --- Formation Processes to Know]
\begin{itemize}
\item \textbf{Magmatic}: early crystallisation and settling of dense crystals $\rightarrow$ layered chromite and magnetite; liquid immiscibility $\rightarrow$ Ni--Cu sulphide ores; concentration of rare elements in late residual fluids $\rightarrow$ \cle{pegmatites} (Sn, W, Li, gemstones).
\item \textbf{Hydrothermal}: hot aqueous fluids circulating through rocks deposit dissolved minerals on cooling, especially in fractures $\rightarrow$ \cle{veins} (Au, Ag, Pb, Zn); reaction of fluids with limestone $\rightarrow$ \cle{skarn} (W, Fe, Cu); large low-grade disseminated bodies $\rightarrow$ \cle{porphyry copper}.
\item \textbf{Sedimentary}: mechanical sorting by water or wind $\rightarrow$ \cle{placers} (gold, diamond, rutile, cassiterite); chemical precipitation $\rightarrow$ \cle{banded iron formations} (BIF) and \cle{evaporites} (halite, gypsum, potash); biochemical accumulation $\rightarrow$ phosphates, coal, petroleum.
\item \textbf{Weathering/residual}: intense tropical leaching removes silica and bases and concentrates aluminium $\rightarrow$ \cle{bauxite} (lateritisation); downward-percolating waters can upgrade sulphide ores $\rightarrow$ \cle{supergene enrichment} of copper.
\item \textbf{Metamorphic}: recrystallisation under heat and pressure $\rightarrow$ marble, graphite, garnet, kyanite; contact metasomatism around intrusions $\rightarrow$ skarn ores.
\end{itemize}
\end{aretenir}

\begin{aretenir}[Summary Sheet 5 --- Occurrence: Cameroon to Remember]
\begin{itemize}
\item \cle{Bauxite}: Minim--Martap and Ngaoundal (Adamawa Plateau; residual, lateritic deposits on the plateau surface).
\item \cle{Iron}: Mbalam--Nabeba (South-East; itabirite/BIF-type ore at the margin of the Congo craton).
\item \cle{Nickel--cobalt}: Lomi\'e (East Region; lateritic deposit developed by deep tropical weathering of ultrabasic rocks).
\item \cle{Gold}: Batouri--B\'etar\'e-Oya (East Region; alluvial placers and primary lode occurrences).
\item \cle{Rutile}: Akonolinga (Centre); \cle{limestone}: Figuil (North; cement manufacture); \cle{oil and gas}: Rio del Rey and Douala sedimentary basins; \cle{cassiterite}: Mayo-Darl\'e (North-West).
\end{itemize}
\end{aretenir}

\begin{methode}[Summary Sheet 6 --- Working Methods]
\textbf{To judge whether a body is an economic deposit:}
\begin{enumerate}
\item Identify the \emph{commodity} and its use.
\item Estimate the \emph{tonnage} and multiply by the \emph{grade} to obtain the contained metal.
\item Compare the grade with the \emph{cut-off grade} under current prices and technology.
\item Conclude: uneconomic concentration (occurrence/resource) or exploitable ore (reserve).
\end{enumerate}
\textbf{To classify a deposit:}
\begin{enumerate}
\item Ask: ``Is the ore the same age as the host rock?'' Yes $\rightarrow$ syngenetic; no $\rightarrow$ epigenetic.
\item Read the textures and geological setting (layering, veins, alteration, sedimentary structures).
\item Assign the genetic process (magmatic, hydrothermal, sedimentary, residual, metamorphic) and justify it with one observation.
\end{enumerate}
\end{methode}

\begin{attention}[Summary Sheet 7 --- Common Mistakes]
\begin{itemize}
\item Confusing \emph{mineral}, \emph{ore mineral} and \emph{ore}: an ore is a \textbf{rock} (a mixture of ore minerals and gangue), not a single mineral.
\item Stating that reserves are fixed quantities: they depend on price, technology and costs, and therefore change with time.
\item Calling every deposit syngenetic: hydrothermal veins and skarns are \textbf{epigenetic} because they post-date their host rocks.
\item Expressing the grade of gold in \%: precious metals are quoted in $\mathrm{g/t}$ (ppm).
\item Writing ``bauxite mineral'': bauxite is a \textbf{rock} (a mixture of aluminium hydroxides such as gibbsite), used as the ore of aluminium.
\end{itemize}
\end{attention}

\begin{aretenir}[Summary Sheet 8 --- GCE Examination Tips]
\begin{itemize}
\item Define each term precisely, then give \textbf{one Cameroonian example} per process --- examiners reward local illustration (e.g.\ lateritisation $\rightarrow$ Minim--Martap bauxite).
\item Be ready to draw quick labelled sketches: a hydrothermal vein cross-cutting sedimentary strata; a laterite profile (topsoil -- bauxite/ferricrete -- saprolite -- fresh bedrock); a layered intrusion with a chromite seam.
\item In resource/reserve questions, always structure the answer along the two axes: \emph{geological confidence} (inferred--indicated--measured) and \emph{economic viability} (resource--probable--proven reserve).
\item Show the reasoning in calculations: contained metal $=$ tonnage $\times$ grade; reserve life $=$ reserves $\div$ annual production; always carry the units.
\end{itemize}
\end{aretenir}

\findecours

\end{document}
